% !TEX program = xelatex
% !TEX encoding = UTF-8 Unicode
%
% 제어공학개론 -- 강의노트 (한글판), 2026년 2학기
% 회차를 따로 나누지 않고 하나의 문서로 이어 붙인다. 새 강의는 해당 절에 녹인다.
% 영문판과 공유하는 설정은 전부 controlnotes.sty 에 있다.
% 컴파일:  xelatex control_notes_KR.tex   (목차 때문에 세 번 — publish.sh 가 해 준다)
%
% 앞으로의 규칙:
%   1. 새 강의 내용은 주제가 맞는 대단원의 절에 녹인다 — 새 장으로 덧붙이지
%      않는다.
%   2. 정말로 새 주제가 시작되면 \partdivider 로 새 대단원을 연다. 절 번호는
%      대단원과 무관하게 이어서 붙인다 — 다시 1부터 세지 않는다.
%   3. 그 강의의 새 내용이 시작되는 자리에 \lecturebreak 를 넣고, 앞부분 회차
%      표에 같은 자리를 가리키는 행을 한 줄 추가한다. (lectures.py 가 국문판
%      표의 마지막 행을 읽어 manifest 의 lectures 필드를 만든다.)
%   4. "강의의 흐름" 목록에서 해당 대단원을 진행 중/완료로 고친다. 예정 항목은
%      근거가 있는 것 — 강의에서 실제로 밝힌 계획 — 만 적는다.
%   5. 두 판의 대단원·절·label 구조는 1:1 로 유지한다.
%
\documentclass[11pt]{article}
\usepackage[korean]{controlnotes}

%% ---- 머리표에 들어가는 항목 ------------------------------------------------
%% \lectitle 은 기존 제목 블록의 "강의노트" 와 아래 회차↔절 대응표의 범위를
%% 합친 것이고, 학기와 날짜는 저자 밑에 있던 그 줄의 나머지 두 조각이다.
%% \lecnum 과 \lecturer 는 비워 둔다 — 회차별이 아닌 단일 문서이고 교수님
%% 성함은 노트 어디에도 적혀 있지 않다.
\coursename{제어공학개론}
\term{2026학년도 2학기}
\lectitle{강의노트: 제1강$\sim$제10강}
\lecdate{최종 수정 2026년 10월 10일}
\scribe{우수한}

\begin{document}
\makelecheader

\section*{강의의 흐름}

제1강에서 밝힌 수업의 흐름은 세 단계였다. 먼저 동적 시스템을 전달함수와 상태변수
표현법으로 \emph{모델링}하고 그 모델을 \emph{해석}하는 법을 배우고, 중반부에는
선형대수와 MATLAB에 집중하며, 후반부에는 안정도, 제어기 설계, 고전 제어를 다룬다.
구체적으로 배울 것으로는 PID 제어, 상태궤환 제어기, 상태변수 관측기, 추종 제어기,
Bode 선도와 Nyquist 선도가 꼽혔다. 이 노트의 대단원도 그 계획을 따른다.

\begin{itemize}[leftmargin=1.2em,itemsep=1pt,topsep=3pt]
  \item 제1부 --- 되먹임과 모델링 \hfill {\small\S\S1--2}
  \item 제2부 --- 상태공간 모델과 선형화 \hfill {\small\S\S3--4}
  \item 제3부 --- 전달함수 \hfill {\small\S5}
  \item 제4부 --- 주파수 응답과 안정도 \hfill {\small\S\S6--7}
  \item 제5부 --- 상태공간을 위한 선형대수 \hfill {\small\S\S8--11, 진행 중}
  \item \emph{안정도, 제어기 설계, 고전 제어} --- 상태궤환 제어기, 상태변수 관측기,
        추종 제어기, PID 제어 \hfill {\small 예정}
\end{itemize}

예정 항목은 제1강에서 설명한 그대로 적은 것이며, 이것이 대단원으로 어떻게 나뉠지는
강의가 진행되는 대로 정한다. 질문~\ref{q:goal}의 목표 --- 두 수레를 인자 값을 모르는
채로 5초 뒤 $1$\,m 앞에 세우는 것 --- 는 제어기 설계 문제이고, 지금까지의 대단원은
그 설계에 필요한 모델과 해석을 쌓는 과정이다.

\section*{이 노트의 구성}

이 노트는 날짜순이 아니라 주제순으로 엮었다. 수업의 단계마다 대단원(부)을 하나씩
두고, 새 강의는 새 장으로 덧붙이지 않고 해당하는 대단원과 절 안에 녹여 넣으므로 같은
말이 두 번 나오지 않는다. 절 번호는 대단원과 상관없이 문서 처음부터 끝까지 이어진다.
각 강의의 새 내용이 시작되는 자리에는 강의 번호와 날짜를 단 가는 가로줄을 그어 두었고,
아래 표도 같은 자리를 가리킨다. 여러 강의를 한 단위로 묶어 정리한 곳에는 가로줄
하나가 그 범위 전체를 맡는다. \textbf{보충}과 \textbf{예제}는 수업 시간에 칠판에 다
적지 못한 유도·예제·그림을 채운 것으로, 주제에 맞는 자리에 놓았다. 그러니 가로줄은
그 강의가 시작된 자리를 표시할 뿐, 그 뒤의 내용이 모두 그 강의에서 나왔다는 뜻은
아니다.

\smallskip
\begin{center}\small
\begin{tabular}{@{}lll@{}}
\toprule
회차 & 날짜 & 시작하는 곳\\
\midrule
제1강 & 2026.\,09.\,02 (수) & \S1\\
제2강 & 2026.\,09.\,07 (월) & \S2\\
제3강 & 2026.\,09.\,09 (수) & \S4\\
제4강 & 2026.\,09.\,14 (월) & \S5\\
제5강 & 2026.\,09.\,16 (수) & \S\ref{sec:polezero}\\
제6강 & 2026.\,09.\,21 (월) & \S\ref{sec:second}\\
제7$\sim$10강 & 2026.\,09.\,23 (수) $\sim$ 10.\,07 (수) & \S\ref{sec:vspace}\\
\bottomrule
\end{tabular}
\end{center}

{\small\tableofcontents}

%% ================================================================ 제1부
\partdivider{제1부 --- 되먹임과 모델링}

\lecturebreak{제1강}{2026년 9월 2일}

%%%%%%%%%%%%%%%%%%%%%%%%%%%%%%%%%%%%%%%%%%%%%%%%%%%%%%%%%%%%%%%%%%%%%%%%%%%%%%
\section{제어공학이란 무엇을 하는 학문인가}

\subsection{질문 하나} 제어공학은 질문 하나에서 출발한다.
\begin{center}
\emph{원하는 출력을 얻으려면 입력을 어떻게 주어야 하는가?}
\end{center}
이 질문에 답하려면 먼저 입력과 출력이 어떤 관계인지를 알아야 하고, 물리적인
시스템에서 그 관계는 거의 언제나 \emph{미분방정식}으로 주어진다. 그래서 이
과목은 제어기가 아니라 미분방정식에서 시작한다.

한 학기 내내 따라다닐 단어가 둘 있다. 첫째는 \textbf{불확실성}(uncertainty)이다.
우리가 적어 놓은 미분방정식은 실제 세계가 따르는 그 방정식이 결코 아니다. 둘째는
\textbf{되먹임}(feedback)이다. 출력을 재서 입력으로 되돌리는 이 구조가 설계를
불확실성 속에서 살아남게 한다.

실제 제어 시스템은 센서, 액추에이터, 컴퓨터, 제어기, 그리고 이들 사이의 통신으로
이루어진다. 이 노트에서는 그 전부를 들어가는 신호 $u(t)$와 나오는 신호 $y(t)$로
추상화한다. 출력 $y$가 반드시 연속적으로 변하며 따라가야 할 양일 필요는 없다.
``저 지점에 가서 서 있어라'' 같은 것도 똑같이 출력으로 잡을 수 있고, 같은
도구로 다룬다.

\subsection{동적 대 정적} \emph{정적}(static) 관계는 시각 $t$의 $y$가 같은 시각의
$u$만으로 정해지는 관계다. \emph{동적}(dynamic) 관계에는 기억이 있다. 지금의
출력이 과거의 입력에 의존한다. 가장 단순한 예는 1차 선형 미분방정식
\begin{equation}\label{eq:rl}
   L\dot y + Ry = u(t)
\end{equation}
로, 예컨대 전압원 $u$로 구동되는 $RL$ 회로에서 $y$를 전류로 본 것이다. 앞으로
$y^{*}(t)$는 \emph{원하는} 출력을, $u(t)$는 제어 입력을 뜻한다.

\begin{note}
$y$는 그때그때 출력인 양을 가리키는 기호일 뿐이다. \eqref{eq:rl}에서는 전류였고
바로 다음 예제에서는 위치다. 수식은 신경 쓰지 않지만 단위는 신경 쓰므로, 무엇의
$y$인지는 늘 확인해야 한다.
\end{note}

\subsection{첫 예제: 중력을 이기고 물체를 붙들어 두기}
연직 방향으로만 움직이는 질량 $m$의 물체를 생각하자(그림~\ref{fig:mass}).
위쪽을 양으로 하여 위치를 $y(t)$라 하고, 제어력 $u(t)$와 중력이 작용한다.
\begin{equation}\label{eq:mass}
   m\ddot y = u - mg .
\end{equation}
목표는 높이 $h$에 붙들어 두는 것이다. 즉 원하는 출력은 상수 $y^{*}(t)=h$이고,
물체는 거기서 정지한 채로 출발한다. 즉 $y(0)=h$, $\dot y(0)=0$.

\begin{figure}[htb]
\centering
\begin{tikzpicture}[>=Stealth,scale=1.0]
  \draw[thick] (-0.15,0) -- (5.0,0);
  \foreach \x in {0.05,0.35,...,4.95}{\draw[thin] (\x,0) -- ++(-0.22,-0.22);}
  \draw[->,thick] (0.25,0) -- (0.25,3.55) node[left] {$y$};
  \draw[dashed] (0.25,2.50) -- (4.95,2.50);
  \node[right] at (4.95,2.50) {$y^{*}=h$};
  \draw[thick,fill=black!5] (1.95,2.10) rectangle (3.15,2.90);
  \node at (2.55,2.50) {$m$};
  \draw[very thick,->] (2.20,1.15) -- (2.20,2.07);
  \node[left,sig] at (2.10,1.58) {$u(t)$};
  \draw[very thick,->] (2.90,2.07) -- (2.90,1.15);
  \node[right,sig] at (2.92,1.58) {$mg$};
  \draw[<->] (1.15,0) -- (1.15,2.50);
  \node[left,sig] at (1.13,1.25) {$y(t)$};
\end{tikzpicture}
\caption{중력을 이기고 제어력으로 붙들어 둔 물체. $y(t)$가 $h$에 머물도록
$u(t)$를 고르는 것이 제어 문제다.}
\label{fig:mass}
\end{figure}

\smallskip
\noindent\textbf{(a) 개루프.} \eqref{eq:mass}에 $\ddot y=0$을 넣으면 $u=mg$가
나온다. 이 상수 입력을 주면 $\ddot y \equiv 0$이고, 위 초기조건과 함께 보면
$y(t)\equiv h$가 영원히 유지된다. 종이 위에서는 되먹임 없이 문제가 풀렸다.

\smallskip
\noindent\textbf{(b) 문제: 불확실성.} 그런데 우리는 $m$을 모른다. 아는 것은
추정값 $\hat m$뿐이고, 그래서 실제로 주는 입력은 $u=\hat m g$다. 그러면
\eqref{eq:mass}는 $m\ddot y = \hat m g - mg = (\hat m - m)g$가 되어 \emph{0이
아닌 일정한 가속도}가 남는다. 두 번 적분하면
\[
   y(t) = h + \frac{(\hat m-m)g}{2m}\,t^{2}\;\longrightarrow\;\pm\infty .
\]
모델링 오차 $\hat m - m$이 아무리 작아도 물체는 결국 떠나 버린다. 개루프 설계는
단지 부정확한 것이 아니라 질적으로 실패한다.

\smallskip
\noindent\textbf{(c) 처방: 되먹임.} $y$와 $\dot y$를 재서
\begin{equation}\label{eq:pd}
   u(t) = -c\,\dot y - k\,(y-h) + \hat m g
\end{equation}
를 준다. 여기서 $c$와 $k$는 \emph{이득}(gain), 즉 우리가 정하는 숫자다. 우리가
직접 고르므로 여기에는 불확실성이 없다. \eqref{eq:pd}를 \eqref{eq:mass}에 넣으면
\begin{equation}\label{eq:massfb}
   m\ddot y + c\,\dot y + k\,(y-h) = (\hat m-m)g
\end{equation}
이고, $c,k>0$이면 상수로 구동되는 감쇠 2차 시스템이므로 $y$는 어딘가에 정착한다.
$\ddot y=\dot y=0$으로 두면 그 정착값은
\begin{equation}\label{eq:offset}
   y(\infty) = h + \frac{(\hat m - m)g}{k}
\end{equation}
이다. 오차가 이제 \emph{유계}이고, $k$를 키우면 줄어든다. 되먹임이 불확실성을
없앤 것이 아니라, 발산하던 오차를 작은 오차로 바꾼 것이다.

\begin{question}
왜 이게 잘 되는가? $c$와 $k$는 얼마로 잡아야 하는가? $k$를 얼마나 키워도 되는가
(어디서부터 다른 것이 망가지는가)? 이 질문들에 답하는 것이 남은 학기의 내용이다.
\end{question}

\subsection{불확실성, 강인성, 그리고 학습}
공칭 모델에 대해 제어기를 설계하는 것은 일의 절반이다. 나머지 절반은 실제
시스템이 될 수 있는 플랜트의 \emph{집합} 전체에 대해 여전히 잘 동작함을
보장하는 것이고, 이것이 \emph{강인}(robust) 제어공학이 다루는 문제다.

수업에서는 이것을 강화학습과 대비해서 설명했다. 강화학습은 모델 대신 보상
신호로 정책을 학습한다. 그 특유의 난점이 \emph{sim-to-real gap}이다. 시뮬레이터
안에서 잘 맞춘 정책이 실제 하드웨어에서는 그대로 동작하지 않는다. 표준적인
처방인 domain randomization 계열 --- 무작위화한 시뮬레이터 집합 위에서 학습시켜
어느 한 불일치도 치명적이지 않게 만드는 것 --- 은 사실 강인성 아이디어를 다른
이름으로 쓴 것이다. 휴리스틱하지만, 쓸 만한 물리 모델을 얻기 어려운 곳에서는
바로 그 점 때문에 강력하다. 반대로 모델 기반 설계는 수학적이어서, 원리적으로는
요구사항에서 거꾸로 제어기를 유도해 갈 수 있다. 물론 그것도 쉽지 않고, 이 과목의
상당 부분이 그 방법을 배우는 데 쓰인다.

%% ------------------------------------------------------------------ 제2강
\lecturebreak{제2강}{2026년 9월 7일}

%%%%%%%%%%%%%%%%%%%%%%%%%%%%%%%%%%%%%%%%%%%%%%%%%%%%%%%%%%%%%%%%%%%%%%%%%%%%%%
\section{모델링: Two-Cart 시스템}

\subsection{한 학기 동안 쓸 예제}
이 과목을 끝까지 따라다닐 예제는 그림~\ref{fig:twocart}의 \emph{two-cart
시스템}이다. 질량 $m_1$인 수레가 위치 $p_1(t)$에서 힘 $u(t)$로 밀리고, 강성
$k$인 스프링과 감쇠계수 $b$인 댐퍼를 통해 질량 $m_2$, 위치 $p_2(t)$인 두 번째
수레와 연결되어 있다. 두 수레 모두 계수 $c$의 지면 마찰을 받는다. 측정하는
출력은 두 번째 수레의 위치 $y(t) = p_2(t)$다. 그리고 $p_1(0)=p_2(0)=0$으로 둔다.
즉 위치는 출발점 기준으로 잰다.

\begin{figure}[htb]
\centering
\begin{tikzpicture}[>=Stealth,scale=1.05]
  \draw[thick] (-0.85,0) -- (7.75,0);
  \foreach \x in {-0.6,-0.25,...,7.7}{\draw[thin] (\x,0) -- ++(-0.25,-0.25);}
  \draw[thick,fill=black!5] (0.40,0.36) rectangle (2.00,1.60);
  \node at (1.20,0.98) {$m_1$};
  \draw[thick,fill=black!5] (5.00,0.36) rectangle (6.60,1.60);
  \node at (5.80,0.98) {$m_2$};
  \foreach \x in {0.80,1.60,5.40,6.20}{\draw[thick,fill=white] (\x,0.18) circle (0.18);}
  \draw[thick] (2.00,1.28) -- (2.35,1.28);
  \draw[thick,decorate,decoration={coil,aspect=0.62,segment length=3.8mm,amplitude=2.2mm}]
       (2.35,1.28) -- (4.65,1.28);
  \draw[thick] (4.65,1.28) -- (5.00,1.28);
  \node[above] at (3.50,1.44) {$k$};
  \draw[thick] (2.00,0.68) -- (3.20,0.68);
  \draw[thick] (3.20,0.46) rectangle (3.78,0.90);
  \draw[thick] (3.62,0.46) -- (3.62,0.90);
  \draw[thick] (3.62,0.68) -- (5.00,0.68);
  \node[below] at (2.58,0.63) {$b$};
  \draw[very thick,->] (-0.80,0.98) -- (0.38,0.98);
  \node[above] at (-0.22,1.02) {$u(t)$};
  \draw[dashed,gray] (1.20,0) -- (1.20,-0.62);
  \draw[->] (1.20,-0.62) -- (0.35,-0.62);
  \node[right,lbl] at (1.26,-0.62) {$c\,\dot p_1$};
  \draw[dashed,gray] (5.80,0) -- (5.80,-0.62);
  \draw[->] (5.80,-0.62) -- (4.95,-0.62);
  \node[right,lbl] at (5.86,-0.62) {$c\,\dot p_2$};
  \draw[dashed] (1.20,1.60) -- (1.20,2.42);
  \draw[dashed] (5.80,1.60) -- (5.80,2.42);
  \draw[->] (1.20,2.22) -- ++(0.90,0); \node[above] at (1.65,2.22) {$p_1(t)$};
  \draw[->] (5.80,2.22) -- ++(0.90,0); \node[above] at (6.25,2.22) {$p_2(t)=y(t)$};
\end{tikzpicture}
\caption{Two-cart 시스템. 스프링 $k$와 댐퍼 $b$는 두 수레 \emph{사이}에 걸려
있고, 지면 마찰 $c$는 각 수레에 따로 작용하여 그 속도를 거스른다.}
\label{fig:twocart}
\end{figure}

\begin{question}[학기 전체의 목표]\label{q:goal}
\emph{5초 뒤에 1\,m 앞에서 정지}하게 만들어라. 그것도 $m_1$, $k$, $m_2$의 값을
모르는 채로 할 수 있게.
\end{question}

질문~\ref{q:goal}의 뒷문장이 진짜 내용이다. 플랜트를 정확히 알고 있으면 목표를
맞히는 것은 누구나 할 수 있다. 그 값을 몰라도 동작하는 것이 제어 설계다.

\subsection{운동방정식}
각 수레에 Newton의 법칙을 쓴다. 수레 1은 $u$로 밀리고, 스프링에 의해 수레 2 쪽으로
당겨지고, 두 수레의 상대운동에 의해 감쇠되며, 지면 마찰로 끌린다. 수레 2도
입력만 빼면 같다.
\begin{equation}\label{eq:eom}
\left\{
\begin{aligned}
   m_1\ddot p_1 &= u(t) + k\,(p_2-p_1) + b\,(\dot p_2-\dot p_1) - c\,\dot p_1,\\
   m_2\ddot p_2 &= \phantom{u(t)\;+{}} k\,(p_1-p_2) + b\,(\dot p_1-\dot p_2) - c\,\dot p_2,\\
   y &= p_2 .
\end{aligned}\right.
\end{equation}

\begin{remark}[결합항의 부호]
스프링과 나란히 달린 댐퍼는 양 끝의 \emph{상대} 운동에 저항하므로, 수레 1에
작용하는 댐퍼 힘은 스프링 힘과 정확히 같은 부호 구조를 갖는다. 즉
$+\,k(p_2-p_1)$과 $+\,b(\dot p_2-\dot p_1)$이고, 수레 2에는 그 거울상이
작용한다. 이 부호가 물리적으로 맞는 부호임을 빠르게 확인하는 방법은 스프링의
늘어난 양 $q=p_1-p_2$를 보는 것이다. 두 식을 빼면 (간단히 $c=0$으로 두고)
\[
   \ddot q + b\Bigl(\tfrac1{m_1}+\tfrac1{m_2}\Bigr)\dot q
           + k\Bigl(\tfrac1{m_1}+\tfrac1{m_2}\Bigr)q = \frac{u}{m_1}
\]
로 감쇠 진동자가 된다. 두 수레가 서로 엇갈려 흔들리는 내부 모드가 댐퍼를 통해
에너지를 잃는 것이고, 당연히 그래야 한다. 두 결합항 중 어느 하나의 부호라도
뒤집으면 그 모드는 음의 강성이나 음의 감쇠를 갖게 된다.
\end{remark}

\begin{note}[기호]
이 과목 전체에서 $c$는 \emph{지면 마찰} 계수, $b$는 \emph{두 수레 사이의
댐퍼}다.
\end{note}

\subsection{같은 시스템을 적는 세 가지 방법}
\begin{figure}[htb]
\centering
\begin{tikzpicture}[>=Stealth]
  \node (ph) [blk] {물리\\ 시스템};
  \node (de) [blk,right=17mm of ph] {미분\\ 방정식};
  \node (tf) [blk,above right=6mm and 20mm of de] {전달함수 표현형\\ $P(s)=Y(s)/U(s)$};
  \node (ss) [blk,below right=6mm and 20mm of de] {상태공간 표현형\\ $\dot x=Ax+Bu$};
  \draw[->,thick] (ph) -- node[above,lbl] {모델링} (de);
  \draw[->,thick] (de.east) -- node[above,sloped,lbl,pos=0.60] {라플라스} (tf.west);
  \draw[->,thick] (de.east) -- node[below,sloped,lbl,pos=0.60] {상태변수} (ss.west);
  \draw[<->,thick] (tf) -- node[right,lbl,align=left]
        {\ref{sec:ss2tf}절\\ 실현} (ss);
  \draw[->,dashed] (ph.north) to[bend left=26]
        node[above,lbl] {data-driven} (tf.north west);
  \draw[->,dashed] (ph.south) to[bend right=26]
        node[below,lbl] {data-driven} (ss.south west);
\end{tikzpicture}
\caption{물리 시스템에서 수학적 모델로 가는 경로. 실선이 이 과목이 가는 길이다.
점선은 데이터로부터의 시스템 식별$\cdot$학습(신경망 등)으로, 미분방정식 단계를
아예 건너뛴다.}
\label{fig:routes}
\end{figure}

전달함수 표현형과 상태공간 표현형은 같은 미분방정식을 적는 두 가지 방법이고,
선형 시스템에서는 서로 오갈 수 있다. 그림~\ref{fig:routes}의 양방향 화살표가
그것이다. 두 방향 중 하나는 \ref{sec:ss2tf}절에서 실제로 계산한다.

%% ================================================================ 제2부
\partdivider{제2부 --- 상태공간 모델과 선형화}

%%%%%%%%%%%%%%%%%%%%%%%%%%%%%%%%%%%%%%%%%%%%%%%%%%%%%%%%%%%%%%%%%%%%%%%%%%%%%%
\section{상태공간 표현}

\subsection{상태변수}
\eqref{eq:eom}은 2계 방정식이다. 시스템이 가진 독립적인 ``기억'' 하나마다 변수를
하나씩 도입하자. 여기서는 수레마다 위치 하나와 속도 하나다.
\[
   x_1 = p_1,\qquad x_2=\dot p_1,\qquad x_3=p_2,\qquad x_4=\dot p_2 .
\]
그러면 \eqref{eq:eom}은 네 개의 \emph{1계} 방정식이 된다.
\begin{equation}\label{eq:statelines}
\left\{
\begin{aligned}
   \dot x_1 &= x_2, \\
   \dot x_2 &= \frac{k}{m_1}(x_3-x_1) - \frac{c}{m_1}x_2 + \frac{b}{m_1}(x_4-x_2) + \frac{1}{m_1}u, \\
   \dot x_3 &= x_4, \\
   \dot x_4 &= \frac{k}{m_2}(x_1-x_3) - \frac{c}{m_2}x_4 + \frac{b}{m_2}(x_2-x_4), \\
   y &= x_3 .
\end{aligned}\right.
\end{equation}
네 식은 서로 잔뜩 얽혀 있지만 하나하나는 모두 1계이고, 바로 그 점이 이 식들을
다루기 쉽게 만든다. 컴퓨터로 풀기에도, 선형대수로 다루기에도 그렇다.

\subsection{표준형}\label{sec:stdform}
$x=(x_1,x_2,x_3,x_4)\tp$로 모으면 \textbf{상태변수 표현형의 표준형}에 도달한다.
\begin{equation}\label{eq:std}
   \boxed{\;\dot x(t) = A\,x(t) + B\,u(t),\qquad y(t) = C\,x(t) + D\,u(t).\;}
\end{equation}
각 신호의 차원은
\begin{equation}\label{eq:dims}
   x(t)\in\R^{n},\qquad u(t)\in\R^{m},\qquad y(t)\in\R^{p}
\end{equation}
이고, 따라서 $A\in\R^{n\times n}$, $B\in\R^{n\times m}$, $C\in\R^{p\times n}$,
$D\in\R^{p\times m}$이다. 공간 $\R^{n}$을 \textbf{상태공간}(state space)이라
부르고, 상태변수의 개수 $n$이 시스템의 \emph{차수}(order)다. $m=p=1$인 시스템을
SISO(단일입력 단일출력)라 하고, 그렇지 않으면 MIMO(다중입력 다중출력)라 한다.
세 차원은 서로 독립이라는 점에 유의하자. 상태는 많은데 입력은 하나뿐인 플랜트도
있고, 입력은 여럿인데 출력은 하나인 플랜트도 있다.

Two-cart 시스템에서는 $n=4$, $m=p=1$이다. 성분을 짧게 적기 위해 $b=0$으로 두면
\begin{equation}\label{eq:ABCD}
 A=\begin{bmatrix}
     0 & 1 & 0 & 0\\[2pt]
     -\dfrac{k}{m_1} & -\dfrac{c}{m_1} & \dfrac{k}{m_1} & 0\\[6pt]
     0 & 0 & 0 & 1\\[2pt]
     \dfrac{k}{m_2} & 0 & -\dfrac{k}{m_2} & -\dfrac{c}{m_2}
   \end{bmatrix},\quad
 B=\begin{bmatrix}0\\[2pt] \dfrac{1}{m_1}\\[6pt] 0\\[2pt] 0\end{bmatrix},\quad
 C=\begin{bmatrix}0&0&1&0\end{bmatrix},\quad D = 0 .
\end{equation}

\begin{remark}[불확실성이 어디에 있는가]
$A$의 1행과 3행에는 파라미터가 하나도 없다. 이 두 행은 물리가 아니라 정의
$\dot x_1=x_2$, $\dot x_3=x_4$이고, 수레의 질량이 얼마든 정확히 성립한다. 모델의
불확실성은 전부 2행과 4행에 갇혀 있다. 기계 시스템의 상태공간 모델이 일반적으로
갖는 유용한 성질이다. 운동학을 적은 행은 정확하고, Newton의 법칙을 담은 행만이
우리가 모를 수도 있는 파라미터를 품는다.
\end{remark}

\begin{remark}[댐퍼를 남겨 두면]
$b$를 버리지 않으면 2행과 4행만 바뀐다.
\[
   A_{2\bullet}=\Bigl[-\tfrac{k}{m_1},\;-\tfrac{c+b}{m_1},\;\tfrac{k}{m_1},\;\tfrac{b}{m_1}\Bigr],
   \qquad
   A_{4\bullet}=\Bigl[\tfrac{k}{m_2},\;\tfrac{b}{m_2},\;-\tfrac{k}{m_2},\;-\tfrac{c+b}{m_2}\Bigr].
\]
댐퍼는 대각선의 속도 항에 더해지고, 스프링이 두 위치를 결합하는 것과 똑같은
방식으로 두 속도를 결합한다.
\end{remark}

MATLAB에서는 이 네 행렬을 \verb|sys = ss(A,B,C,D)|로 묶는다.

\subsection{``상태''란 무엇인가}
\begin{definition}[상태]\label{def:state}
\emph{상태} $x(t)$란, 시각 $t$ 이후의 입력 $u(\tau)$ ($\tau\ge t$)가 주어졌을 때
이후의 출력 $y(\tau)$를 모두 결정하기 위해 시각 $t$에 필요한 \emph{최소한의}
정보량이다.
\end{definition}
시각 $t$ 이전에 무슨 일이 있었든, 그것은 오직 $x(t)$를 통해서만 미래에 영향을
준다. 두 수레의 경우 지금 두 수레가 어디에 있고 얼마나 빠른지만 알면 충분하며,
지나온 경로는 알 필요가 없다. 정의에서 진짜 일하는 단어는 \emph{최소한의}다.
상태가 여기서 5차원이 아니라 4차원인 이유가 그것이다.

\begin{remark}\label{rem:nonlinear}
정의~\ref{def:state}에는 선형성에 관한 말이 한 마디도 없다. 실제로 상태공간
표현형은 비선형 시스템에도 그대로 쓸 수 있다(\ref{sec:nonlinear}절). 전달함수는
그렇지 않다. 전달함수는 선형 시스템에만 존재한다.
\end{remark}

\subsection{표준형의 수치 적분}
\eqref{eq:std}이 1계 연립방정식이기 때문에 컴퓨터는 이것을 그대로 앞으로 굴릴 수
있다. 그중 가장 소박하고, 이 과목이 출발점으로 삼는 방법이 \emph{전진 오일러}
(forward Euler)다. 작은 스텝 $h>0$에 대해
\begin{equation}\label{eq:euler}
   x(t+h) \;\approx\; x(t) + h\,\bigl(Ax(t)+Bu(t)\bigr),
\end{equation}
비선형 우변이면 $x(t+h)\approx x(t)+h\,f(x(t),u(t))$. 이것은 $x$를 $t$ 둘레에서
테일러 전개한 첫 두 항에 지나지 않으므로, 한 스텝에서 생기는 오차는 $h^{2}$
차수다. 그렇다고 $h$를 아무렇게나 고를 수는 없다. 너무 크면 시스템 자체는
발산하지 않는데 시뮬레이션만 발산한다.

\begin{remark}[얼마나 큰 스텝까지]\label{rem:eulerstab}
$h$의 한계는 스칼라 시험방정식 $\dot x=\lambda x$에서 정확히 계산할 수 있다.
참해는 $\operatorname{Re}\lambda<0$이면 언제나 감쇠한다. 오일러 한 스텝은 $x$에
일정한 인수를 곱하므로
\[
   x_{k+1} = (1+h\lambda)\,x_k ,\qquad x_k = (1+h\lambda)^{k}x_0 ,
\]
시뮬레이션이 감쇠하려면 $|1+h\lambda|<1$, 즉 $h\lambda$가 $-1$을 중심으로 하는
단위원판 안에 있어야만 한다. 실수 $\lambda<0$이면 이는
\[
   h < \frac{2}{|\lambda|}
\]
라는 뜻이다. 선형 시스템 $\dot x=Ax$에서는 같은 조건이 $A$의 모든 고윳값에 대해
동시에 성립해야 하므로, 가장 빠른 모드가 스텝을 정한다. 최악의 경우는 허수축
위나 그 가까이에 있는 고윳값, 즉 감쇠가 약한 진동이다. $\sigma$가 작은
$\lambda=-\sigma\pm j\omega$에 대해 조건 $|1+h\lambda|<1$은
$h<2\sigma/(\sigma^{2}+\omega^{2})$가 되고, $\sigma\ll\omega$이면 이 값은 아주
작다. 순허수 $\lambda$는 이 조건을 결코 만족할 수 없다. 전진 오일러는 감쇠 없는
진동을 모두 커지는 진동으로 바꿔 버린다.
\end{remark}

%% ------------------------------------------------------------------ 제3강
\lecturebreak{제3강}{2026년 9월 9일}

%%%%%%%%%%%%%%%%%%%%%%%%%%%%%%%%%%%%%%%%%%%%%%%%%%%%%%%%%%%%%%%%%%%%%%%%%%%%%%
\section{비선형 시스템, 평형점, 선형화}\label{sec:nonlinear}

\subsection{비선형 표준형}
실제 시스템은 대개 선형이 아니다. 앞으로 쓸 일반형은
\begin{equation}\label{eq:nl}
   \dot z = f(z,u),\qquad w = h(z,u)
\end{equation}
이며 $z(t)\in\R^{n}$, $u(t)\in\R^{m}$, $w(t)\in\R^{p}$이고
\[
   f:\R^{n}\times\R^{m}\to\R^{n},\qquad h:\R^{n}\times\R^{m}\to\R^{p}
\]
는 \emph{regular}하다고 --- 아래에서 쓸 미분이 존재할 만큼 매끄럽다고 --- 가정한다.
풀어 쓰면 \eqref{eq:nl}은
\[
   \dot z=\begin{pmatrix}\dot z_1\\ \vdots\\ \dot z_n\end{pmatrix}
        =\begin{pmatrix}f_1(z(t),u(t))\\ \vdots\\ f_n(z(t),u(t))\end{pmatrix}
\]
이다.

\begin{note}[이 절의 기호]
비선형 시스템에는 $(z,u,w)$를, 선형 시스템에는 $(x,u,y)$를 쓴다. 선형화한
모델의 변수는 물리량 자체가 아니라 동작점으로부터의 \emph{편차}가 되기 때문이다.
이 둘을 섞는 것이 이 단원에서 가장 흔한 실수다.
\end{note}

\begin{example}[진자]\label{ex:pend}
질량이 $m$, 길이 $\ell$인 무게 없는 막대에 달린 진자가 있고, 회전축에 점성마찰
$b$가 있으며 토크 $u(t)$를 걸 수 있다고 하자(그림~\ref{fig:pend}). 운동방정식은
\[
   m\ell^{2}\ddot\theta = u(t) - mg\ell\sin\theta - b\,\dot\theta,
   \qquad\text{즉}\qquad
   m\ell^{2}\ddot\theta + b\,\dot\theta + mg\ell\sin\theta = u(t)
\]
이고, $\theta$는 아래쪽 연직 방향에서 잰 각이다. $z_1=\theta$, $z_2=\dot\theta$로
두면 정확히 \eqref{eq:nl}의 꼴이 된다.
\begin{equation}\label{eq:pendstate}
\left\{
\begin{aligned}
   \dot z_1 &= z_2,\\
   \dot z_2 &= -\frac{b}{m\ell^{2}}\,z_2 - \frac{g}{\ell}\sin z_1
               + \frac{1}{m\ell^{2}}\,u,\\
   w &= z_1 .
\end{aligned}\right.
\end{equation}
어려운 것은 오직 $\sin z_1$ 하나다. 나머지는 전부 선형이다.
\end{example}

\begin{figure}[htb]
\centering
\begin{tikzpicture}[>=Stealth,scale=1.25]
  \draw[thick] (-1.45,0) -- (1.45,0);
  \foreach \x in {-1.35,-1.05,...,1.42}{\draw[thin] (\x,0) -- ++(-0.22,0.22);}
  \draw[dashed] (0,0) -- (0,-2.95);
  \draw[thick] (0,0) -- (1.30,-2.25);
  \draw[thick,fill=black!10] (1.30,-2.25) circle (0.28);
  \node[sig] at (1.30,-2.25) {$m$};
  \fill (0,0) circle (2.4pt);
  \draw[->] (0,-1.25) arc (-90:-60:1.25);
  \node at (0.38,-1.46) {$\theta$};
  \node[right,sig] at (1.00,-1.66) {$\ell$};
  \draw[very thick,->] (1.30,-2.62) -- (1.30,-3.38);
  \node[right,sig] at (1.34,-3.02) {$mg$};
  \draw[->,thick] (190:0.88) arc (190:256:0.88);
  \node[left,sig] at (-0.94,-0.52) {$u(t)$};
  \node[left,sig] at (-0.06,-0.24) {$b$};
\end{tikzpicture}
\caption{예제~\ref{ex:pend}의 진자. $\theta$는 아래쪽 연직 방향에서 재고, $b$는
회전축의 점성마찰, $u(t)$는 가해 주는 토크다.}
\label{fig:pend}
\end{figure}

\subsection{평형점} \eqref{eq:nl}을 다루는 요령은 아무 일도 일어나지 않는 점
근처에서 보는 것이다.

\begin{definition}[평형점]\label{def:eq}
$f(z^{*},u^{*}) = 0$을 만족하는 쌍 $(z^{*},u^{*})$을 \emph{평형점}이라 한다.
입력을 상수 $u^{*}$로 고정하고 $z(0)=z^{*}$에서 출발하면 $\dot z=0$이므로
$z(t)\equiv z^{*}$가 영원히 유지된다.
\end{definition}
이 말이 성립하려면 입력이 \emph{상수}여야 한다는 점에 유의하자. 평형점은 언제나
``어떤 특정한 상수 입력에 대한'' 평형점이다.

진자에서 $u^{*}=0$이면 $f=0$으로부터 $z_2^{*}=0$이고 $\sin z_1^{*}=0$이므로
$z_1^{*}=0$(아래로 매달린 상태) 또는 $z_1^{*}=\pi$(위로 선 상태)다. 둘 다
평형점이지만, 손을 떼고 놔둘 수 있는 것은 하나뿐이다.

\subsection{자코비안}
$f:\R^{n}\to\R^{n}$에 대해 (잠시 $u$는 숨기고) \emph{자코비안}은 1계 편도함수를
모두 모은 행렬이다.
\[
   \Dop f \;=\;
   \begin{pmatrix}
     \dfrac{\partial f_1}{\partial x_1} & \dfrac{\partial f_1}{\partial x_2} & \cdots & \dfrac{\partial f_1}{\partial x_n}\\[10pt]
     \vdots & \vdots & \ddots & \vdots\\[6pt]
     \dfrac{\partial f_n}{\partial x_1} & \dfrac{\partial f_n}{\partial x_2} & \cdots & \dfrac{\partial f_n}{\partial x_n}
   \end{pmatrix}
   \;\in\;\R^{n\times n}.
\]
한 점 근처에서 $f$를 가장 잘 근사하는 선형사상, 즉 도함수의 다변수 판이다.

\subsection{선형화}\label{sec:lin}
평형점 $(z^{*},u^{*})$ 둘레에서 $f$를 테일러 전개한다.
\[
   \dot z = f(z,u)
   = \underbrace{f(z^{*},u^{*})}_{=\,0}
   + \underbrace{\left.\frac{\partial f}{\partial z}\right|_{z^{*},u^{*}}}_{=:\,A}
     \underbrace{(z-z^{*})}_{=:\,x}
   + \underbrace{\left.\frac{\partial f}{\partial u}\right|_{z^{*},u^{*}}}_{=:\,B}
     \underbrace{(u-u^{*})}_{=:\,\tilde u}
   + \cdots
\]
상수항이 사라지는 것은 평형점에서 전개했기 때문이다. 평형점이 필요한 이유가
바로 이것이다. $z^{*}$가 상수이므로 $\dot x=\dot z$이고, 고차항을 버리면
\begin{equation}\label{eq:linAB}
   \dot x = A\,x + B\,\tilde u,\qquad
   A = \Dop_{z}f\bigl|_{z^{*},u^{*}},\qquad
   B = \Dop_{u}f\bigl|_{z^{*},u^{*}}
\end{equation}
가 남는다. 출력도 똑같이 다룬다.
$w = h(z,u) = h(z^{*},u^{*}) + \Dop_{z}h|_{*}\,x + \Dop_{u}h|_{*}\,\tilde u + \cdots$
에서 출력 편차를 $y := w - h(z^{*},u^{*})$로 정의하면
\begin{equation}\label{eq:linCD}
   y = C\,x + D\,\tilde u,\qquad
   C = \Dop_{z}h\bigl|_{z^{*},u^{*}},\qquad
   D = \Dop_{u}h\bigl|_{z^{*},u^{*}}
\end{equation}
이다. 결과가 정확히 표준형 \eqref{eq:std}이고, 그래서 선형 시스템용으로 만든
도구를 비선형 시스템에도 --- 국소적으로 --- 그대로 겨눌 수 있다.

\begin{caution}[변수가 무엇을 뜻하는지]
$A,B,C,D$는 \emph{상수} 행렬이다. 평형점에서 값을 매긴 자코비안이지 $z$의
함수가 아니다. 그리고 $x$, $\tilde u$, $y$는 모두 \emph{편차}다.
$x=z-z^{*}$, $\tilde u=u-u^{*}$, $y=w-h(z^{*},u^{*})$. (입력 편차도 물결표를
떼고 그냥 $u$로 적는 경우가 흔하다.) 선형화 모델에 편차 대신 절대 입력을 넣으면
그 뒤가 전부 상수만큼 어긋난다.
\end{caution}

\begin{example}[진자, 이어서]
예제~\ref{ex:pend}에 대해
\[
   \Dop_{z}f = \begin{bmatrix} 0 & 1\\[4pt]
     -\dfrac{g}{\ell}\cos z_1 & -\dfrac{b}{m\ell^{2}}\end{bmatrix},
   \qquad
   \Dop_{u}f = \begin{bmatrix}0\\[4pt] \dfrac{1}{m\ell^{2}}\end{bmatrix},
   \qquad
   x = \begin{bmatrix}z_1\\ z_2\end{bmatrix} - z^{*} .
\]
아래쪽 평형점 $z^{*}=(0,0)$, $u^{*}=0$에서는 $\cos z_1^{*}=1$이므로
\[
   A = \begin{bmatrix}0&1\\[3pt] -\dfrac{g}{\ell} & -\dfrac{b}{m\ell^{2}}\end{bmatrix},
   \qquad
   B = \begin{bmatrix}0\\[3pt] \dfrac{1}{m\ell^{2}}\end{bmatrix},
   \qquad C=\begin{bmatrix}1&0\end{bmatrix},\qquad D=0 .
\]
\end{example}

\begin{example}[진자, 위로 선 평형점]\label{ex:pendup}
$\Dop_z f$의 $(2,1)$ 성분은 $-\frac{g}{\ell}\cos z_1^{*}$이므로 어디서
선형화하느냐에 따라 달라진다. 위로 선 평형점 $z^{*}=(\pi,0)$, $u^{*}=0$에서는
이 성분이 $+g/\ell$이 된다.
\[
   A_{\text{up}} = \begin{bmatrix}0&1\\[3pt] \dfrac{g}{\ell} & -\dfrac{b}{m\ell^{2}}\end{bmatrix},
   \qquad
   \det(\lambda I-A_{\text{up}}) = \lambda^{2}+\frac{b}{m\ell^{2}}\lambda-\frac{g}{\ell}.
\]
이제 특성다항식의 상수항이 음수이므로, $b$가 얼마든 근 하나는 양수다. $b=0$이면
고윳값은 $\pm\sqrt{g/\ell}$이다. 선형화가 불안정하고, 당연히 그래야 한다. 끝으로
세워 둔 진자는 넘어진다. 매달린 평형점에서는 같은 다항식이
$\lambda^{2}+\frac{b}{m\ell^{2}}\lambda+\frac{g}{\ell}$이 되고, 모든 $b>0$에 대해
근의 실수부가 음수다. 비선형 모델 하나, 평형점 둘, 질적으로 다른 선형 모델 둘.
선형화를 적을 때 어느 평형점인지를 언제나 밝혀야 하는 이유가 바로 이것이다.
\end{example}

\subsection{선형 모델이 놓이는 자리}
선형화한 시스템은 따로 떨어진 물건이 아니라 실제 시스템을 들여다보는 창이다.
그림~\ref{fig:lin}이 그 장부 정리다. 들어갈 때 $u^{*}$만큼 옮기고, 나올 때
평형 출력 $h(z^{*},u^{*})$를 빼면 된다.

\begin{figure}[htb]
\centering
\begin{tikzpicture}[>=Stealth]
  \node (sum1) [sumc] {};
  \node (plant) [draw,align=center,inner sep=6pt,minimum height=13mm,
                 right=19mm of sum1] {$\dot z = f(z,u)$\\[2pt] $w = h(z,u)$};
  \node (sum2) [sumc,right=19mm of plant] {};
  \node (in)  [left=14mm of sum1] {$\tilde u$};
  \node (out) [right=14mm of sum2] {$y$};
  \draw[->] (in) -- node[pos=0.62,above,lbl] {$+$} (sum1);
  \draw[->] (sum1) -- node[above,sig] {$u$} (plant);
  \draw[->] (plant) -- node[above,sig] {$w$} (sum2);
  \draw[->] (sum2) -- (out);
  \draw[->] ($(sum1)+(0,13mm)$) node[above,sig] {$u^{*}$}
        -- node[pos=0.72,right,lbl] {$+$} (sum1);
  \draw[->] ($(sum2)+(0,13mm)$) node[above,sig] {$h(z^{*},u^{*})$}
        -- node[pos=0.72,right,lbl] {$-$} (sum2);
  \node[below=2.5mm of plant,lbl] {실제 비선형 시스템};
\end{tikzpicture}
\caption{선형화 모델 $\dot x=Ax+B\tilde u$, $y=Cx+D\tilde u$는 $\tilde u$가 작을 때
여기 그려진 사상 $\tilde u\mapsto y$를 근사한다.}
\label{fig:lin}
\end{figure}

\subsection{위상평면}
2계 자율 시스템 $\dot x_1 = f_1(x_1,x_2)$, $\dot x_2 = f_2(x_1,x_2)$에
대해서는 해의 집합 전체를 한 그림에 담을 수 있다. $(x_1,x_2)$ 평면에 궤적을
그리는 것이다. 이 그림을 \emph{위상평면}(phase portrait)이라 한다. 시간 정보는
그림에 나타나지 않지만 나머지는 다 보인다. 평형점은 벡터장이 0이 되는 점이고,
해의 유일성 때문에 두 궤적은 서로 만나지 않는다. 그림~\ref{fig:phase}는 감쇠
나선의 예다.

\begin{figure}[htb]
\centering
\begin{tikzpicture}[>=Stealth,scale=1.2]
  \foreach \a in {-2.4,-1.6,...,2.5}{
    \foreach \bb in {-2.4,-1.6,...,2.5}{
      \pgfmathsetmacro{\dx}{\bb}
      \pgfmathsetmacro{\dy}{-\a-0.7*\bb}
      \pgfmathsetmacro{\nn}{sqrt(\dx*\dx+\dy*\dy)+0.001}
      \draw[->,gray!50,thin] (\a,\bb) -- ++({0.44*\dx/\nn},{0.44*\dy/\nn});
    }}
  \draw[->] (-3.0,0) -- (3.0,0) node[right] {$x_1$};
  \draw[->] (0,-2.95) -- (0,2.95) node[above] {$x_2$};
  \foreach \A in {2.4,1.3}{
    \draw[thick] plot[domain=0:720,samples=100,variable=\t,smooth]
      ({\A*exp(-0.0065216*\t)*(cos(\t)+0.37364*sin(\t))},
       {-1.0676*\A*exp(-0.0065216*\t)*sin(\t)});}
  \foreach \A/\T in {2.4/40, 2.4/180, 1.3/110}{
    \pgfmathsetmacro{\xa}{\A*exp(-0.0065216*\T)*(cos(\T)+0.37364*sin(\T))}
    \pgfmathsetmacro{\ya}{-1.0676*\A*exp(-0.0065216*\T)*sin(\T)}
    \pgfmathsetmacro{\xb}{\A*exp(-0.0065216*(\T+16))*(cos(\T+16)+0.37364*sin(\T+16))}
    \pgfmathsetmacro{\yb}{-1.0676*\A*exp(-0.0065216*(\T+16))*sin(\T+16)}
    \draw[->,thick] (\xa,\ya) -- (\xb,\yb);}
  \foreach \A in {2.4,1.3}{\fill (\A,0) circle (1.7pt);}
  \node[above right,lbl] at (2.42,0.04) {$x(0)$};
  \fill (0,0) circle (2.4pt);
\end{tikzpicture}
\caption{$\dot x_1=x_2$, $\dot x_2=-x_1-0.7x_2$의 위상평면. 회색 화살표가
벡터장이고 굵은 두 곡선이 궤적이다. 원점의 검은 점이 평형점이며 두 궤적 모두
그리로 감겨 들어간다.}
\label{fig:phase}
\end{figure}

\emph{선형} 2계 시스템 $\dot x = Ax$에서는 위상평면 전체가 $A$의 두 고윳값
$\lambda_{1,2}$로 정해지고, 가능한 그림은 몇 가지뿐이다.

\begin{center}\small
\renewcommand{\arraystretch}{1.3}
\begin{tabular}{@{}lll@{}}
\toprule
고윳값 & 위상평면의 모양 & 평형점 근처의 거동\\
\midrule
실수, 둘 다 음수          & 안정 노드(node)     & 모든 궤적이 돌지 않고 들어온다\\
실수, 둘 다 양수          & 불안정 노드   & 모든 궤적이 빠져나간다\\
실수, 부호가 반대         & 안장점(saddle)          & 한 방향으로 들어와 다른 방향으로 나간다\\
복소수, $\operatorname{Re}\lambda<0$ & 안정 초점(focus, 나선) & 궤적이 감겨 들어간다\\
복소수, $\operatorname{Re}\lambda>0$ & 불안정 초점 & 궤적이 감겨 나간다\\
순허수             & 중심(centre)          & 닫힌 궤도, 들어가지도 나가지도 않는다\\
\bottomrule
\end{tabular}
\end{center}

\noindent 고윳값이 실수이면 고유벡터를 따라가는 직선 궤적이 그림의 뼈대다. 해는
$x(t)=c_1e^{\lambda_1t}v_1+c_2e^{\lambda_2t}v_2$이고, 각 성분이 제 고유벡터를
따라 커지거나 줄어든다(그림~\ref{fig:portraits}). 복소수 쌍 $\sigma\pm j\omega$이면
해는 주파수 $\omega$의 회전을 $e^{\sigma t}$ 포락선이 감싼 것이다. 정확히
\ref{sec:pfe}절의 부분분수 항들이다.

\begin{figure}[htb]
\centering
\begin{tikzpicture}[>=Stealth]
  \foreach \px/\ma/\mb/\mc/\md/\name in {0/-1/0/0/-2.5/안정 노드,
                                     3.35/1/0/0/-1.5/안장점,
                                     6.7/-0.35/1/-1/-0.35/안정 초점,
                                     10.05/0/1/-1/0/중심}{
    \begin{scope}[xshift=\px cm,scale=0.62]
      \draw[gray!60] (-2,0) -- (2,0); \draw[gray!60] (0,-2) -- (0,2);
      \foreach \u in {-1.6,-0.8,...,1.7}{
        \foreach \w in {-1.6,-0.8,...,1.7}{
          \pgfmathsetmacro{\dx}{\ma*\u+\mb*\w}
          \pgfmathsetmacro{\dy}{\mc*\u+\md*\w}
          \pgfmathsetmacro{\nn}{sqrt(\dx*\dx+\dy*\dy)+0.001}
          \draw[->,thin] (\u,\w) -- ++({0.34*\dx/\nn},{0.34*\dy/\nn});
        }}
      \fill (0,0) circle (2.2pt);
      \node[lbl] at (0,-2.45) {\name};
    \end{scope}}
\end{tikzpicture}
\caption{표의 선형 위상평면 가운데 넷을 $\dot x=Ax$의 벡터장으로 그린 것.
안정 노드($A=\operatorname{diag}(-1,-2.5)$), 안장점
($\operatorname{diag}(1,-1.5)$), 안정 초점(고윳값 $-0.35\pm j$), 중심(고윳값
$\pm j$).}
\label{fig:portraits}
\end{figure}

\begin{remark}[진자로 돌아가서]
선형화는 비선형 시스템이 각 평형점 \emph{근처}에서 어떤 그림을 보이는지 알려
준다. 예제~\ref{ex:pend}의 진자에서 매달린 평형점의 특성다항식은
$\lambda^{2}+\frac{b}{m\ell^{2}}\lambda+\frac g\ell$로, 마찰이 약하면 안정 초점,
강하면 안정 노드다. 위로 선 평형점은
$\lambda^{2}+\frac{b}{m\ell^{2}}\lambda-\frac g\ell$이고 그 근은 부호가 반대인
실수이므로 안장점이다. 비선형 위상평면 하나에 두 그림이 모두 들어 있다.
$\theta=0$의 나선형 흡인점과 $\theta=\pi$의 안장점이, 둘 모두에서 먼 곳에서
이어 붙어 있다.
\end{remark}

%% ================================================================ 제3부
\partdivider{제3부 --- 전달함수}

\lecturebreak{제4강}{2026년 9월 14일}

%%%%%%%%%%%%%%%%%%%%%%%%%%%%%%%%%%%%%%%%%%%%%%%%%%%%%%%%%%%%%%%%%%%%%%%%%%%%%%
\section{라플라스 변환과 전달함수}

\subsection{라플라스 변환: 필요한 만큼만}
신호 $f(t)$의 (단측) 라플라스 변환은
\begin{equation}\label{eq:laplace}
   F(s) = \Lap\{f\}(s) = \int_{0^{-}}^{\infty} f(t)\,e^{-st}\,\d t ,
   \qquad s\in\mathbb{C}
\end{equation}
이다. 적분 하한을 $0^{-}$로 적는 이유는, 원점에 놓인 임펄스와 초기 상태까지
세기 위해 $t=0$ \emph{이전}부터 적분을 시작하기 때문이다.

\smallskip
\begin{center}
\renewcommand{\arraystretch}{2.0}
\begin{tabular}{@{}l@{\qquad\qquad}l@{\qquad\qquad}l@{}}
\toprule
$\delta(t) \;\mapsto\; 1$
  & $u_{s}(t) \;\mapsto\; \dfrac{1}{s}$
  & $t\,u_{s}(t) \;\mapsto\; \dfrac{1}{s^{2}}$\\[6pt]
$t^{n}u_{s}(t) \;\mapsto\; \dfrac{n!}{s^{\,n+1}}$
  & $e^{-\alpha t} \;\mapsto\; \dfrac{1}{s+\alpha}$
  & $f(t-a)\,u_{s}(t-a) \;\mapsto\; e^{-as}F(s)$\\[6pt]
$\sin\omega t \;\mapsto\; \dfrac{\omega}{s^{2}+\omega^{2}}$
  & $\cos\omega t \;\mapsto\; \dfrac{s}{s^{2}+\omega^{2}}$
  & $f*g \;\mapsto\; F(s)\,G(s)$\\
\bottomrule
\end{tabular}
\end{center}

\smallskip\noindent
여기서 $u_{s}$는 단위계단함수다. 변환 쌍은 도구의 절반일 뿐이다. 나머지 절반은
변환이 대수로 바꿔 주는 \emph{연산}의 목록이다.

\begin{center}\small
\renewcommand{\arraystretch}{1.55}
\begin{tabular}{@{}lll@{}}
\toprule
& 시간 영역 & $s$ 영역\\
\midrule
선형성        & $a f(t)+b g(t)$            & $aF(s)+bG(s)$\\
미분  & $\dot f(t)$                & $sF(s)-f(0^{-})$\\
적분      & $\int_{0}^{t} f(\tau)\,\d\tau$ & $F(s)/s$\\
시간 이동       & $f(t-a)\,u_s(t-a)$, $a>0$  & $e^{-as}F(s)$\\
주파수 이동  & $e^{-\alpha t}f(t)$        & $F(s+\alpha)$\\
$t$ 곱하기 & $t\,f(t)$             & $-F'(s)$\\
합성곱      & $(f*g)(t)=\int_0^t f(\tau)g(t-\tau)\,\d\tau$ & $F(s)\,G(s)$\\
\bottomrule
\end{tabular}
\end{center}

이 가운데 실제로 일을 하는 것은 미분 규칙이다.
\begin{equation}\label{eq:diffrule}
   \Lap\{\dot f\} = sF(s) - f(0),
   \qquad
   \Lap\{\ddot f\} = s^{2}F(s) - s f(0) - \dot f(0).
\end{equation}
외울 것은 그 결론 한 줄이다. \emph{초기조건이 모두 0이면 미분이 $s$ 곱하기로
바뀐다.} 그리고 적분 규칙에 따라 적분은 $s$ 나누기로 바뀐다. 이 사실 하나가
선형 미분방정식을 대수 방정식으로 바꾸고, 전달함수가 존재하는 이유 전부가
여기에 있다.

\begin{remark}[표 읽기]
두 줄은 그림으로 볼 만하다. 주파수 이동은 신호를 $e^{-\alpha t}$로 감쇠시키면
그 변환이 $s$-평면에서 왼쪽으로 $\alpha$만큼 미끄러진다는 뜻이다.
$\sin\omega t$의 극점은 $\pm j\omega$에 있고, $e^{-\alpha t}\sin\omega t$의
극점은 $-\alpha\pm j\omega$에 있다. 시간 이동은 $a$초의 순수한 지연이
$e^{-as}$를 곱한다는 뜻이다. 이 인수는 다항식의 비가 아니고, 그래서 지연은
유리 전달함수 위에 세운 모든 것과 잘 어울리지 않는다.
\end{remark}

\subsection{초기값과 최종값}\label{sec:fvt}
극한 두 개만 있으면 신호를 역변환하지 않고도 변환에서 곧바로 그 시작과 끝을
읽을 수 있다.

\begin{proposition}[초기값 정리와 최종값 정리]\label{prop:fvt}
$f$의 라플라스 변환을 $F$라 하자.
\begin{enumerate}[label=(\roman*),itemsep=2pt]
\item $F$가 strictly proper이면 $f(0^{+}) = \lim_{s\to\infty} sF(s)$이다.
\item $sF(s)$의 모든 극점의 실수부가 음수이면 $t\to\infty$일 때 $f(t)$의 극한이
존재하고
\[
   \lim_{t\to\infty} f(t) = \lim_{s\to 0} sF(s)
\]
이다.
\end{enumerate}
\end{proposition}

둘 다 미분 규칙에서 나온다. $\int_{0}^{\infty}\dot f(t)e^{-st}\,\d t
= sF(s)-f(0)$에서 $s\to0$으로 보내면(적분이 $f(\infty)-f(0)$이 된다), 또는
$s\to\infty$로 보내면(적분이 사라진다) 두 명제가 나온다.

\begin{caution}[최종값 정리의 가정]
극점에 대한 조건은 장식이 아니다. $F(s)=1/(s-1)$이면 공식은
$\lim_{s\to0}sF(s)=0$을 주지만 $f(t)=e^{t}$는 한없이 커진다. 마찬가지로
$F(s)=\omega/(s^{2}+\omega^{2})$도 $0$을 주지만 $\sin\omega t$에는 극한 자체가
없다. 이 정리는 존재하는 최종값을 알려 줄 뿐, 최종값이 존재하는지는 판정하지
못한다. 그 질문에 답하는 것은 극점이고, 그것이 안정도의 주제다.
\end{caution}

\begin{example}[붙들린 물체의 정상 오차]
1.3절의 되먹임 법칙은 \eqref{eq:massfb},
$m\ddot y+c\dot y+k(y-h)=(\hat m-m)g$를 낳았다. $e=y-h$로 두자. $e(0)=\dot e(0)=0$이고
우변이 상수 $d=(\hat m-m)g$이므로
\[
   E(s) = \frac{d}{s\,(ms^{2}+cs+k)} ,
   \qquad
   \lim_{t\to\infty} e(t) = \lim_{s\to0} sE(s) = \frac{d}{k}
\]
이다. $sE(s)$의 극점은 $ms^{2}+cs+k$의 근이고 $c,k>0$이면 좌반평면에 있으므로
정리를 쓸 수 있다. 미분방정식을 풀지 않고도 \eqref{eq:offset}이 그대로 다시
나온다.
\end{example}

\subsection{전달함수}
\begin{definition}[전달함수]\label{def:tf}
입력 $u(t)$, 출력 $y(t)$인 선형 시불변(LTI) 시스템이 \emph{정지 상태에서
출발}(initially at rest, 모든 초기조건이 0)할 때, \emph{전달함수}는
\[
   P(s) \;=\; \frac{\Lap\{y(t)\}}{\Lap\{u(t)\}} \;=\; \frac{Y(s)}{U(s)}
\]
이다.
\end{definition}

\begin{figure}[htb]
\centering
\begin{tikzpicture}[>=Stealth]
  \node (b) [draw,minimum width=24mm,minimum height=12mm] {$P(s)$};
  \draw[->,thick] ($(b.west)+(-24mm,0)$) node[left] {$u(t)$}
        -- node[above,sig] {$U(s)$} (b.west);
  \draw[->,thick] (b.east)
        -- node[above,sig] {$Y(s)$} ($(b.east)+(24mm,0)$) node[right] {$y(t)$};
\end{tikzpicture}
\caption{블록으로 본 전달함수. 시스템이 정지 상태에서 출발할 때
$Y(s)=P(s)U(s)$가 성립한다.}
\label{fig:tfblock}
\end{figure}

같은 말로, $P(s)$는 임펄스 응답의 라플라스 변환이다. $u=\delta$를 넣으면
$U\equiv1$이므로 $Y=P$가 되기 때문이다. 신호가 여럿 등장할 때는 이것을
$T_{uy}(s)$로 적어 어느 신호가 입력이고 어느 것이 출력인지 표시하기도 한다.
단일입력 단일출력이면 그냥 $P(s)$로 써도 헷갈릴 일이 없다.

\begin{remark}[``정지 상태에서 출발''이 왜 사소하지 않은가]
이 가정이 없으면 변환된 식에 초기조건에서 오는 $-sy(0)-\dot y(0)-\cdots$ 항이
남고, 그러면 $Y(s)/U(s)$는 시스템만이 아니라 ``시스템이 어떻게 출발했는지''에도
의존하게 된다. 플랜트의 성질이 아니게 되므로 그것을 시스템의 전달\emph{함수}라고
부를 수 없다.
\end{remark}

\begin{example}\label{ex:tf}
$\ddot y + 3\dot y + y = 2\dot u + u$를 보자. 영초기조건에서 변환하면
$s^{2}Y+3sY+Y = 2sU+U$이므로 $P(s) = (2s+1)/(s^{2}+3s+1)$이다.
\end{example}
요령은 예제에 그대로 보인다.
$a_n y^{(n)}+\cdots+a_1\dot y+a_0y = b_m u^{(m)}+\cdots+b_1\dot u+b_0u$에서
$\d/\d t$를 전부 $s$로 바꾸고 $Y/U$에 대해 풀면
\[
   P(s) = \frac{b_ms^{m}+\cdots+b_1s+b_0}{a_ns^{n}+\cdots+a_1s+a_0}
\]
이다. 분모의 차수가 시스템의 차수다.

\subsection{예제: two-cart 시스템 (이상화)}\label{sec:twocarttf}
Two-cart 시스템에서 마찰과 결합 댐퍼를 뺀($c=b=0$) 경우를 보자.
\begin{equation}\label{eq:twocartideal}
   m_1\ddot p_1 = k\,(p_2-p_1) + u,\qquad
   m_2\ddot p_2 = k\,(p_1-p_2),\qquad y=p_2 .
\end{equation}
영초기조건에서 변환하고 $Y=P_2$로 쓰면
\[
   m_1 s^{2}P_1 = kY - kP_1 + U,\qquad
   m_2 s^{2}Y = kP_1 - kY
\]
이다. 두 번째 식에서 $P_1 = \dfrac{m_2s^{2}+k}{k}\,Y$이므로, 이것을 첫 번째 식
$(m_1s^{2}+k)P_1 = kY+U$에 넣으면
\[
   \frac{(m_1s^{2}+k)(m_2s^{2}+k)}{k}\,Y = kY + U
   \quad\Longrightarrow\quad
   \Bigl[(m_1s^{2}+k)(m_2s^{2}+k) - k^{2}\Bigr]Y = k\,U
\]
가 된다. $(m_1s^{2}+k)(m_2s^{2}+k)-k^{2} = s^{2}\bigl(m_1m_2s^{2}+k(m_1+m_2)\bigr)$
이므로
\begin{equation}\label{eq:twocartP}
   \boxed{\;
   \frac{Y(s)}{U(s)} \;=\; \frac{k}{\,s^{2}\bigl(m_1m_2\,s^{2}+k(m_1+m_2)\bigr)\,}\;}
\end{equation}
이다.

\begin{remark}[\eqref{eq:twocartP}을 읽는 법]
분모의 $s^{2}$은 자유 강체 모드다. 마찰이 없으니 두 수레를 지면에 붙들어 두는
것이 아무것도 없고, 알짜 충격량이 들어가면 둘은 영원히 미끄러져 간다. 남은 인수는
$s = \pm j\sqrt{k(m_1+m_2)/(m_1m_2)}$에서 0이 되는데, 이것이 두 수레가 서로에
대해 진동하는 주파수다. 둘 다 허수축 위에 있다. $c=b=0$이면 감쇠시키는 것이
없기 때문이고, 2절의 모델이 $c$와 $b$를 남겨 둔 이유가 바로 이것이다.
\end{remark}

\begin{caution}[two-cart 모델이 두 개다]
\eqref{eq:twocartideal}은 \eqref{eq:eom}이 \emph{아니다}. 계산을 읽을 만하게
하려고 마찰 $c$와 댐퍼 $b$를 뺀 것이다. 두 모델에서 나온 결과를 섞어 쓰면 안
된다. $c>0$이면 \eqref{eq:twocartP}의 $s^{2}$ 인수가 $s(\,\cdot\,)$이 되어,
극점 하나가 원점을 떠나고 자유 적분기는 하나만 남는다.
\end{caution}

\begin{remark}[$c$와 $b$를 남긴 채 같은 계산을 하면]\label{rem:fullTF}
전체 모델 \eqref{eq:eom}을 그대로 끌고 가면 행렬 꼴로
\[
   \begin{bmatrix}
     m_1s^{2}+(b+c)s+k & -(bs+k)\\[3pt]
     -(bs+k) & m_2s^{2}+(b+c)s+k
   \end{bmatrix}
   \begin{bmatrix}P_1\\[3pt] P_2\end{bmatrix}
   = \begin{bmatrix}U\\[3pt] 0\end{bmatrix}
\]
이 되고,
$\Delta(s) = \bigl[m_1s^{2}+(b+c)s+k\bigr]\bigl[m_2s^{2}+(b+c)s+k\bigr]-(bs+k)^{2}$
로 두면 $Y(s)/U(s) = (bs+k)/\Delta(s)$다. $b=c=0$을 넣으면 \eqref{eq:twocartP}이
정확히 복원되므로, 양쪽을 서로 검산해 주는 셈이다. 눈여겨볼 것이 둘 있다. 결합
댐퍼는 $s=-k/b$에 유한 영점을 만든다. 즉 분자는 $bs+k$이지 결코 $0$이 아니다.
그리고 $c=0$이면서 $b>0$이면
$\Delta(s)=s^{2}\bigl[m_1m_2s^{2}+(m_1+m_2)bs+(m_1+m_2)k\bigr]$로 원점의 이중극점이
살아남는다. 그 극점은 댐퍼가 없어서가 아니라 \emph{지면} 마찰이 없어서 생기는
것이기 때문이다.
\end{remark}

\subsection{상태공간에서 전달함수로}\label{sec:ss2tf}
그림~\ref{fig:routes}의 두 표현형은 계산 한 번으로 연결된다. $x(0)=0$인
\eqref{eq:std}에서 출발해 변환하면 $s\,X(s) = A\,X(s) + B\,U(s)$이다. $X$ 항을
모으면 $(sI-A)X(s) = B\,U(s)$이고, 역행렬을 취하면
$X(s) = (sI-A)^{-1}B\,U(s)$이다. 이것을 $Y = CX + DU$에 넣으면
\begin{equation}\label{eq:ss2tf}
   \boxed{\;Y(s) = \bigl[\,C\,(sI-A)^{-1}B + D\,\bigr]\,U(s),
   \qquad P(s) = C(sI-A)^{-1}B + D. \;}
\end{equation}
SISO 시스템이면 이것은 $s$에 대한 스칼라 유리함수이고, 일반적으로는 유리함수를
성분으로 갖는 $p\times m$ 행렬이다.

\begin{remark}[$x(0)$을 남겨 두면]
초기조건을 0으로 두지 않고 한 번 계산해 보면, 영초기조건 가정이 정확히 무엇을
사 주는지가 드러난다. $\dot x = Ax+Bu$를 그대로 변환하면
$sX(s)-x(0) = AX(s)+BU(s)$, 즉 $(sI-A)X(s) = x(0) + BU(s)$이므로
\[
   Y(s) = \underbrace{C(sI-A)^{-1}x(0)}_{\text{초기값에 의한 응답}}
        + \underbrace{\bigl[C(sI-A)^{-1}B+D\bigr]U(s)}_{\text{입력에 의한 응답}}
\]
가 된다. 전달함수는 두 번째 항만을 가리킨다. 첫 번째 항이 현실에서 사라지는
것이 아니라, 정의에서 떼어낸 것일 뿐이다. 정의~\ref{def:tf}가 ``정지 상태에서
출발''을 명시해야 하는 이유가 이것이다.
\end{remark}

\begin{remark}
$(sI-A)^{-1} = \operatorname{adj}(sI-A)/\det(sI-A)$이므로
$P(s) = C\,\operatorname{adj}(sI-A)\,B/\det(sI-A) + D$로도 쓸 수 있고, 따라서
$P(s)$의 모든 극점은 $A$의 고윳값이다. 역은 성립하지 않는다. 약분되어 사라질 수
있기 때문이다. 즉 상태공간 모델은 전달함수가 담지 못하는 정보를 가질 수 있고,
이 과목이 전달함수만 쓰지 않는 이유 중 하나가 그것이다. 반대 방향, 즉 $P(s)$에서
사중쌍 $(A,B,C,D)$을 만드는 것은 \emph{실현}(realization)이라 부른다. $P$가
proper이면 언제나 가능하지만 답이 전혀 유일하지 않다. 그림~\ref{fig:routes}에서
두 상자 사이 화살표가 양방향인 이유다.
\end{remark}

\begin{example}[실현, 한 단계씩]\label{ex:realization}
$P(s)=\dfrac{s^{2}+s+2}{s^{3}+3s^{2}+2s-1}$을 실현해 보자. 중간 신호 $\bar y$를
두어 둘로 나눈다.
\[
   U \;\longrightarrow\; \boxed{\dfrac{1}{s^{3}+3s^{2}+2s-1}} \;\longrightarrow\;
   \bar Y \;\longrightarrow\; \boxed{s^{2}+s+2} \;\longrightarrow\; Y .
\]
첫 블록은 미분방정식 $\dddot{\bar y}+3\ddot{\bar y}+2\dot{\bar y}-\bar y=u$이고
입력의 미분이 없으므로, 3.1절처럼 $x_1=\bar y$, $x_2=\dot{\bar y}$,
$x_3=\ddot{\bar y}$로 고르면
\[
   \dot x_1=x_2,\qquad \dot x_2=x_3,\qquad \dot x_3 = x_1-2x_2-3x_3+u
\]
가 된다. 두 번째 블록은 $\bar y$를 미분하는 것뿐이고, 필요한 미분은 이미 모두
상태다. $y=\ddot{\bar y}+\dot{\bar y}+2\bar y = 2x_1+x_2+x_3$. 따라서
\[
   A=\begin{bmatrix}0&1&0\\0&0&1\\1&-2&-3\end{bmatrix},\qquad
   B=\begin{bmatrix}0\\0\\1\end{bmatrix},\qquad
   C=\begin{bmatrix}2&1&1\end{bmatrix},\qquad D=0 .
\]
이 규칙은 일반적이다. 분모의 계수가 (부호가 뒤집혀) $A$의 마지막 행을 채우고,
분자의 계수가 $C$가 된다. 검산하면
$(sI-A)^{-1}B=(1,\;s,\;s^{2})\tp/(s^{3}+3s^{2}+2s-1)$이고, 여기에 $C$를 곱하면
$P(s)$다.
\end{example}

%% ------------------------------------------------------------------ 제5강
\lecturebreak{제5강}{2026년 9월 16일}

\subsection{극점, 영점, properness}\label{sec:polezero}
앞으로 만날 전달함수는 모두 다항식의 비, 즉
\begin{equation}\label{eq:rational}
   G(s) = \frac{b_ms^{m}+\cdots+b_1s+b_0}
               {s^{n}+a_{n-1}s^{n-1}+\cdots+a_1s+a_0}
\end{equation}
꼴의 \emph{유리}(rational) 전달함수다. 분모의 차수 $n$을 $G$의 \emph{차수}라
하고, $n-m$을 \emph{상대차수}(relative degree)라 한다. $n$과 $m$을 견주면 이름이
넷 나온다.

\begin{center}\small
\begin{tabular}{@{}ll@{}}
\toprule
$n \ge m$ & proper\\
$n > m$   & strictly proper\\
$n = m$   & biproper\\
$n < m$   & improper\\
\bottomrule
\end{tabular}
\end{center}

\begin{remark}[improper 가 왜 반갑지 않은가]
improper 하다는 것은 출력 쪽 미분 횟수보다 입력 쪽 미분 횟수가 많다는 뜻이고,
그것을 실제로 구현하려면 입력을 \emph{미분}해야 한다. 순수 미분기는
$|j\omega|=\omega$이므로 주파수가 올라갈수록 이득이 한없이 커진다. 정작 없애고
싶은 고주파 측정잡음을 그만큼 증폭한다는 뜻이다. 실제 플랜트는 strictly proper
인데, 이는 즉각적으로 반응하는 것은 없다는 말과 같다. 우리가 직접 만드는
제어기도 같은 이유로 적어도 proper로 유지한다.
\end{remark}

분자와 분모를 인수분해하면
\begin{equation}\label{eq:factored}
   G(s) = \frac{N(s)}{D(s)}
        = k\,\frac{\prod_{i=1}^{m}(s-z_i)}{\prod_{i=1}^{n}(s-p_i)}
\end{equation}
이고 $N$, $D$는 실계수 다항식이다. $D(s)=0$이 되는 $s\in\mathbb{C}$가
\textbf{극점}(pole)으로 $n$개, $N(s)=0$이 되는 $s$가 \textbf{영점}(zero)으로
$m$개다. $N$과 $D$의 계수가 \emph{실수}이므로, 실수가 아닌 극점과 영점은 반드시
켤레복소수와 짝을 이루어 함께 나타난다. 이들을 복소평면에 극점은 $\times$,
영점은 $\circ$로 찍은 것이 \emph{pole--zero map}이다(그림~\ref{fig:pzmap}).

\begin{figure}[htb]
\centering
\begin{tikzpicture}[>=Stealth,scale=0.95,
   polex/.pic={\draw[thick] (-0.13,-0.13)--(0.13,0.13) (-0.13,0.13)--(0.13,-0.13);}]
  \draw[->] (-3.5,0) -- (1.7,0) node[right] {$\operatorname{Re}s$};
  \draw[->] (0,-2.1) -- (0,2.2) node[above] {$\operatorname{Im}s$};
  \pic at (-2.7,0) {polex};
  \pic at (-1.15,1.35) {polex};
  \pic at (-1.15,-1.35) {polex};
  \draw[thick] (-0.40,0.75) circle (0.14);
  \draw[thick] (-0.40,-0.75) circle (0.14);
  \draw[dashed,gray] (-1.15,1.35) -- (-1.15,-1.35);
  \draw[dashed,gray] (-0.40,0.75) -- (-0.40,-0.75);
  \node[anchor=west,lbl] at (0.30,1.80) {$\times$ 극점 \qquad $\circ$ 영점};
\end{tikzpicture}
\caption{Pole--zero map. $N$과 $D$의 계수가 실수이므로 실수가 아닌 극점과 영점은
실축에 대해 대칭인 짝과 함께 나타난다(점선). 실수인 것은 혼자 선다.}
\label{fig:pzmap}
\end{figure}

\subsection{부분분수로 역변환하기}\label{sec:pfe}
$Y(s)$에서 $y(t)$로 돌아갈 때 역변환 적분을 직접 계산하는 일은 거의 없다. 대신
유리함수 $Y(s)$를 간단한 항들의 합으로 쪼개는데, 각 항이 변환 쌍 표의 한 줄이다.
$Y$가 strictly proper이고 극점이 $p_i$이면 항은 세 종류다.

\smallskip\noindent\textbf{서로 다른 실수 극점.} 단순 극점 $p_i$ 하나가
$A_i/(s-p_i)$를 내놓고, 그 계수는 \emph{가림 규칙}(cover-up rule)으로 구한다.
$(s-p_i)$를 곱하고 $s=p_i$를 넣으면 된다.
\begin{equation}\label{eq:coverup}
   A_i = \bigl[(s-p_i)\,Y(s)\bigr]_{s=p_i} .
\end{equation}
이런 항은 각각 $A_i e^{p_i t}$로 역변환된다.

\smallskip\noindent\textbf{중복 극점.} 중복도 $r$인 극점은 항 $r$개,
$B_1/(s-p)+B_2/(s-p)^{2}+\dots+B_r/(s-p)^{r}$를 내놓는다. 맨 위 계수 $B_r$은
$(s-p)^{r}$로 가림 규칙을 써서 구하고, 아래 계수들은 미분해서 구한다.
$B_{r-k} = \frac{1}{k!}\frac{\d^{k}}{\d s^{k}}
\bigl[(s-p)^{r}Y(s)\bigr]_{s=p}$. 항 $B/(s-p)^{j}$는
$B\,t^{j-1}e^{pt}/(j-1)!$로 역변환된다. 중복 극점은 지수함수 앞에 $t$의 다항식을
붙인다.

\smallskip\noindent\textbf{복소수 쌍.} 극점 $-\alpha\pm j\omega$는 한데 묶어
실수 2차식 $(s+\alpha)^{2}+\omega^{2}$ 하나로 두고, 분자는 1차식으로 둔다. 이를
표의 이동된 두 쌍으로 맞춰 쓰면
\[
   \frac{C(s+\alpha)+D\omega}{(s+\alpha)^{2}+\omega^{2}}
   \;\longmapsto\;
   e^{-\alpha t}\bigl(C\cos\omega t + D\sin\omega t\bigr)
\]
로, ($\alpha>0$이면) 감쇠하는 진동이 나온다. 극점의 실수부가 포락선을, 허수부가
주파수를 정한다.

\begin{example}[세 경우]\label{ex:pfe}
\begin{enumerate}[label=(\alph*),itemsep=3pt,leftmargin=*]
\item $Y(s)=\dfrac{s+5}{(s+1)(s+2)}=\dfrac{A}{s+1}+\dfrac{B}{s+2}$. 가림 규칙으로
$A=\frac{s+5}{s+2}\big|_{s=-1}=4$, $B=\frac{s+5}{s+1}\big|_{s=-2}=-3$이므로
$y(t)=4e^{-t}-3e^{-2t}$.
\item $Y(s)=\dfrac{1}{s(s+1)^{2}}=\dfrac{A}{s}+\dfrac{B_1}{s+1}+\dfrac{B_2}{(s+1)^{2}}$.
가림 규칙으로 $A=1$, $B_2=\frac1s\big|_{s=-1}=-1$이고, 미분하면
$B_1=\frac{\d}{\d s}\frac1s\big|_{s=-1}=-1$이다. 따라서
$y(t)=1-e^{-t}-t\,e^{-t}$.
\item $Y(s)=\dfrac{1}{s^{2}+2s+5}=\dfrac{1}{(s+1)^{2}+2^{2}}
=\dfrac12\cdot\dfrac{2}{(s+1)^{2}+2^{2}}$이므로 $y(t)=\tfrac12e^{-t}\sin 2t$.
\end{enumerate}
(b)에서는 \ref{sec:fvt}절의 두 정리로 답을 검산할 수 있다.
$\lim_{s\to\infty}sY(s)=0=y(0^{+})$이고 $\lim_{s\to0}sY(s)=1=y(\infty)$이다.
\end{example}

\subsection{출력의 극점은 어디서 오는가}\label{sec:outputpoles}
플랜트를 손에 넣었으면 입력에 대한 응답은 $y(t) = \Lap^{-1}\{G(s)\,U(s)\}$이고,
이 역변환을 읽는 방법은 부분분수로 쪼개는 것이다. 예컨대
\begin{equation}\label{eq:pfe}
   Y(s) = \frac{A}{s-1} + \frac{B}{s-2} + \frac{C}{s+3}
        + \frac{Ds+E}{s^{2}+\omega^{2}}
\end{equation}
처럼 나왔다고 하자. 앞의 세 분모는 \emph{플랜트}의 극점이고, 마지막 하나는
\emph{입력}의 극점이다. 입력이 $\sin\omega t$였다면 정확히 저 인수가 붙는다. 항별로
역변환하면 극점 하나당 시간함수 하나가 나오므로,
\begin{center}
\emph{출력의 극점은 플랜트의 극점과 입력의 극점을 합한 것이다.}
\end{center}
플랜트 자신의 극점은 무엇으로 구동하든 반드시 나타난다. 그 극점이 제자리에
있는지를 따지는 일이 입력과 무관하게 플랜트만의 문제인 이유가 이것이다.

\subsection{블록선도와 되먹임 루프}\label{sec:blocks}
$Y(s)=P(s)U(s)$가 그냥 곱셈이므로, 서로 연결된 시스템은 대수로 합쳐진다. 거의
모든 블록선도는 세 가지 연결로 설명된다(그림~\ref{fig:blocks}).
\begin{equation}\label{eq:blockalg}
   \underbrace{\;Q(s)P(s)\;}_{\text{직렬}},
   \qquad
   \underbrace{\;P(s)+Q(s)\;}_{\text{병렬}},
   \qquad
   \underbrace{\;\frac{G(s)}{1+G(s)H(s)}\;}_{\text{음의 되먹임}} .
\end{equation}
앞의 둘은 바로 보인다. 중요한 것은 세 번째이고, 한 번은 유도해 둘 만하다.
그림~\ref{fig:blocks}(c)의 루프에서 오차는 $E=R-HY$, 출력은 $Y=GE$이므로
$Y=G(R-HY)$다. 출력이 양변에 나타나는데, 이것이 되먹임이 대수에 남긴 흔적이다.
정리하면 $(1+GH)Y=GR$이고 공식이 나온다.

2절의 운동방정식과 비교했을 때 무엇을 얻고 무엇을 잃는지 짚어 둘 만하다.
블록선도에서는 시스템을 \emph{연결}하는 일이 사소하다. 곱하거나 더하거나
\eqref{eq:blockalg}를 쓰면 된다. 대신 물리가 가려진다. $P(s)$라고 적힌 블록은
그 안에 어떤 질량과 스프링과 댐퍼가 있는지 더는 말해 주지 않는다. 미분방정식은
그 반대다. 모든 항에 물리적 의미가 있지만, 두 시스템을 연결하려면 식을 다시
써야 한다. 둘을 나란히 쓰는 이유가 이것이다.

\begin{figure}[htb]
\centering
\begin{tikzpicture}[>=Stealth,node distance=9mm]
  % (a) series
  \node[lbl] (la) at (-1.1,0) {(a)};
  \node (ua) at (0,0) {$U$};
  \node[draw,minimum width=10mm,minimum height=7mm] (pa) at (1.5,0) {$P$};
  \node[draw,minimum width=10mm,minimum height=7mm] (qa) at (3.2,0) {$Q$};
  \node (ya) at (4.7,0) {$Y$};
  \draw[->] (ua) -- (pa); \draw[->] (pa) -- (qa); \draw[->] (qa) -- (ya);
  % (b) parallel
  \node[lbl] at (6.0,0) {(b)};
  \node (ub) at (6.8,0) {$U$};
  \coordinate (sb) at (7.4,0);
  \node[draw,minimum width=10mm,minimum height=7mm] (pb) at (8.6,0.55) {$P$};
  \node[draw,minimum width=10mm,minimum height=7mm] (qb) at (8.6,-0.55) {$Q$};
  \node[sumc] (jb) at (9.9,0) {};
  \node (yb) at (10.9,0) {$Y$};
  \draw (ub) -- (sb); \fill (sb) circle (1.2pt);
  \draw[->] (sb) |- (pb); \draw[->] (sb) |- (qb);
  \draw[->] (pb) -| node[pos=0.8,right,lbl] {$+$} (jb);
  \draw[->] (qb) -| node[pos=0.8,right,lbl] {$+$} (jb);
  \draw[->] (jb) -- (yb);
  % (c) feedback
  \node[lbl] at (1.2,-2.2) {(c)};
  \node (rc) at (2.0,-2.2) {$R$};
  \node[sumc] (sc) at (3.1,-2.2) {};
  \node[draw,minimum width=10mm,minimum height=7mm] (gc) at (4.9,-2.2) {$G$};
  \coordinate (tc) at (6.5,-2.2);
  \node (yc) at (7.4,-2.2) {$Y$};
  \node[draw,minimum width=10mm,minimum height=7mm] (hc) at (4.9,-3.3) {$H$};
  \draw[->] (rc) -- node[pos=0.75,above,lbl] {$+$} (sc);
  \draw[->] (sc) -- node[above,sig] {$E$} (gc);
  \draw[->] (gc) -- (yc);
  \fill (tc) circle (1.2pt);
  \draw[->] (tc) |- (hc);
  \draw[->] (hc) -| node[pos=0.85,left,lbl] {$-$} (sc);
\end{tikzpicture}
\caption{(a) 직렬: $Y=QPU$. (b) 병렬: $Y=(P+Q)U$. (c) 음의 되먹임:
$Y=\frac{G}{1+GH}R$.}
\label{fig:blocks}
\end{figure}

표준적인 제어 루프에서 앞먹임 경로는 제어기 다음에 플랜트가 오는 $G=P(s)C(s)$이고,
출력은 그대로 되먹여진다($H=1$). 기준입력 $r$은 $y$가 되었으면 하는 값이고,
제어기는 오차 $e=r-y$에 작용한다.
\begin{equation}\label{eq:closedloop}
   \frac{Y(s)}{R(s)} = \frac{P(s)C(s)}{1+P(s)C(s)},
   \qquad
   \frac{E(s)}{R(s)} = \frac{1}{1+P(s)C(s)} .
\end{equation}

\begin{definition}[루프 전달함수]\label{def:loop}
루프를 한 바퀴 돌며 만나는 블록을 모두 곱한 것
\[
   L(s) = P(s)\,C(s)
\]
(일반적으로는 $G(s)H(s)$)를 \emph{루프 전달함수}(loop transfer function)라 한다.
\emph{루프 이득}(loop gain) 또는 \emph{개루프} 전달함수라고도 부른다. 방정식
\begin{equation}\label{eq:chareq}
   1+L(s)=0
\end{equation}
을 폐루프의 \emph{특성방정식}(characteristic equation)이라 한다.
\end{definition}

\eqref{eq:closedloop}의 폐루프 전달함수는 모두 분모가 $1+L$로 같다. $L=N_L/D_L$을
다항식의 비로 쓰면
\[
   \frac{L}{1+L} = \frac{N_L}{D_L+N_L}
\]
이므로 폐루프 극점은 $D_L+N_L$의 근, 즉 \eqref{eq:chareq}의 근이다.
\emph{되먹임이 극점을 옮긴다}는 말의 정확한 뜻이 이것이다. 플랜트 혼자서는
$D_L$의 근을 극점으로 갖는데, 루프를 닫으면 그것이 $D_L+N_L$의 근으로 바뀌고,
그 차이를 조절하는 손잡이가 제어기 $C$다.

\begin{example}[붙들린 물체, 루프로 보기]\label{ex:massloop}
1.3절로 돌아가 모든 것을 목표 높이에서 재자. $e=h-y$. 힘에서 높이로 가는
플랜트는 $P(s)=1/(ms^{2})$이고, 되먹임 법칙 \eqref{eq:pd}는 중력을 보상하는
상수를 빼면 제어기 $C(s)=k+cs$다. 그러면
\[
   L(s) = \frac{k+cs}{ms^{2}},
   \qquad
   1+L(s) = \frac{ms^{2}+cs+k}{ms^{2}}
\]
이다. 개루프는 원점에 이중극점을 갖는다. 떠내려가는 자유 질량이다. 폐루프의
극점은 $ms^{2}+cs+k$의 근이고, $c,k>0$이면 좌반평면에 있다. 이득이 자유 질량에는
없던 \emph{유효 감쇠} $c$와 \emph{유효 강성} $k$가 된 것이고, ``되먹임이
안정화했다''는 말의 내용은 이것이 전부다.
\end{example}

\begin{question}
플랜트 $P$와 제안된 제어기 $C$가 주어졌을 때 폐루프는 안정한가? 한 가지 답은
$D_L+N_L$의 근을 계산하는 것이다. \ref{sec:nyquist}절은 두 번째 답을 준다. 이
답에는 $L(j\omega)$의 그래프만 있으면 되고 --- 그것은 측정할 수 있다 ---, 루프가
안정한지 \emph{여부}뿐 아니라 안정성을 잃기까지 \emph{얼마나 여유가 있는지}도
알려 준다.
\end{question}

%% ================================================================ 제4부
\partdivider{제4부 --- 주파수 응답과 안정도}

%%%%%%%%%%%%%%%%%%%%%%%%%%%%%%%%%%%%%%%%%%%%%%%%%%%%%%%%%%%%%%%%%%%%%%%%%%%%%%
\section{주파수 응답과 Bode 선도}\label{sec:freq}

\subsection{안정한 플랜트}
\begin{definition}[안정한 플랜트]\label{def:stable}
모든 극점의 실수부가 음수일 때 그 플랜트를 \emph{안정하다}(stable)고 한다.
\end{definition}

\begin{remark}[이 절이 안정성에서 시작하는 이유]
아래에서 하는 일은 전부 $P$를 허수축 위, 즉 $s=j\omega$에서 값매김하는 것이다.
다시 말해 라플라스 변환을 푸리에 변환으로 바꿔 쓰는 것인데, 그 대입은 라플라스
적분이 거기서 실제로 수렴할 때만 정당하다. 안정한 플랜트라면 수렴하고,
$P(j\omega)$는 임펄스 응답의 푸리에 변환이 된다. 사인파를 넣어 실제로 측정할 수
있는 양이라는 뜻이다. 반대로 극점이 우반평면에 있으면 허수축은 수렴영역 밖이고,
사인파에 대한 응답은 한없이 커지며, $P(j\omega)$라는 기호는 무엇의 주파수응답도
아니다. 안정성은 이 절 전체의 전제이지 곁가지가 아니다.
\end{remark}

\subsection{사인파를 넣으면 사인파가 나온다}
안정한 플랜트에 $\sin\omega t$를 넣고 과도응답이 잦아들기를 기다리면, 남는 것은
\emph{같은 주파수}의 사인파다.
\begin{equation}\label{eq:freqresp}
   \sin\omega t \;\longmapsto\; A\sin(\omega t+\theta),
   \qquad A = |P(j\omega)|,\qquad \theta = \angle P(j\omega).
\end{equation}
바뀌는 것은 숫자 둘뿐이고, 둘 다 복소수 $P(j\omega)$ 하나에서 읽어낸다. 크기는
진폭을 곱하고 편각은 위상을 옮긴다. 이 둘을 $\omega$의 함수로 본 것이
\emph{주파수응답}이다. 전달함수를 손에 쥘 만한 실질적인 이유가 여기에 있다.
미분방정식이 주파수마다 이득 하나와 위상 하나로 바뀐다.

\begin{proof}[\eqref{eq:freqresp}가 성립하는 이유]
입력의 변환은 $\omega/(s^{2}+\omega^{2})$이므로 \ref{sec:outputpoles}절에 의해
\[
   Y(s) = P(s)\,\frac{\omega}{s^{2}+\omega^{2}}
        = \frac{as+b}{s^{2}+\omega^{2}} + \bigl(P\text{의 극점에서 오는 항들}\bigr)
\]
이다. 플랜트가 안정하므로 두 번째 무리는 사라지고, 남는 것은
$y_{ss}(t) = a\cos\omega t+(b/\omega)\sin\omega t$다. $a$와 $b$를 구하려면 양변에
$s-j\omega$를 곱하고 $s=j\omega$를 넣는다. 플랜트 항들은 사라지고
\[
   P(j\omega)\,\frac{\omega}{2j\omega} = \frac{a\,j\omega+b}{2j\omega}
   \quad\Longrightarrow\quad
   P(j\omega) = \frac{b}{\omega} + j a
\]
가 된다. 따라서 $b/\omega=|P(j\omega)|\cos\angle P(j\omega)$,
$a=|P(j\omega)|\sin\angle P(j\omega)$이고, 삼각함수의 덧셈정리가 $y_{ss}$를
$|P(j\omega)|\sin\bigl(\omega t+\angle P(j\omega)\bigr)$로 합쳐 준다.
\end{proof}

\begin{example}[두 주파수를 한꺼번에]
$u(t)=2\sin t+3\cos(2t+1)$을 $P(s)=1/(s+1)$에 넣자. 선형성에 의해 각 항을 따로
다루면 되고, 코사인도 사인과 같은 규칙을 따른다. $\omega=1$에서
$P(j)=1/(1+j)$의 크기는 $1/\sqrt2$, 위상은 $-\pi/4$이고, $\omega=2$에서
$P(2j)=1/(1+2j)$의 크기는 $1/\sqrt5$, 위상은 $-\arctan2\approx-63.4^\circ$다.
따라서
\[
   y_{ss}(t) = \sqrt2\,\sin\!\Bigl(t-\frac\pi4\Bigr)
             + \frac{3}{\sqrt5}\cos\bigl(2t+1-\arctan 2\bigr)
\]
이다. 높은 주파수일수록 더 많이 줄어들고 더 많이 늦어진다. $1/(s+1)$은 저역통과
필터이고, 다음 소절의 Bode 선도가 그것을 한눈에 보여 준다.
\end{example}

\subsection{Bode 선도}
주파수응답은 가로축을 $\log_{10}\omega$로 놓고 두 장으로 그린다. 하나는 크기를
데시벨로 나타낸
\[
   20\log_{10}|P(j\omega)| ,
\]
다른 하나는 위상 $\angle P(j\omega)$다. 이 두 장을 합쳐 \emph{Bode 선도}라 한다.

두 로그 모두 제값을 한다. 전달함수는 인수들의 \emph{곱}인데, 크기에 로그를 씌우면
그 곱이 \emph{합}으로 바뀐다. 그래서 인수 하나하나가 따로 기여하고, 전체는 조각을
더하기만 하면 얻어진다. 그다음 가로축을 $\log\omega$로 두면 그 조각들이 대부분의
구간에서 직선이 된다.

\subsection{인수로부터 크기를 쌓아 올리기}
예를 들어
\begin{equation}\label{eq:bodeex}
   P(s) = \frac{K(s+z_1)(s+z_2)}
               {(s+p_1)(s+p_2)^{2}\bigl(s^{2}+2\zeta\omega_n s+\omega_n^{2}\bigr)}
\end{equation}
를 보자. 크기에 $20\log_{10}$을 씌우면 인수마다 항 하나가 된다.
\begin{equation}\label{eq:bodesum}
\begin{aligned}
   20\log|P| = 20\bigl(&\log|s+z_1| + \log|s+z_2| - \log|s+p_1| \\
                       &- 2\log|s+p_2| - \log|s^{2}+2\zeta\omega_n s+\omega_n^{2}|\bigr)
\end{aligned}
\end{equation}
이며 $s=j\omega$에서 값매김한다. 분자 인수는 $+$로, 분모 인수는 $-$로 들어가고,
거듭제곱된 인수는 그 지수만큼 들어간다. $\log|s+p_2|$ 앞의 $2$가 그것이다.

1차 인수 하나는 꺾임 하나에 대응한다(그림~\ref{fig:bode}). $s+z_1$을
$s=j\omega$에서 보면
\[
   20\log|j\omega+z_1| = 10\log_{10}\bigl(\omega^{2}+z_1^{2}\bigr)
   \;\approx\;
   \begin{cases}
      20\log z_1, & \omega \ll z_1,\\[2pt]
      20\log \omega, & \omega \gg z_1
   \end{cases}
\]
이므로 점근선은 수평선 하나와 기울기 $+20$\,dB/decade 인 직선 하나이고, 둘은
\emph{꺾임 주파수}(corner frequency) $\omega=z_1$에서 만난다. 정작 그 지점의 참값은
$20\log(z_1\sqrt2) = 20\log z_1 + 3.01$\,dB이다. 즉 곡선은 꺾임점보다 $3$\,dB
\emph{위}로 지난다. 분모 인수는 같은 그림을 뒤집은 것으로 기울기가
$-20$\,dB/decade이고 $3$\,dB \emph{아래}로 지나며, 거듭제곱된 $(s+p_2)^{2}$는 그
둘을 두 배로 해서 $-40$\,dB/decade가 된다.

\begin{figure}[htb]
\centering
\begin{tikzpicture}[>=Stealth,x=1.75cm,y=0.085cm]
  % ---- axes ----
  \draw[->] (-1.62,0) -- (1.80,0) node[below right,lbl] {$\log_{10}\omega$};
  \draw[->] (-1.5,-6) -- (-1.5,35);
  \node[rotate=90,anchor=south,lbl] at (-1.78,16) {$20\log_{10}|j\omega+z_1|$};
  % ---- corner ----
  \draw[dashed,gray] (0,-6) -- (0,33);
  \node[below,lbl] at (0,-1.5) {$\log_{10}z_1$};
  % ---- true curve ----
  \draw[gray!65,thick] plot[domain=-1.5:1.5,samples=90,smooth,variable=\u]
        (\u,{10*log10(1+pow(10,2*\u))});
  % ---- asymptotes ----
  \draw[thick] (-1.5,0) -- (0,0);
  \draw[thick] (0,0) -- (1.5,30);
  % ---- 3 dB at the corner ----
  \node[circle,fill,inner sep=1.1pt] at (0,3.01) {};
  \node[left,lbl] at (-0.05,7) {$3$\,dB};
  \node[above left,lbl] at (1.28,29) {$+20$\,dB/decade};
\end{tikzpicture}
\caption{1차 인수 하나. 두 점근선은 꺾임 주파수 $\omega=z_1$에서 만나고, 참값
(회색)은 꺾임점보다 $3$\,dB 위를 지난다. Bode 크기 선도는 인수마다 이런 그림을
하나씩 더해서 만든다.}
\label{fig:bode}
\end{figure}

\subsection{위상은 각을 더해서}\label{sec:phase}
위상은 더 쉽다. 같은 인수분해가 이미 복소수의 곱이기 때문이다. $s=j\omega$에서
각 인수를 극형식으로 쓰면
\[
   P(j\omega)
   = \frac{\prod_i M_i e^{\,j\theta_i}}{\prod_j M'_j e^{\,j\theta'_j}} ,
   \qquad\text{따라서}\qquad
   \angle P(j\omega) = \sum_i \theta_i - \sum_j \theta'_j
\]
이다. 인수마다 각을 따로 재서 더하고 빼기만 하면 된다. 위상 선도도 크기 선도와
마찬가지로 인수 하나씩 쌓아 만든다. 1차 인수의 각은
\[
   \angle(j\omega+z_1) = \arctan\frac{\omega}{z_1}
\]
로, 낮은 주파수의 $0^\circ$에서 높은 주파수의 $90^\circ$까지 올라가며 꺾임점에서
정확히 $45^\circ$를 지난다. 변화는 거의 전부 꺾임점 양쪽 한 decade 안에서
일어나므로, 흔히 쓰는 직선 근사는 $z_1/10$까지 $0^\circ$를 유지하고, $10z_1$까지
$\log\omega$에 대해 직선으로 $90^\circ$까지 올라간 뒤, 그대로 머문다. 분모
인수는 이것의 부호를 뒤집은 만큼 기여한다.

\begin{remark}[기하학적으로 읽기]\label{rem:vectors}
$j\omega+z_1 = j\omega-(-z_1)$은 영점 $-z_1$에서 허수축 위의 점 $j\omega$로 가는
벡터다. $\omega$가 허수축을 따라 올라가는 동안 이 벡터의 길이가 그 인수의
크기이고, 각이 그 인수의 위상이다. $P(j\omega)$의 인수는 모두 극점이나 영점에서
$j\omega$로 그린 이런 벡터다. $P$의 크기는 영점 벡터 길이의 곱을 극점 벡터
길이의 곱으로 나눈 것이고, 위상은 영점 벡터 각의 합에서 극점 벡터 각의 합을 뺀
것이다. $j\omega$를 닫힌 곡선을 따라 도는 점으로 바꾼 같은 그림이
\ref{sec:nyquist}절의 열쇠다.
\end{remark}

%% ------------------------------------------------------------------ 제6강
\lecturebreak{제6강}{2026년 9월 21일}

\subsection{2차 시스템과 감쇠비}\label{sec:second}
\eqref{eq:bodeex}의 2차 인수는 따로 다룰 만하다. 점근선이 Bode 선도를 잘
설명하지 \emph{못하는} 유일한 인수이고, 모든 감쇠 진동자의 전달함수이기도 하다.
표준 2차 시스템은
\begin{equation}\label{eq:second}
   G(s) = \frac{\omega_n^{2}}{s^{2}+2\zeta\omega_n s+\omega_n^{2}},
   \qquad \omega_n>0,\quad \zeta\ge 0
\end{equation}
로, $G(0)=1$이 되도록 정규화한 것이다. 두 파라미터에는 이름이 있다. $\omega_n$은
\emph{고유진동수}(natural frequency), $\zeta$는 \emph{감쇠비}(damping ratio)다.
둘 다 물리가 정해 준다. 댐퍼가 달린 스프링--질량계 $m\ddot q+b\dot q+kq=ku$를
$m$으로 나누면 정확히 \eqref{eq:second}이고
\begin{equation}\label{eq:zetawn}
   \omega_n = \sqrt{k/m},\qquad
   \zeta = \frac{b}{2\sqrt{km}}
\end{equation}
이다. 즉 $\omega_n$은 감쇠가 없었다면 시스템이 진동했을 주파수이고, $\zeta$는
실제 감쇠를 진동을 막 억누르는 양 $2\sqrt{km}$에 견준 것이다. 감쇠항을
$2\zeta s$가 아니라 $2\zeta\omega_n s$로 적기 때문에 $\zeta$는 차원 없는 수이고,
아래에서 $\zeta$에 대해 하는 말은 시스템의 시간 척도와 무관하게 성립한다.

\smallskip\noindent\textbf{극점.} $s^{2}+2\zeta\omega_n s+\omega_n^{2}$의 근은
\[
   s = -\zeta\omega_n \pm \omega_n\sqrt{\zeta^{2}-1}
\]
이고, $\zeta$의 값에 따라 시스템은 네 영역으로 나뉜다.
\begin{center}\small
\renewcommand{\arraystretch}{1.3}
\begin{tabular}{@{}lll@{}}
\toprule
$\zeta>1$   & 과감쇠(overdamped)        & 서로 다른 음의 실수 극점 두 개\\
$\zeta=1$   & 임계감쇠(critically damped) & $-\omega_n$의 이중극점\\
$0<\zeta<1$ & 부족감쇠(underdamped)       & $-\zeta\omega_n\pm j\omega_d$, \quad $\omega_d=\omega_n\sqrt{1-\zeta^{2}}$\\
$\zeta=0$   & 비감쇠(undamped)          & 허수축 위의 $\pm j\omega_n$\\
\bottomrule
\end{tabular}
\end{center}
흥미로운 것은 부족감쇠인 경우다. 이때 극점은 원 $|s|=\omega_n$ 위에 있고, 원점에서
극점으로 그은 선분이 음의 실축과 이루는 각 $\theta$는
\begin{equation}\label{eq:zetaangle}
   \cos\theta = \zeta
\end{equation}
로 정해진다(그림~\ref{fig:secondpoles}). 원점으로부터의 거리가 주파수이고, 각이
감쇠다. $\zeta$가 일정한 선은 원점에서 나가는 반직선이다.

\begin{figure}[htb]
\centering
\begin{tikzpicture}[>=Stealth,scale=1.0,
   polex/.pic={\draw[thick] (-0.12,-0.12)--(0.12,0.12) (-0.12,0.12)--(0.12,-0.12);}]
  \draw[->] (-2.9,0) -- (1.3,0) node[right] {$\operatorname{Re}s$};
  \draw[->] (0,-2.4) -- (0,2.5) node[above] {$\operatorname{Im}s$};
  \draw[dashed,gray] (90:2) arc (90:270:2);
  \draw[thin] (0,0) -- (-1,1.732);
  \draw[thin] (0,0) -- (-1,-1.732);
  \pic at (-1,1.732) {polex};
  \pic at (-1,-1.732) {polex};
  \draw[dotted] (-1,1.732) -- (-1,0);
  \draw[dotted] (-1,1.732) -- (0,1.732);
  \node[below,lbl] at (-1,-0.02) {$-\zeta\omega_n$};
  \node[right,lbl] at (0.02,1.732) {$j\omega_d$};
  \draw (-0.6,0) arc (180:120:0.6);
  \node[lbl] at (-0.78,0.33) {$\theta$};
  \node[lbl,above left] at (-0.48,0.86) {$\omega_n$};
  \node[lbl,anchor=west] at (0.35,-1.25) {$\cos\theta=\zeta$};
\end{tikzpicture}
\caption{\eqref{eq:second}의 부족감쇠 극점 쌍을 $\zeta=0.5$로 그린 것. 극점은
반지름 $\omega_n$인 원 위에 있고, 극점이 음의 실축과 이루는 각 $\theta$는 감쇠비
하나만으로 정해진다.}
\label{fig:secondpoles}
\end{figure}

\smallskip\noindent\textbf{시간 응답.} $0<\zeta<1$일 때 단위계단입력
$Y(s)=G(s)/s$에 \ref{sec:pfe}절의 부분분수를 적용하면
\begin{equation}\label{eq:secondstep}
   y(t) = 1 - \frac{e^{-\zeta\omega_n t}}{\sqrt{1-\zeta^{2}}}\,
          \sin\bigl(\omega_d t+\theta\bigr),
   \qquad \theta=\arccos\zeta
\end{equation}
를 얻는다. 이 식의 모든 조각을 극점에서 읽을 수 있다. 실수부 $-\zeta\omega_n$이
감쇠하는 포락선을, 허수부 $\omega_d$가 출렁임의 주파수를 정한다
(그림~\ref{fig:secondstep}). $\zeta$가 작으면 빠르지만 심하게 출렁이는 응답이,
$\zeta$가 $1$에 가까우면 거의 넘치지 않지만 굼뜬 응답이 나온다. 곡선을 요약하는
수는 둘, 최대 오버슈트(overshoot)와 정착시간(settling time)이다.
\begin{equation}\label{eq:overshoot}
   M_p = e^{-\pi\zeta/\sqrt{1-\zeta^{2}}},
   \qquad
   t_s \approx \frac{4}{\zeta\omega_n}\quad\text{($2\%$ 이내 기준)} .
\end{equation}
오버슈트는 $\zeta$에만, 즉 극점의 \emph{각}에만 의존하고, 정착시간은
$\zeta\omega_n$에만, 즉 극점이 \emph{허수축에서 떨어진 거리}에만 의존한다.

\begin{figure}[htb]
\centering
\begin{tikzpicture}[>=Stealth,x=0.46cm,y=2.1cm]
  \draw[->] (0,0) -- (14.7,0) node[right,lbl] {$\omega_n t$};
  \draw[->] (0,0) -- (0,1.95) node[above,lbl] {$y(t)$};
  \draw[dotted,gray] (0,1) -- (14,1);
  \foreach \tk in {0,2,...,14}{\draw (\tk,0) -- (\tk,-0.04) node[below,lbl] {\tk};}
  \foreach \yk in {0.5,1,1.5}{\draw (0,\yk) -- (-0.18,\yk) node[left,lbl] {\yk};}
  \draw[gray] plot[domain=0:14,samples=150,smooth]
    (\x,{1-exp(-0.1*\x)/sqrt(0.99)*sin(deg(sqrt(0.99)*\x)+acos(0.1))});
  \draw[dashed] plot[domain=0:14,samples=120,smooth]
    (\x,{1-exp(-0.3*\x)/sqrt(0.91)*sin(deg(sqrt(0.91)*\x)+acos(0.3))});
  \draw[thick] plot[domain=0:14,samples=100,smooth]
    (\x,{1-exp(-0.7*\x)/sqrt(0.51)*sin(deg(sqrt(0.51)*\x)+acos(0.7))});
  \draw[thick,dotted] plot[domain=0:14,samples=80,smooth] (\x,{1-(1+\x)*exp(-\x)});
  % legend
  \begin{scope}[shift={(9.2,0.14)}]
    \draw[gray] (0,0.36) -- (1.4,0.36); \node[right,lbl] at (1.5,0.36) {$\zeta=0.1$};
    \draw[dashed] (0,0.24) -- (1.4,0.24); \node[right,lbl] at (1.5,0.24) {$\zeta=0.3$};
    \draw[thick] (0,0.12) -- (1.4,0.12); \node[right,lbl] at (1.5,0.12) {$\zeta=0.7$};
    \draw[thick,dotted] (0,0) -- (1.4,0); \node[right,lbl] at (1.5,0) {$\zeta=1$};
  \end{scope}
\end{tikzpicture}
\caption{\eqref{eq:second}의 단위계단응답을 $\omega_n t$에 대해 그린 것. 감쇠가
적을수록 첫 봉우리가 높고 출렁임이 늦게 잦아든다. $\zeta=1$은 전혀 넘치지 않는다.}
\label{fig:secondstep}
\end{figure}

\begin{remark}[학기 목표, 다시 보기]
질문~\ref{q:goal}은 수레가 $5$초 안에 정착하기를 요구한다. \eqref{eq:overshoot}에
따르면 응답이 얼마나 빨리 정착하는지와 얼마나 넘치는지는 둘 다 \emph{우세
극점이 어디에 있는지}에 관한 말이다. $5$\,s 안에 정착하려면
$\zeta\omega_n \gtrsim 4/5$여야 하고, 오버슈트는 각 $\theta$가 정한다. 개루프
two-cart 시스템에게는 극점을 어디에 두라고 말할 수가 없다. 폐루프는 할 수 있다.
\ref{sec:blocks}절에서 되먹임이 한 일이 그것이다.
\end{remark}

\smallskip\noindent\textbf{주파수응답.} 허수축 위에서
\begin{equation}\label{eq:secondfreq}
   G(j\omega) = \frac{\omega_n^{2}}{(\omega_n^{2}-\omega^{2})+j\,2\zeta\omega_n\omega}
\end{equation}
이다. 정규화한 주파수 $u=\omega/\omega_n$을 쓰면
\[
   |G(j\omega)| = \frac{1}{\sqrt{(1-u^{2})^{2}+(2\zeta u)^{2}}},
   \qquad
   \angle G(j\omega) = -\operatorname{atan2}\bigl(2\zeta u,\;1-u^{2}\bigr).
\]
점근선은 2차식에서 기대하는 그대로다. $u\ll1$이면 크기는 $1$(0\,dB), 위상은
$0^\circ$이고, $u\gg1$이면 크기는 $1/u^{2}$로 기울기 $-40$\,dB/decade인 직선,
위상은 $-180^\circ$다. 둘은 $\omega=\omega_n$에서 만난다. 그러나 꺾임점
\emph{에서} 무슨 일이 일어나는지는 $\zeta$가 정한다. \eqref{eq:secondfreq}에
$u=1$을 넣으면
\[
   G(j\omega_n) = \frac{1}{j\,2\zeta},
   \qquad
   |G(j\omega_n)| = \frac{1}{2\zeta},
   \qquad
   \angle G(j\omega_n) = -90^\circ
\]
이다. 꺾임점의 위상은 모든 $\zeta$에 대해 $-90^\circ$이지만, 크기는
$1/(2\zeta)$로 감쇠가 작으면 커진다. $\zeta=0.1$이면 곡선이 점근선보다 $14$\,dB
위에 온다(그림~\ref{fig:secondbode}). 1차 인수는 점근선에서 $3$\,dB 넘게 벗어나는
일이 없지만, 2차 인수는 높은 \emph{공진 봉우리}(resonant peak)를 가질 수 있다.
$|G|^{2}$를 $u$로 미분한 것을 0으로 두면, $\zeta<1/\sqrt2$일 때 봉우리가
\begin{equation}\label{eq:resonance}
   \omega_r = \omega_n\sqrt{1-2\zeta^{2}},
   \qquad
   M_r = |G(j\omega_r)| = \frac{1}{2\zeta\sqrt{1-\zeta^{2}}}
\end{equation}
에 있고, $\zeta\ge1/\sqrt2\approx0.707$이면 봉우리가 아예 없음을 알 수 있다. 한편
위상은 $\zeta$가 작아질수록 $0^\circ$에서 $-180^\circ$로 점점 더 급하게 떨어지고,
$\zeta\to0$인 극한에서는 $\omega_n$에서 $180^\circ$만큼 한 번에 뛴다.

\begin{figure}[htb]
\centering
\begin{tikzpicture}[>=Stealth,x=2.7cm]
  % ---------------- magnitude ----------------
  \begin{scope}[y=0.058cm]
    \draw[->] (-1,0) -- (1.13,0) node[right,lbl] {$\log_{10}(\omega/\omega_n)$};
    \draw[->] (-1,-42) -- (-1,20) node[above,lbl] {dB};
    \foreach \d in {-40,-20,0}{\draw (-1,\d) -- (-1.03,\d) node[left,lbl] {\d};}
    \draw[dotted,gray] (0,-42) -- (0,18);
    \draw[gray!70,densely dashed] (-1,0) -- (0,0) -- (1,-40);
    \draw[gray] plot[domain=-1:1,samples=160,smooth]
      (\x,{-10*log10(pow(1-pow(10,2*\x),2)+pow(0.2*pow(10,\x),2))});
    \draw[dashed] plot[domain=-1:1,samples=120,smooth]
      (\x,{-10*log10(pow(1-pow(10,2*\x),2)+pow(0.6*pow(10,\x),2))});
    \draw[thick] plot[domain=-1:1,samples=100,smooth]
      (\x,{-10*log10(pow(1-pow(10,2*\x),2)+pow(1.4*pow(10,\x),2))});
    \draw[thick,dotted] plot[domain=-1:1,samples=100,smooth]
      (\x,{-10*log10(pow(1-pow(10,2*\x),2)+pow(2.0*pow(10,\x),2))});
    \begin{scope}[shift={(-0.93,-33)}]
      \draw[gray] (0,9) -- (0.2,9); \node[right,lbl] at (0.22,9) {$\zeta=0.1$};
      \draw[dashed] (0,4.5) -- (0.2,4.5); \node[right,lbl] at (0.22,4.5) {$\zeta=0.3$};
      \draw[thick] (0.78,9) -- (0.98,9); \node[right,lbl] at (1.0,9) {$\zeta=0.7$};
      \draw[thick,dotted] (0.78,4.5) -- (0.98,4.5); \node[right,lbl] at (1.0,4.5) {$\zeta=1$};
    \end{scope}
  \end{scope}
  % ---------------- phase ----------------
  \begin{scope}[yshift=-4.3cm,y=0.017cm]
    \draw[->] (-1,0) -- (1.13,0) node[right,lbl] {$\log_{10}(\omega/\omega_n)$};
    \draw[->] (-1,-190) -- (-1,12) node[above,lbl] {위상};
    \foreach \d in {0,-90,-180}{\draw (-1,\d) -- (-1.03,\d) node[left,lbl] {$\d^\circ$};}
    \draw[dotted,gray] (0,-190) -- (0,5);
    \draw[dotted,gray] (-1,-90) -- (1,-90);
    \draw[gray] plot[domain=-1:1,samples=160,smooth]
      (\x,{-atan2(0.2*pow(10,\x),1-pow(10,2*\x))});
    \draw[dashed] plot[domain=-1:1,samples=120,smooth]
      (\x,{-atan2(0.6*pow(10,\x),1-pow(10,2*\x))});
    \draw[thick] plot[domain=-1:1,samples=100,smooth]
      (\x,{-atan2(1.4*pow(10,\x),1-pow(10,2*\x))});
    \draw[thick,dotted] plot[domain=-1:1,samples=100,smooth]
      (\x,{-atan2(2.0*pow(10,\x),1-pow(10,2*\x))});
  \end{scope}
\end{tikzpicture}
\caption{2차 시스템 \eqref{eq:second}의 Bode 크기 선도(위)와 위상 선도(아래).
위 그림의 회색 파선이 점근선이다. 모든 곡선이 $\omega=\omega_n$에서
$-90^\circ$를 지나고, 거기서의 봉우리 $20\log_{10}(1/2\zeta)$는 $\zeta\to0$일 때
한없이 커진다.}
\label{fig:secondbode}
\end{figure}

\begin{remark}[한 사실, 세 그림]
약한 감쇠는 세 번 나타난다. 허수축에 가까운 극점으로, 계단응답의 출렁임으로,
Bode 선도의 공진 봉우리로. 이것들은 같은 사실을 세 방향에서 본 것이다.
$-\zeta\omega_n+j\omega_d$의 극점은 허수축에서 $\zeta\omega_n$만큼 떨어져 있으므로,
$j\omega$가 그 옆을 지날 때 보충~\ref{rem:vectors}의 극점 벡터가 아주 짧아진다.
그리고 분모의 짧은 벡터는 큰 크기를 뜻한다. 실제로 $\omega=\omega_d$에서 그
벡터의 길이는 $\zeta\omega_n$이고 다른 극점 벡터의 길이는 약 $2\omega_d$이므로
$|G|\approx\omega_n^{2}/(2\zeta\omega_n\omega_d)=M_r$이 된다.
\end{remark}

\begin{example}[감쇠 없는 공진]
$\zeta\to0$의 극한을 끝까지 밀어 보자. $\zeta=0$이면
$G(s)=\omega_n^{2}/(s^{2}+\omega_n^{2})$의 극점이 허수축 위에 있으므로 안정하지
않고, \eqref{eq:freqresp}는 적용되지 않는다. 정확히 고유진동수로,
$u=\sin\omega_n t$로 구동해 보자.
\[
   Y(s) = \frac{\omega_n^{2}}{s^{2}+\omega_n^{2}}\cdot\frac{\omega_n}{s^{2}+\omega_n^{2}}
        = \frac{\omega_n^{3}}{(s^{2}+\omega_n^{2})^{2}},
   \qquad
   y(t) = \tfrac12\bigl(\sin\omega_n t-\omega_n t\cos\omega_n t\bigr).
\]
입력의 극점이 플랜트의 극점과 겹쳐 $\pm j\omega_n$의 극점이 이중극점이 되고,
\ref{sec:pfe}절에 따르면 이중극점은 $t$를 곱해 준다. 진동이 선형으로, 한없이
커진다. 이것이 가장 순수한 형태의 공진이고, $\zeta\to0$인 Bode 곡선의 무한한
봉우리는 그 주파수 영역의 그림자다.
\end{example}

\begin{remark}[우세 극점]\label{rem:dominant}
실제 플랜트에는 극점이 많은데 왜 2차 시스템을 이렇게 자세히 보는가? 안정한
시스템에서 극점마다 항 $e^{p_it}$가 하나씩 생기는데, 멀리 왼쪽에 있는 극점의
항은 허수축 가까운 극점의 항보다 훨씬 먼저 사라진다. 그래서 허수축에 가장 가까운
극점, 즉 \emph{우세 극점}(dominant poles)이 처음 잠깐 이후에 보이는 모든 것을
결정한다. 우세 극점이 나머지와 충분히 떨어진 복소수 쌍이면(그리고 영점과 거의
약분되어 유수가 작아지는 경우가 아니면), $n$차 시스템은 그 두 극점을 가진 2차
시스템과 거의 똑같이 움직인다. 2차가 아닌 플랜트에서도 $\zeta$, $\omega_n$,
오버슈트, 정착시간을 실무 용어로 쓰는 이유가 이것이다.
\end{remark}

\subsection{기본 인수 목록}\label{sec:catalogue}
모든 유리 전달함수는 아래 인수들의 곱이므로, 모든 Bode 선도는 이들 그림의 합이다.
1차 인수는 $0$\,dB에서 출발하는 정규화된 꼴로 적었다.

\begin{center}\small
\renewcommand{\arraystretch}{1.6}
\begin{tabular}{@{}llll@{}}
\toprule
인수 & 크기 & 위상 & 꺾임점에서\\
\midrule
$K>0$ & $20\log_{10}K$로 평평 & $0^\circ$ & ($K<0$이면 $\pm180^\circ$ 추가)\\
$s$ & $+20$\,dB/dec, $\omega=1$에서 $0$\,dB & $+90^\circ$ & ---\\
$1/s$ & $-20$\,dB/dec, $\omega=1$에서 $0$\,dB & $-90^\circ$ & ---\\
$1+s/z$ & $0$, $\omega=z$부터 $+20$\,dB/dec & $0^\circ\to+90^\circ$ & $+3$\,dB, $+45^\circ$\\
$1/(1+s/p)$ & $0$, $\omega=p$부터 $-20$\,dB/dec & $0^\circ\to-90^\circ$ & $-3$\,dB, $-45^\circ$\\
$\dfrac{\omega_n^{2}}{s^{2}+2\zeta\omega_n s+\omega_n^{2}}$ & $0$, $\omega_n$부터 $-40$\,dB/dec & $0^\circ\to-180^\circ$ & $1/(2\zeta)$, $-90^\circ$\\
\bottomrule
\end{tabular}
\end{center}

\begin{example}[완전한 스케치]
$P(s)=\dfrac{10(s+1)}{s(s+10)}=\dfrac{1+s}{s\,(1+s/10)}$을 보자. 적분기 하나,
꺾임 주파수 $1$인 1차 영점, 꺾임 주파수 $10$인 1차 극점의 곱이다. 왼쪽부터
읽으면, $\omega=1$ 아래에서는 적분기만 작동하므로 크기는 $-20$\,dB/decade로
떨어지며 $\omega=1$에서 $0$\,dB를 지나고 위상은 $-90^\circ$다. $\omega=1$에서
영점이 기울기를 상쇄해 크기는 $0$\,dB에서 평평해지고 위상은 $0^\circ$ 쪽으로
올라간다. $\omega=10$에서 극점이 기울기를 다시 $-20$\,dB/decade로, 위상을 다시
$-90^\circ$ 쪽으로 되돌린다. 전체 선도는 꺾임점 $1$과 $10$에서 이어 붙인 직선
세 조각이다.
\end{example}

%%%%%%%%%%%%%%%%%%%%%%%%%%%%%%%%%%%%%%%%%%%%%%%%%%%%%%%%%%%%%%%%%%%%%%%%%%%%%%
\section{Nyquist 선도와 폐루프 안정도}\label{sec:nyquist}

\subsection{질문, 다시}
\ref{sec:blocks}절에 따르면 루프 전달함수가 $L=PC$인 폐루프는 특성방정식
$1+L(s)=0$이 $\operatorname{Re}s\ge0$인 근을 갖지 않을 때, 그리고 그때에만
안정하다.
\[
   F(s) = 1+L(s) = \frac{D_L(s)+N_L(s)}{D_L(s)}
\]
로 두자. $F$의 \emph{영점}은 폐루프 극점이다. 모르는 것이고, 바로 우리가 찾고
싶은 것이다. $F$의 \emph{극점}은 $L$의 극점, 곧 개루프 극점이다. 플랜트를 알고
제어기는 우리가 골랐으니 아는 것이다. Nyquist의 아이디어는 우반평면에 있는 모르는
영점을 계산하지 않고 세는 것이다. $s$가 그 반평면의 경계를 따라 도는 동안 $F$의
위상을 지켜보면 된다.

\subsection{각으로 세기: 편각 원리}\label{sec:argument}
$\Gamma$를 $s$-평면의 단순 닫힌 곡선으로, 한 바퀴 \emph{시계방향}으로 도는
것이라 하고, $F$는 $\Gamma$ 위에 극점도 영점도 없는 유리함수라 하자. $s$가
$\Gamma$를 도는 동안 점 $F(s)$는 $F$-평면에 닫힌 곡선 $F(\Gamma)$를 그린다
(그림~\ref{fig:argument}). $F$를 \eqref{eq:factored}처럼 인수분해하면
\ref{sec:phase}절에 의해
\[
   \angle F(s) = \angle k + \sum_i \angle(s-z_i) - \sum_j \angle(s-p_j)
\]
이다. 이제 인수 하나만 따라가 보자. $z_i$가 $\Gamma$ \emph{안}에 있으면, $z_i$에서
움직이는 점 $s$로 가는 벡터 $s-z_i$는 $s$가 $\Gamma$를 도는 동안 한 바퀴 온전히
돈다. 시계방향이므로 각의 변화는 $-2\pi$다. $z_i$가 \emph{밖}에 있으면 벡터는
이리저리 흔들리지만 결국 제자리로 돌아온다. 알짜 변화는 $0$이다. 극점도 마찬가지인데,
$\angle F$에서 부호만 반대다. 인수들을 모두 더하면 다음을 얻는다.

\begin{theorem}[편각 원리]\label{thm:argument}
$\Gamma$ 안에 있는 $F$의 영점과 극점의 개수를 중복도까지 세어 각각 $Z$와 $P$라
하자. $s$가 $\Gamma$를 시계방향으로 한 바퀴 돌면 $\angle F(s)$는 모두
$-2\pi(Z-P)$만큼 변한다. 즉 곡선 $F(\Gamma)$는 원점을 시계방향으로
\[
   N = Z - P
\]
번 감싼다(포위한다). 반시계방향으로 한 번 감싸는 것은 $-1$로 센다.
\end{theorem}

예를 들어 $\Gamma$ 안에 두 점 $a$, $b$를, 밖에 세 번째 점 $c$를 두자. 그러면
$F=s-a$는 원점을 시계방향으로 한 번 감고, $F=(s-a)(s-b)$는 $\theta_a$와
$\theta_b$가 각각 한 바퀴씩 돌므로 두 번 감으며, $F=(s-a)(s-b)(s-c)$도
\emph{여전히} 두 번 감는다. $\theta_c$는 $c$에서 $\Gamma$에 그은 두 접선 사이의
쐐기 안에서 흔들리기만 하기 때문이다. 극점은 방향을 뒤집는다. $\angle F=-\theta_a$
이므로 $F=1/(s-a)$는 \emph{반시계방향}으로 한 번 감는다.

이 각의 셈이 \emph{곧} 증명이다. \ref{sec:phase}절의 위상 규칙을, 허수축을 따라
올라가는 점 대신 고리를 따라 도는 점에 적용한 것이다. (복소해석에서는 같은 셈이
반시계방향 $\Gamma$에 대한
$\frac{1}{2\pi j}\oint_\Gamma F'(s)/F(s)\,\d s = Z-P$로 나타난다.
$F'/F=\sum_i 1/(s-z_i)-\sum_j 1/(s-p_j)$이고, $\Gamma$ 안에 든 항은 각각 유수
$+1$ 또는 $-1$을 보탠다.)

\begin{figure}[htb]
\centering
\begin{tikzpicture}[>=Stealth,
   polex/.pic={\draw[thick] (-0.11,-0.11)--(0.11,0.11) (-0.11,0.11)--(0.11,-0.11);}]
  % ---- s-plane ----
  \begin{scope}
    \node[lbl] at (-1.9,1.55) {$s$-평면};
    \draw[thick] (0,0) ellipse (1.6 and 1.1);
    \draw[->,thick] (100:1.6 and 1.1) arc[start angle=100,end angle=80,x radius=1.6,y radius=1.1];
    \draw[->,thick] (-80:1.6 and 1.1) arc[start angle=-80,end angle=-100,x radius=1.6,y radius=1.1];
    \node[lbl] at (1.35,-1.05) {$\Gamma$};
    \draw[thick] (-0.35,0.15) circle (0.1);
    \node[lbl,below] at (-0.35,0.03) {$z_1$};
    \pic at (2.35,0.75) {polex};
    \node[lbl,right] at (2.45,0.75) {$p_1$};
    \coordinate (s) at (40:1.6 and 1.1);
    \fill (s) circle (1.3pt); \node[lbl,above right] at (s) {$s$};
    \draw[->] (-0.35,0.15) -- ($(s)+(-0.05,-0.03)$);
    \draw[->,gray] (2.35,0.75) -- ($(s)+(0.06,0.0)$);
  \end{scope}
  % ---- F-plane ----
  \begin{scope}[xshift=6.4cm]
    \node[lbl] at (-1.6,1.55) {$F$-평면};
    \draw[->] (-1.9,0) -- (1.9,0);
    \draw[->] (0,-1.5) -- (0,1.6);
    \draw[thick] plot[domain=0:360,samples=90,smooth,variable=\a]
      ({0.25+1.25*cos(\a)+0.25*cos(2*\a)},{-0.95*sin(\a)+0.12*sin(3*\a)});
    \draw[->,thick] (1.5,-0.1) -- (1.48,-0.3);
    \fill (0,0) circle (1.3pt);
    \node[lbl] at (1.45,1.05) {$F(\Gamma)$};
  \end{scope}
\end{tikzpicture}
\caption{편각 원리. $s$가 $\Gamma$를 시계방향으로 한 바퀴 돌면, 안에 든 영점에서
나가는 벡터 $s-z_1$은 시계방향으로 한 바퀴 돌고, 밖에 있는 극점에서 나가는 벡터
$s-p_1$은 앞뒤로 흔들리기만 한다. 안에 영점 하나, 극점은 없으므로 상
$F(\Gamma)$는 원점을 시계방향으로 한 번 감싼다. $N=Z-P=1$.}
\label{fig:argument}
\end{figure}

\subsection{Nyquist 경로와 Nyquist 선도}
우반평면의 근을 세려면 $\Gamma$를 우반평면의 경계로 잡으면 된다. \emph{Nyquist
경로}(Nyquist contour) $\Gamma_N$은 허수축을 따라 $-jR$에서 $+jR$까지 올라간 뒤,
우반평면을 지나는 반원 $|s|=R$을 따라 시계방향으로 돌아오며, $R\to\infty$로
보낸다(그림~\ref{fig:contour}). 그 안쪽은 열린 우반평면 전체다. $L$이 허수축
\emph{위}에 극점을 가지면 --- 이를테면 적분기 $1/s$ --- $\Gamma_N$은 그 극점을
지나갈 수 없다. 대신 우반평면 쪽으로 부푼 반지름 $\varepsilon\to0$의 작은 반원으로
각 극점을 비켜 가고, 그러면 그 극점들은 바깥에 남는다.

\begin{figure}[htb]
\centering
\begin{tikzpicture}[>=Stealth,scale=0.95,
   polex/.pic={\draw[thick] (-0.11,-0.11)--(0.11,0.11) (-0.11,0.11)--(0.11,-0.11);}]
  \fill[black!5] (0,-2.2) -- (0,-0.32) arc (-90:90:0.32) -- (0,2.2) arc (90:-90:2.2) -- cycle;
  \draw[->] (-1.4,0) -- (3.0,0) node[right] {$\operatorname{Re}s$};
  \draw[->] (0,-2.7) -- (0,2.75) node[above] {$\operatorname{Im}s$};
  \draw[thick] (0,-2.2) -- (0,-0.32);
  \draw[->,thick] (0,-2.2) -- (0,-1.15);
  \draw[thick] (0,-0.32) arc (-90:90:0.32);
  \draw[thick] (0,0.32) -- (0,2.2);
  \draw[->,thick] (0,0.32) -- (0,1.35);
  \draw[thick] (0,2.2) arc (90:-90:2.2);
  \draw[->,thick] (45:2.2) arc (45:35:2.2);
  \pic at (0,0) {polex};
  \node[lbl,anchor=west] at (1.6,1.95) {$R\to\infty$};
  \node[lbl,anchor=west] at (0.42,0.36) {$\varepsilon\to0$};
  \node[lbl] at (1.05,-0.95) {우반평면};
\end{tikzpicture}
\caption{시계방향으로 도는 Nyquist 경로 $\Gamma_N$. 허수축을 따라 올라간 뒤 반지름이
무한대인 반원을 따라 돌아온다. 그 안쪽(음영)이 열린 우반평면이다. 허수축 위의
극점 --- 여기서는 원점 --- 은 우반평면 쪽으로 작게 비켜 가서 피한다.}
\label{fig:contour}
\end{figure}

$L$에 의한 $\Gamma_N$의 상을 $L$의 \emph{Nyquist 선도}(Nyquist plot)라 한다. 두
가지 사실 덕분에 그리기가 쉽다.
\begin{itemize}[leftmargin=*,itemsep=2pt]
\item $L$이 strictly proper이면 $|s|\to\infty$일 때 $L(s)\to0$이므로, 큰 반원 전체가
한 점 $0$으로 간다. 그러면 선도는 곡선 $L(j\omega)$, $-\infty<\omega<\infty$일
뿐이다. 주파수응답을 Bode 선도 두 장 대신 복소평면에 그린 것이다.
\item $L$의 계수가 실수이므로 $L(-j\omega)=\overline{L(j\omega)}$이다. 즉
$\omega<0$인 절반은 $\omega>0$인 절반을 실축에 대해 뒤집은 것이다.
\end{itemize}
실제로는 $\omega\ge0$에 대해 $L(j\omega)$를 그리고 --- 원점까지의 거리가
$|L(j\omega)|$, 각이 $\angle L(j\omega)$이므로 Bode 선도에서 바로 읽어낸다 ---
그것을 실축에 대해 뒤집은 뒤, 작은 우회로가 보태는 부분을 더한다.

\subsection{Nyquist 안정도 판별법}
정리~\ref{thm:argument}를 $\Gamma_N$ 위의 $F=1+L$에 적용하자. $\Gamma_N$ 안에 있는
$F$의 영점은 우반평면의 폐루프 극점이고, 안에 있는 $F$의 극점은 거기 있는 개루프
극점이다. 그리고 $F(\Gamma_N)$은 $L$의 Nyquist 선도를 오른쪽으로 한 칸 옮긴
것이므로, $F$가 원점을 감싸는 것은 $L$이 점 $-1$을 감싸는 것과 같다. 판별법은
이것이 전부다.

\begin{keybox}
\textbf{Nyquist 안정도 판별법.} 열린 우반평면에 있는 $L(s)$의 극점 수를 $P$,
$L$의 Nyquist 선도가 점 $-1$을 시계방향으로 감싸는 횟수를 $N$이라 하자
(반시계방향으로 감싸는 것은 음수로 센다). 그러면 우반평면에 있는 폐루프 극점의
수는
\[
   Z = N + P
\]
이다. 폐루프는 $Z=0$, 즉 $N=-P$이고 선도가 $-1$을 지나지 않을 때, 그리고
그때에만 안정하다(선도가 $-1$을 지나면 폐루프는 허수축 위에 극점을 갖는다).
특히 $L$ 자체가 안정하면($P=0$) 폐루프는 Nyquist 선도가 $-1$을 감싸지 않을 때,
그리고 그때에만 안정하다.

\smallskip\noindent
반대 방향으로 세면 같은 말이 이렇게 된다. \emph{$-1$을 반시계방향으로 감싸는
횟수는 $P-Z$와 같다.}
\end{keybox}

\begin{remark}[왜 이런 수고를 할 만한가]
$D_L+N_L$의 근을 그냥 계산하는 것과 비교해 이 판별법에는 두 가지가 다르다.
첫째, $L(j\omega)$만 있으면 되고 그것은 \emph{측정}할 수 있다. 개루프를
사인파로 구동해 이득과 위상을 기록하면 되므로, 플랜트 모델이 전혀 필요 없다.
둘째, 루프가 안정한지뿐 아니라 불안정에 얼마나 \emph{가까운지}도 보여 준다.
$-1$ 근처를 지나는 선도는 이득이나 위상이 조금만 바뀌어도 --- 1절에서 반드시
예상해야 한다고 했던 바로 그런 모델링 오차로 --- $-1$ 건너편으로 밀려날 수 있다.
\end{remark}

\subsection{예제}
\begin{example}[안정한 개루프]\label{ex:nyq1}
$K>0$인 $L(s)=\dfrac{K}{(s+1)(s+2)(s+3)}$. 개루프 극점이 모두 좌반평면에 있으므로
$P=0$이다. $s=j\omega$에서 분모는
\[
   (6-6\omega^{2}) + j\,(11\omega-\omega^{3})
\]
이다. 선도는 양의 실축 위의 $L(0)=K/6$에서 출발해, 위상이 $0^\circ$에서
$-270^\circ$ 쪽으로 떨어지는 동안 제4사분면과 제3사분면을 지나 시계방향으로
움직인다. 음의 실축은 분모의 허수부가 사라지는 $\omega^{2}=11$에서 가로지르고,
거기서 $L=K/(6-66)=-K/60$이다(그림~\ref{fig:nyqex} 왼쪽, $K=20$). $K<60$이면
교점이 $-1$의 오른쪽에 있어 감싸지 않는다. $N=0$, $Z=0$, 안정. $K>60$이면 선도가
$-1$을 시계방향으로 두 번 감싼다. $N=2$이므로 폐루프 극점 두 개가 우반평면에
있다. 경계인 $K=60$에서는 선도가 $\omega=\sqrt{11}$에서 $-1$을 지나며, 실제로
\[
   D_L+N_L = s^{3}+6s^{2}+11s+66 = (s+6)(s^{2}+11)
\]
의 폐루프 극점이 정확히 $\pm j\sqrt{11}$, 즉 선도가 $-1$을 만나는 주파수에 있다.
\end{example}

\begin{example}[루프 안의 적분기]\label{ex:nyq2}
$K>0$인 $L(s)=\dfrac{K}{s(s+1)}$. 원점의 극점이 허수축 위에 있으므로 $\Gamma_N$은
$\phi$가 $-90^\circ$에서 $+90^\circ$까지 가는 $s=\varepsilon e^{j\phi}$로 그것을
비켜 간다. 거기서 $L(s)\approx K/s = (K/\varepsilon)e^{-j\phi}$이므로, 반지름이
엄청나게 큰 $K/\varepsilon$인 반원을 $+90^\circ$에서 $-90^\circ$까지, 즉 평면의
오른쪽 절반을 지나 시계방향으로 그린다. 허수축 위에서는
\[
   L(j\omega) = -\frac{K}{1+\omega^{2}} - j\,\frac{K}{\omega(1+\omega^{2})}
\]
이므로, $\omega>0$에 대해 선도는 점근선 $\operatorname{Re}L=-K$를 따라 아래
멀리서 올라와 제3사분면에 머물다가 $\omega\to\infty$일 때 위상 $-180^\circ$로
원점에 들어간다. 허수부가 결코 0이 되지 않으므로 음의 실축을 가로지르지 않고,
$-1$은 결코 감싸지지 않는다. 비켜 간 극점은 세지 않으므로 $P=0$이고 $N=0$이다.
따라서 $Z=0$이고 모든 $K>0$에 대해 루프가 안정하다. 직접 확인하면
$D_L+N_L=s^{2}+s+K$이고, 모든 $K>0$에 대해 그 근의 실수부는 음수다.
\end{example}

우회로는 언제나 이 예제처럼 움직인다. 허수축 위의 단순 극점 $p$ 근처에서
$L(s)\approx c/(s-p)$이고, 우회로가 $s-p$를 $+180^\circ$ 돌리는 동안(아래를
가리키다가 오른쪽을 거쳐 위를 가리키도록) 상은 무한대 반지름에서 $-180^\circ$
돈다. 시계방향 반 바퀴다. 허수축 위의 이중극점은 시계방향 한 바퀴,
$-360^\circ$를 주며, 그것만으로도 $-1$을 감쌀 수 있다.

\begin{example}[허수축 위의 극점 쌍]\label{ex:nyq4}
$L(s)=\dfrac{1}{s^{2}+4}=\dfrac{1}{(s-2j)(s+2j)}$. 두 극점 모두 허수축 위에 있어
비켜 가므로 $P=0$이다. 그 밖에서는 $L(j\omega)=1/(4-\omega^{2})$가 실수다.
$|\omega|\to2$일 때 양의 실축을 따라 $\tfrac14$에서 $+\infty$까지 가고,
$|\omega|>2$에서는 음의 실축을 따라 $-\infty$에서 $0$으로 돌아온다. 각 우회로는
무한대에서 시계방향 반원을 보태는데 --- $+2j$의 것은 아래 반평면을, $-2j$의 것은
위 반평면을 지난다 ---, 둘이 합쳐 선도를 무한히 큰 원으로 닫는다. 음의 실축
전체가 덮이므로 선도는 $4-\omega^{2}=-1$, 즉 $\omega=\sqrt5$에서 $-1$을
\emph{지난다}. 판별법에 따르면 폐루프는 허수축 위에 극점을 갖고, 실제로 $1+L=0$은
$s^{2}+5=0$, $s=\pm j\sqrt5$를 뜻한다. $\sqrt5$\,rad/s의 진동이 줄지도 커지지도
않고 계속된다.
\end{example}

\begin{example}[불안정한 개루프]\label{ex:nyq3}
$K>0$인 $L(s)=\dfrac{K}{s-1}$. 이제 개루프가 $s=1$에 극점을 가지므로 $P=1$이다.
허수축 위에서
\[
   L(j\omega) = \frac{K}{j\omega-1}
   = -\frac{K}{1+\omega^{2}} - j\,\frac{K\omega}{1+\omega^{2}}
\]
이고, 이것은 $-K/2$를 중심으로 하는 반지름 $K/2$인 원을 그린다. 원점
($\omega=-\infty$)에서 출발해 $-\tfrac K2+j\tfrac K2$, 이어서 $-K$($\omega=0$),
이어서 $-\tfrac K2-j\tfrac K2$를 거쳐 원점으로 돌아오는 반시계방향이다
(그림~\ref{fig:nyqex} 오른쪽). $K>1$이면 원이 $-1$을 품고 반시계방향으로 한 번
감싼다. $N=-1$이므로 $Z=N+P=0$이고 루프는 안정하다. $K<1$이면 $-1$을 비켜 간다.
$N=0$, $Z=1$, 불안정. 직접 보면 $D_L+N_L=s-1+K$이고, 정확히 $K>1$일 때 안정하다.
불안정한 플랜트를 안정화하려면 감싸기가 \emph{필요하다}. ``선도가 $-1$을 감싸면
안 된다''는 안정한 개루프에만 해당하는 규칙이다.
\end{example}

\begin{remark}[Bode 선도에서 판별법 읽기]\label{rem:bodenyq}
$L$이 안정하고($P=0$) 그 위상이 한 주파수 $\omega_{pc}$에서 단 한 번
$-180^\circ$를 지나 떨어지면, $\omega>0$인 Nyquist 선도는 점 $-|L(j\omega_{pc})|$
에서 음의 실축을 단 한 번 가로지른다. 선도는 그 점이 $-1$의 왼쪽에 있을 때 정확히
$-1$을 감싼다. 그래서 판별법은 Bode 선도에서 확인할 수 있는 명제가 된다.
\[
   \text{폐루프 안정}
   \quad\Longleftrightarrow\quad
   |L(j\omega_{pc})| < 1
   \quad\text{단, } \angle L(j\omega_{pc}) = -180^\circ .
\]
위상이 루프 신호를 완전히 뒤집어 놓을 즈음에는 이득이 이미 1보다 작아야 한다.
예제~\ref{ex:nyq1}에서는 위상이 $\omega_{pc}=\sqrt{11}$에서 $-180^\circ$에
이르고, 거기서 $|L|=K/60$이다. 다시 같은 조건 $K<60$이다. 거기서 이득이 너무
크면 이득을 낮추거나(지상 보상기의 아이디어), 위상을 다시 끌어올려 $-180^\circ$에
더 늦게 도달하게 할 수 있다(진상 보상기).
\end{remark}

\begin{figure}[htb]
\centering
\begin{tikzpicture}[>=Stealth,
   declare function={Lm(\w)=20/sqrt(\w*\w+1)/sqrt(\w*\w+4)/sqrt(\w*\w+9);
                     La(\w)=-(atan(\w)+atan(\w/2)+atan(\w/3));}]
  % ---------------- Example: K/((s+1)(s+2)(s+3)), K=20 ----------------
  \begin{scope}[scale=0.8]
    \draw[->] (-1.7,0) -- (4.2,0) node[right,lbl] {$\operatorname{Re}L$};
    \draw[->] (0,-2.8) -- (0,2.8) node[above,lbl] {$\operatorname{Im}L$};
    \draw[thick] plot[domain=-2:1.6,samples=120,smooth]
      ({Lm(pow(10,\x))*cos(La(pow(10,\x)))},{Lm(pow(10,\x))*sin(La(pow(10,\x)))});
    \draw[dashed] plot[domain=-2:1.6,samples=120,smooth]
      ({Lm(pow(10,\x))*cos(La(pow(10,\x)))},{-Lm(pow(10,\x))*sin(La(pow(10,\x)))});
    \draw[->,thick] ({Lm(0.30)*cos(La(0.30))},{Lm(0.30)*sin(La(0.30))})
                 -- ({Lm(0.34)*cos(La(0.34))},{Lm(0.34)*sin(La(0.34))});
    \fill (-1,0) circle (2.2pt);
    \node[lbl,above left] at (-1,0.05) {$-1$};
    \draw (-0.333,0.1) -- (-0.333,-0.1);
    \node[lbl,below left] at (-0.30,-0.08) {$-\tfrac13$};
    \node[lbl,above right] at (3.333,0) {$\tfrac{10}{3}$};
    \node[lbl,right] at (0.12,-2.05) {$\omega=1$};
    \node[lbl] at (3.35,-1.45) {$\omega>0$};
    \node[lbl] at (3.35,1.45) {$\omega<0$};
  \end{scope}
  % ---------------- Example: K/(s-1), K=1.5 ----------------
  \begin{scope}[xshift=8.6cm,scale=1.45]
    \draw[->] (-2.0,0) -- (0.7,0) node[right,lbl] {$\operatorname{Re}L$};
    \draw[->] (0,-1.1) -- (0,1.15) node[above,lbl] {$\operatorname{Im}L$};
    \draw[thick] (-0.75,0) circle (0.75);
    \draw[->,thick] (-0.75,0.75) arc (90:110:0.75);
    \draw[->,thick] (-0.75,-0.75) arc (270:290:0.75);
    \fill (-1,0) circle (1.5pt);
    \node[lbl,above] at (-1.02,0.04) {$-1$};
    \node[lbl,below left] at (-1.5,0) {$-K$};
    \node[lbl,left] at (-1.52,0.18) {$\omega=0$};
    \node[lbl,above right] at (0,0) {$\omega\to\pm\infty$};
  \end{scope}
\end{tikzpicture}
\caption{왼쪽: 예제~\ref{ex:nyq1}의 $K=20$일 때 Nyquist 선도(실선 $\omega>0$,
파선 $\omega<0$). 음의 실축을 $-1$의 오른쪽인 $-K/60=-\tfrac13$에서 가로지르므로
$-1$을 감싸지 않고 폐루프는 안정하다. 오른쪽: 예제~\ref{ex:nyq3}의 $K=1.5$.
개루프는 불안정하고($P=1$) 선도가 $-1$을 반시계방향으로 한 번 돈다($N=-1$).
폐루프를 안정하게 만드는 것이 바로 이것이다.}
\label{fig:nyqex}
\end{figure}

%% ================================================================ 제5부
\partdivider{제5부 --- 상태공간을 위한 선형대수}

\lecturebreak{제7$\sim$10강}{2026년 9월 23일$\sim$10월 7일}

\section{벡터공간, 계수, 선형방정식}\label{sec:vspace}

상태는 $\R^{n}$의 한 점이고, 입력은 그 점을 특정한 방향으로 밀며, 출력은 그 좌표들의
특정한 조합을 알려 준다. 그러니 이 과목의 나머지 부분이 $\dot x=Ax+Bu$를 두고 하는
모든 일은 벡터, 부분공간, 행렬에 관한 진술이다. 이 대단원은 앞으로 쓸 선형대수를
필요해지는 순서대로 모아 둔다.

\subsection{열로 읽는 행렬}
행렬 $A\in\R^{m\times n}$은 $m$개의 행과 $n$개의 열을 가진다. $m=n$이면
\emph{정방}(square) 행렬, $m>n$이면 \emph{세로로 긴}(tall) 행렬, $m<n$이면
\emph{가로로 긴}(fat) 행렬이라 한다. 전치 $A\tp\in\R^{n\times m}$은
$(ABC)\tp=C\tp B\tp A\tp$를 만족하고, 복소 행렬에서는 \emph{켤레전치}(Hermitian
transpose) $A^{H}=\overline{A}\tp$를 쓴다. $A\in\R^{m\times n}$과
$B\in\R^{n\times k}$의 곱 $C=AB$의 성분은 $c_{ij}=\sum_{l}a_{il}b_{lj}$이지만, 이
곱을 읽는 다른 세 방법이 훨씬 쓸모 있다. $A$의 열을 $a_1,\dots,a_n$, $B$의 행을
$b_1\tp,\dots,b_n\tp$이라 쓰자.
\begin{enumerate}[label=(\roman*),itemsep=2pt]
\item $Ax = x_1a_1+\dots+x_na_n$: 행렬과 벡터의 곱은 \emph{열들의 선형결합}(linear
combination)이다.
\item $y\tp B = y_1b_1\tp+\dots+y_nb_n\tp$: 행벡터와 행렬의 곱은 행들의
선형결합이다.
\item $AB = \sum_{i=1}^{n} a_i\,b_i\tp$: 곱은 \emph{외적}(outer product)\emph{의
합}이고, 각 $a_ib_i\tp$는 계수가 1인 행렬이다.
\end{enumerate}
내내 머릿속에 두어야 할 것은 읽기 (i)이다. $Ax=b$를 풀 수 있느냐고 묻는 것은 $b$가
$A$의 열들의 조합이냐고 묻는 것과 같다.

\subsection{벡터공간, 생성, 기저, 차원}
\begin{definition}[벡터공간]
벡터들의 집합 $S$가 덧셈과 스칼라곱에 대해 닫혀 있으면, 즉 $x,y\in S$이고
$a,b\in\R$일 때 $ax+by\in S$이면, $S$를 \emph{벡터공간}(vector space)($\R^{n}$의
\emph{부분공간}(subspace))이라 한다.
\end{definition}
$a=b=0$으로 두면 모든 벡터공간이 원점을 품는다는 것을 알 수 있다. \emph{원점을
지나는} 직선이나 평면은 부분공간이고, 원점을 비껴가는 것은 부분공간이 아니다.

$a_1,\dots,a_k$의 \emph{선형결합}은 $\sum_i c_ia_i$이고, $\sum_ic_ia_i=0$이면
반드시 모든 $c_i=0$일 때 이 벡터들을 \emph{선형독립}(linearly independent)이라
한다. 어느 하나도 나머지로 만들어 낼 수 없다는 것, 즉 모두가 정말로 서로 다른
방향을 가리킨다는 뜻이다. 선형결합 전체의 집합이 \emph{생성}(span)
\[
   \operatorname{span}\{a_1,\dots,a_k\} = \Bigl\{\,\sum_{i=1}^{k}c_ia_i \;:\; c_i\in\R\Bigr\}
\]
이고, 이것은 벡터공간이다. 두 조합을 더하거나 한 조합에 스칼라를 곱해도 또 하나의
조합이 되기 때문이다. 벡터공간의 \emph{기저}(basis)는 그 공간을 생성하는 선형독립인
벡터들의 집합이다. 기저는 유일하지 않지만 어느 기저든 벡터의 개수는 같고, 그 개수가
\emph{차원}(dimension)이다. $\R^{3}$의 평면 $x_3=0$에서 벡터 하나로는 평면을 생성할
수 없고, 그 안의 평행하지 않은 아무 두 벡터나 기저를 이루므로 이 평면의 차원은 $2$다.
벡터들이 서로 수직이고 길이가 1인 기저를 \emph{정규직교}(orthonormal) 기저라 한다.

부분공간을 기술하는 방법은 두 가지이고, 둘 다 쓰게 된다.
\begin{itemize}[leftmargin=*,itemsep=2pt]
\item \textbf{쌓아 올려서}: 생성으로 기술한다. 평면 $x_3=0$은
$\operatorname{span}\{e_1,e_2\}$이다. 독립인 생성 벡터 하나마다 차원이 하나씩
늘어난다($0\to2$).
\item \textbf{깎아 내려서}: 방정식으로 기술한다. 같은 평면은
$\{x\in\R^{3}:[0\;0\;1]\,x=0\}$이다. 독립인 제약 하나마다 차원이 하나씩 줄어든다
($3\to2$).
\end{itemize}
$x_2x_3$-평면에서 원점을 지나 $45^\circ$로 놓인 직선은
$\operatorname{span}\{(0,1,1)\}$이고, 똑같이
$\{x\in\R^{3}: [1\;0\;0]\,x=0,\ [0\;1\;-1]\,x=0\}$이기도 하다. 독립인 제약 둘이
$3-2=1$차원을 남긴다.

\subsection{치역, 영공간, 방정식 \texorpdfstring{$Ax=b$}{Ax=b}}
\label{sec:rangenull}
모든 행렬 $A\in\R^{m\times n}$에는 두 부분공간이 딸려 있고, 두 기술 방식이 하나씩이다.
\begin{equation}\label{eq:rangenull}
   \mathcal R(A) = \{Ax : x\in\R^{n}\}\subseteq\R^{m},
   \qquad
   \mathcal N(A) = \{x\in\R^{n} : Ax=0\}\subseteq\R^{n}.
\end{equation}
읽기 (i)에 따라 \emph{치역}(range) $\mathcal R(A)$는 열들의 생성이다. 이것을
\emph{열공간}(column space)이라고도 부르며, 쌓아 올리는 기술이다. \emph{영공간}(null
space) $\mathcal N(A)$는 깎아 내리는 기술이다. $A$의 \emph{계수}(rank)는 선형독립인
열의 개수이고, 이는 언제나 선형독립인 행의 개수와 같으며, $\mathcal R(A)$의 차원이다.

두 부분공간은 선형방정식 $Ax=b$에 대해 묻는 두 질문에 답한다.
\begin{itemize}[leftmargin=*,itemsep=2pt]
\item \textbf{존재.} 해가 존재할 필요충분조건은 $b\in\mathcal R(A)$, 곧 $b$가
열들의 조합이라는 것이다.
\item \textbf{유일성.} $Ax=b$이고 $Ax^{*}=b$이면 $A(x-x^{*})=0$이므로 두 해는
$\mathcal N(A)$의 원소만큼 차이 난다. 해가 유일할 필요충분조건은
$\mathcal N(A)=\{0\}$이고, 읽기 (i)로 보면 이는 $x_1a_1+\dots+x_na_n=0$이 $x=0$일
때만 성립한다는 것, 곧 열들이 선형독립이라는 것이다.
\end{itemize}
열들이 독립인 행렬, 즉 $\rank A=n$인 행렬은 \emph{열 완전 계수}(full column
rank)라 하고, 그때 $\mathcal N(A)=\{0\}$이다. 존재와 유일성은 답이 다른 별개의
질문이고, 둘을 떼어 생각하는 편이 이롭다.

\begin{remark}[차원 세기]
위에서 깎아 내리며 센 것은 일반적인 사실이다. $A\in\R^{m\times n}$에 대해
\[
   \dim\mathcal R(A)+\dim\mathcal N(A) = n
\]
이고, 이를 \emph{계수--퇴화차수 관계}(rank--nullity)라 한다. $n$개의 입력 방향 가운데
$\rank A$개는 독립인 출력에 닿고, 나머지 $n-\rank A$개는 $0$으로 보내진다.
\end{remark}

\begin{example}[두 수레의 평형점]\label{ex:cartnull}
\eqref{eq:ABCD}의 두 수레 행렬 $A$에서, 힘을 가하지 않을 때 제자리에 머무는 상태는
$Ax=0$의 해다. $1$행과 $3$행에서 $x_2=x_4=0$, 곧 두 수레가 모두 정지해 있다. 그러면
$2$행은 $\frac{k}{m_1}(x_3-x_1)=0$이 되어 $x_1=x_3$이고, $4$행도 같은 말을 한다. 따라서
\[
   \mathcal N(A) = \operatorname{span}\{(1,0,1,0)\tp\},
   \qquad \rank A = 3 .
\]
두 수레가 같은 거리만큼 나아가 있고 스프링이 늘어나지 않은 위치는 어디든 평형점이다.
수레를 되돌려 당기는 바닥 스프링이 없기 때문이다. 특히 질문~\ref{q:goal}의 목표,
두 수레가 $1$\,m 앞에 정지해 있는 $x=(1,0,1,0)\tp$는 제어하지 않은 시스템의
평형점이다. 일단 거기에 가면 수레는 $u=0$으로 그 자리에 머문다. 댐퍼 $b$를 남겨
두어도 달라지는 것은 없다. $b$가 들어간 항은 모두 속도가 곱해져 있기 때문이다.
\end{example}

\subsection{역행렬과 행렬식}
정방 행렬 $A\in\R^{n\times n}$은 $AA^{-1}=A^{-1}A=I$인 $A^{-1}$이 있으면
\emph{가역}(invertible)이라 한다. $A^{-1}=[x_1\;\cdots\;x_n]$으로 쓰면 방정식
$AA^{-1}=I$는 $n$개의 방정식 $Ax_i=e_i$이므로, 역행렬은 모든 $e_i$, 따라서
$\R^{n}$ 전체가 $\mathcal R(A)$에 들어 있을 때 정확히 존재한다. 정방 행렬에서는
그것이 정확히 $\mathcal N(A)=\{0\}$일 때 일어난다. 그러므로 다음은 모두 동치다.
$A$가 가역이다; $\mathcal R(A)=\R^{n}$; $\mathcal N(A)=\{0\}$; 열들이 독립이다;
$\rank A=n$; $\det A\neq0$; $0$이 $A$의 고윳값이 아니다.

$2\times2$ 행렬에서는
\[
   \begin{bmatrix}a&b\\c&d\end{bmatrix}^{-1}
   = \frac{1}{ad-bc}\begin{bmatrix}d&-b\\-c&a\end{bmatrix}
   = \frac{\operatorname{adj}A}{\det A}
\]
이고, 공식 $A^{-1}=\operatorname{adj}A/\det A$는 모든 차원에서 성립한다. 여기서
$\operatorname{adj}A$는 \emph{수반행렬}(adjugate), $\det A$는
\emph{행렬식}(determinant)이며, 이것이 \ref{sec:ss2tf}절의
$(sI-A)^{-1}=\operatorname{adj}(sI-A)/\det(sI-A)$ 뒤에 있는 공식이다. 행렬식은
$\det A\tp=\det A$와 $\det(ABC)=\det A\,\det B\,\det C$를 만족한다.

\begin{proposition}[슈어 보수]\label{prop:schur}
$A$와 $D$가 정방인 $M=\begin{bmatrix}A&B\\C&D\end{bmatrix}$에 대해, $A$가
가역이면
\[
   \det M = \det A\,\det\bigl(D-CA^{-1}B\bigr)
\]
이고, $D$가 가역이면 $\det M=\det D\,\det\bigl(A-BD^{-1}C\bigr)$이다.
\end{proposition}
\begin{proof}
블록 소거: $\begin{bmatrix}A&B\\C&D\end{bmatrix}
=\begin{bmatrix}I&0\\CA^{-1}&I\end{bmatrix}
\begin{bmatrix}A&B\\0&D-CA^{-1}B\end{bmatrix}$. 첫 인수의 행렬식은 $1$이고 둘째
인수는 블록 삼각이다. 다른 경우도 대칭적이다.
\end{proof}
행렬 $D-CA^{-1}B$를 $M$에서 $A$의 \emph{슈어 보수}(Schur complement)라 한다. 블록
$A$를 소거하고 나서 $D$에 남는 것으로, 스칼라에 대한 가우스 소거와 똑같다.

%%%%%%%%%%%%%%%%%%%%%%%%%%%%%%%%%%%%%%%%%%%%%%%%%%%%%%%%%%%%%%%%%%%%%%%%%%%%%%
\section{좌표, 고윳값, 대각화}\label{sec:eigen}

\subsection{상태 좌표 바꾸기}\label{sec:similarity}
시스템의 상태벡터는 하나의 선택이다. 각 수레의 위치와 속도가 한 가지 선택이고, 두
수레 묶음의 위치와 스프링의 늘어난 길이가 또 다른 선택이다. 가역인 아무
$T\in\R^{n\times n}$이나 새 좌표
\[
   z(t) = T\,x(t)
\]
를 정한다. $x=T^{-1}z$를 \eqref{eq:std}에 대입하면
\begin{equation}\label{eq:simtransform}
   \dot z = TAT^{-1}\,z + TB\,u, \qquad y = CT^{-1}\,z + D\,u .
\end{equation}
사중쌍 $(A,B,C,D)$가 $(TAT^{-1},TB,CT^{-1},D)$로 바뀌는데, 이를
\emph{닮음변환}(similarity transformation)이라 한다. 물리는 바뀌지 않았으니 잴 수
있는 것은 무엇도 바뀌어서는 안 되고, 실제로 전달함수는 그대로다.
\[
   CT^{-1}\bigl(sI-TAT^{-1}\bigr)^{-1}TB + D
   = CT^{-1}\bigl(T(sI-A)T^{-1}\bigr)^{-1}TB + D
   = C(sI-A)^{-1}B + D
\]
여기서 $sI=T(sI)T^{-1}$을 썼다. \ref{sec:ss2tf}절의 실현이 전혀 유일하지 않은
이유 하나가 이것이다. 가역인 $T$ 하나하나가 한 실현을 같은 전달함수를 가진 다른
실현으로 바꾼다.

\subsection{고윳값과 고유벡터}
\begin{definition}[고윳값, 고유벡터]
$Av=\lambda v$이면 수 $\lambda$를 $A\in\R^{n\times n}$의 \emph{고윳값}(eigenvalue),
$v\neq0$을 그 \emph{고유벡터}(eigenvector)라 한다.
\end{definition}
고유벡터는 $A$가 돌리지 않고 $\lambda$배로 늘이기만 하는 방향이다. $Av=\lambda v$를
$(\lambda I-A)v=0$으로 고쳐 쓰면, $\lambda$가 고윳값인 것은 정확히
$\mathcal N(\lambda I-A)\neq\{0\}$일 때, 곧 $\lambda I-A$가 가역이 아닐 때다.
\begin{equation}\label{eq:charpoly}
   \det(\lambda I-A) = \lambda^{n}+a_{n-1}\lambda^{n-1}+\dots+a_1\lambda+a_0 = 0 .
\end{equation}
좌변이 \emph{특성다항식}(characteristic polynomial)이고 차수가 $n$이므로, $A$는
중복을 세어 $n$개의 고윳값을 가진다. 계수가 실수이므로 복소 고윳값은 켤레쌍으로
나온다. $\mathcal N(\lambda I-A)$의 영 아닌 벡터는 모두 고유벡터이고, 이 부분공간이
$\lambda$의 \emph{고유공간}(eigenspace)이다.

고윳값은 좌표에 의존하지 않는다. \eqref{eq:simtransform}에 따라
\[
   \det\bigl(\lambda I-TAT^{-1}\bigr) = \det\bigl(T(\lambda I-A)T^{-1}\bigr)
   = \det T\,\det(\lambda I-A)\,\det T^{-1} = \det(\lambda I-A)
\]
이므로 특성다항식 전체가 불변이다. 그리고
$P(s)=C\operatorname{adj}(sI-A)B/\det(sI-A)+D$의 분모는 정확히 $\det(sI-A)$다.
분자와 약분되는 인수가 있는 경우를 빼면 \emph{$A$의 고윳값이 곧 전달함수의
극점이다}. 그림~\ref{fig:routes}의 두 기술이 여기서 만난다.

\subsection{대각화}\label{sec:diag}
$A$가 고윳값 $\lambda_1,\dots,\lambda_n$에 속하는 선형독립인 고유벡터
$v_1,\dots,v_n$을 $n$개 가진다고 하자. $n$개의 방정식 $Av_i=\lambda_iv_i$를 옆으로
나란히 쌓으면
\[
   A\,[v_1\;\cdots\;v_n] = [\lambda_1v_1\;\cdots\;\lambda_nv_n]
   = [v_1\;\cdots\;v_n]
     \begin{bmatrix}\lambda_1&&\\&\ddots&\\&&\lambda_n\end{bmatrix},
\]
즉
\begin{equation}\label{eq:diagonalize}
   AV = V\Lambda,\qquad V^{-1}AV=\Lambda,\qquad A = V\Lambda V^{-1}.
\end{equation}
$V$가 가역인 것은 바로 그 열들이 독립이기 때문이다. \ref{sec:similarity}절의 말로
하면, 좌표 $z=V^{-1}x$가 상태행렬을 대각으로 만든다. 이것이 \emph{대각화}
(diagonalization)다.

독립인 고유벡터가 언제 $n$개 있는가? 두 개수가 이를 정한다. 고윳값의 \emph{대수적
중복도}(algebraic multiplicity)는 그것이 \eqref{eq:charpoly}의 근으로 몇 번
나오는가이고, \emph{기하적 중복도}(geometric multiplicity)는
$\dim\mathcal N(\lambda I-A)$, 곧 그 고윳값이 가진 독립인 고유벡터의 개수다.
언제나 $1\le$ 기하적 $\le$ 대수적이다.
\begin{itemize}[leftmargin=*,itemsep=2pt]
\item \emph{서로 다른} 고윳값에 속하는 고유벡터들은 선형독립이다. 그러므로 $n$개의
고윳값이 모두 다르면 $A$는 대각화 가능하다.
\item 일반적으로, $A$가 대각화 가능할 필요충분조건은 모든 고윳값에서 기하적 중복도가
대수적 중복도와 같은 것이다.
\end{itemize}

\begin{example}[중복도가 같을 때와 다를 때]
두 행렬
\[
   A=\begin{bmatrix}1&0&0&0\\0&2&0&0\\0&0&2&0\\0&0&0&3\end{bmatrix}
   \quad\text{와}\quad
   B=\begin{bmatrix}1&0&0&0\\0&2&1&0\\0&0&2&0\\0&0&0&3\end{bmatrix}
\]
는 모두 고윳값 $1,2,2,3$을 가지며, 두 행렬 모두에서 고윳값 $2$의 대수적 중복도는
$2$다. $A$에서는 $\mathcal N(2I-A)=\operatorname{span}\{e_2,e_3\}$의 차원이 $2$다.
독립인 고유벡터가 둘이고, $A$는 (이미) 대각이다. $B$에서는
\[
   2I-B=\begin{bmatrix}1&0&0&0\\0&0&-1&0\\0&0&0&0\\0&0&0&-1\end{bmatrix},
   \qquad \mathcal N(2I-B)=\operatorname{span}\{e_2\}
\]
의 차원이 $1$이다. 고윳값의 두 사본이 고유벡터 방향 하나를 나눠 쓰므로 $V$를
가역으로 만들 수 없고, $B$는 대각화 가능하지 않다. 차이를 만드는 것은 대각 밖의
$1$ 하나다.
\end{example}

\begin{example}[두 수레, 분리하기]\label{ex:cartdecouple}
질량이 같다고 하자, $m_1=m_2=m$(그리고 \eqref{eq:ABCD}에서처럼 $b=0$). 두 운동방정식
\eqref{eq:eom}을 더하고 빼면
\[
   m(\ddot p_1+\ddot p_2) = u - c(\dot p_1+\dot p_2),\qquad
   m(\ddot p_1-\ddot p_2) = u - 2k(p_1-p_2) - c(\dot p_1-\dot p_2).
\]
그래서 좌표 $z=(p_1+p_2,\;\dot p_1+\dot p_2,\;p_1-p_2,\;\dot p_1-\dot p_2)$, 곧
\[
   T=\begin{bmatrix}1&0&1&0\\0&1&0&1\\1&0&-1&0\\0&1&0&-1\end{bmatrix},
   \qquad
   TAT^{-1}=\begin{bmatrix}0&1&0&0\\0&-\frac cm&0&0\\0&0&0&1\\0&0&-\frac{2k}m&-\frac cm\end{bmatrix},
   \qquad TB=\frac1m\begin{bmatrix}0\\1\\0\\1\end{bmatrix}
\]
인 $z=Tx$는 시스템을 서로 독립인 두 조각으로 가른다. 하나는 마찰만 느끼는 묶음
전체의 운동이고, 다른 하나는 마찰과 두 배의 스프링을 느끼는, 두 수레가 서로에 대해
하는 진동이다. 이제 고윳값을 블록별로 읽어 낼 수 있다. 첫 블록에서 $0$과 $-c/m$,
둘째 블록에서 $\lambda^{2}+\frac cm\lambda+\frac{2k}{m}$의 근이다. 첫 번째 과제에서
쓴 인자($m=1$, $k=4$, $c=0.4$)로는 $0$, $-0.4$, $-0.2\pm2.8213j$이다. 고윳값
$0$은 예제~\ref{ex:cartnull}의 고유벡터 $(1,0,1,0)\tp$에 속한다. 개루프에서 아무것도
붙잡아 주지 않는, 묶음의 자유로운 표류다.
\end{example}

\subsection{대칭행렬}\label{sec:symmetric}
대칭행렬 $A=A\tp$는 무엇보다 다루기 좋은 행렬이고, \ref{sec:svd}절에서 다시 나온다.
\begin{theorem}[스펙트럼 정리]\label{thm:spectral}
$A=A\tp\in\R^{n\times n}$이라 하자. 그러면 $A$의 고윳값은 모두 실수이고, 서로 다른
고윳값에 속하는 고유벡터는 직교하며, 정규직교인 고유벡터 $n$개를 고를 수 있다. 그것들을
$V$의 열로 두면
\[
   V\tp V=I,\qquad V^{-1}=V\tp,\qquad
   A = V\Lambda V\tp = \sum_{i=1}^{n}\lambda_i\,v_iv_i\tp .
\]
\end{theorem}
\begin{proof}[직교성]
$Av_i=\lambda_iv_i$이고 $Av_j=\lambda_jv_j$이면
$\lambda_i\,v_j\tp v_i = v_j\tp Av_i = v_j\tp A\tp v_i = (Av_j)\tp v_i
= \lambda_j\,v_j\tp v_i$이므로 $(\lambda_i-\lambda_j)\,v_j\tp v_i=0$이고,
$\lambda_i\neq\lambda_j$이면 $v_j\tp v_i=0$이다.
\end{proof}
차원이 $k$인 고유공간 안에서는 언제나 정규직교인 벡터 $k$개를 고를 수 있고, 서로 다른
고윳값의 고유공간끼리는 직교하므로, 이 선택들을 모두 모으면 정규직교 집합이 된다.
대칭행렬에서는 $V$의 역행렬이 공짜다. 전치가 곧 역행렬이고, $V\tp AV=\Lambda$는
$V^{-1}AV$보다 훨씬 계산하기 쉽다. 마지막 꼴 $A=\sum\lambda_iv_iv_i\tp$는 $A$를
고유방향마다 하나씩, 계수가 1인 조각들의 합으로 쓴 것이다.

\subsection{조르당 표준형}\label{sec:jordan}
위의 $B$처럼 고유벡터가 모자란 행렬은 대각화할 수 없지만, 모든 정방 행렬은 그다음으로
좋은 꼴로 바꿀 수 있다.
\begin{theorem}[조르당 표준형]
모든 $A\in\R^{n\times n}$은 어떤 가역(복소일 수도 있는) $V$에 대해 $AV=VJ$를
만족한다. 여기서 $J$는 대각에
\[
   J_k(\lambda)=\begin{bmatrix}\lambda&1&&\\&\lambda&\ddots&\\&&\ddots&1\\&&&\lambda\end{bmatrix}
\]
꼴의 \emph{조르당 블록}(Jordan block)이 놓인 블록 대각행렬이다. 이를 \emph{조르당
표준형}(Jordan form)이라 한다. $A$가 대각화 가능한 것은 정확히 모든 블록이 $1\times1$일
때다.
\end{theorem}
$V$의 열이 무엇인지 보려고 $3\times3$ 예를 들자.
\[
   AV=VJ,\qquad J=\begin{bmatrix}\lambda_1&0&0\\0&\lambda_2&1\\0&0&\lambda_2\end{bmatrix},
   \qquad
   [Av_1\;Av_2\;Av_3]=[\lambda_1v_1\;\;\lambda_2v_2\;\;\lambda_2v_3+v_2].
\]
처음 두 열은 보통의 고유벡터다. 셋째 열은
\[
   (A-\lambda_2I)\,v_3 = v_2
\]
를 만족하므로 고유벡터가 아니다. $A-\lambda_2I$가 이것을 소멸시키지 않고 고유벡터
$v_2$로 보낸다. 이런 $v_3$을 \emph{일반화 고유벡터}(generalized eigenvector)라 한다.
일반화 고유벡터는 정확히 어떤 고윳값의 대수적 중복도가 기하적 중복도보다 클 때
존재하며, 모자란 고유벡터가 주었을 방향을 대신 채워 준다.

%%%%%%%%%%%%%%%%%%%%%%%%%%%%%%%%%%%%%%%%%%%%%%%%%%%%%%%%%%%%%%%%%%%%%%%%%%%%%%
\section{행렬지수와 상태 응답}\label{sec:expm}

\subsection{\texorpdfstring{$\dot x=Ax$}{xdot = Ax} 풀기}
스칼라에서 $\dot x=ax$의 해는 $x(t)=e^{at}x_0$이다. 행렬의 지수를 같은 급수로
정의하고 나면 시스템에서도 똑같다.
\begin{definition}[행렬지수]
$A\in\R^{n\times n}$에 대해 \emph{행렬지수}(matrix exponential)는
\[
   e^{At} = I + At + \frac{(At)^{2}}{2!} + \frac{(At)^{3}}{3!} + \cdots .
\]
\end{definition}
이 급수는 모든 $A$와 $t$에서 수렴한다. 항별로 미분하면
\[
   \frac{\d}{\d t}e^{At} = A + A^{2}t + \frac{A^{3}t^{2}}{2!}+\cdots
   = A\Bigl(I+At+\frac{(At)^{2}}{2!}+\cdots\Bigr) = Ae^{At}
\]
이므로 $x(t)=e^{At}x_0$은 $\dot x=Ax$와 $x(0)=x_0$을 만족한다. 이것이 \emph{바로 그}
해다. 성질들은 스칼라의 것을 따르는데, 예외가 하나 있다.
\begin{equation}\label{eq:expmprops}
   e^{A\cdot0}=I,\qquad e^{At_1}e^{At_2}=e^{A(t_1+t_2)},\qquad
   \bigl(e^{At}\bigr)^{-1}=e^{-At},\qquad
   e^{At}e^{Bt}\neq e^{(A+B)t}\ \text{$AB=BA$가 아니면} .
\end{equation}
셋째는 둘째에서 $t_2=-t_1$로 두면 나온다. 특히 $e^{At}$는 언제나 가역이다. 마지막
것이 성립하지 않는 이유는 급수를 곱해 펼치면 $AB$와 $BA$가 서로 다른 순서로 나오기
때문이다. MATLAB에서 행렬지수는 \verb|exp|(성분별로 지수를 취한다)가 아니라
\verb|expm|이다. $x_0$에서 출발한 둘째 수레의 $t=10$에서의 위치는
\verb|[0 0 1 0]*expm(A*10)*x0|이다.

\subsection{상태 응답}\label{sec:stateresponse}
입력이 있으면 $\dot x-Ax=Bu$에 $e^{-At}$를 곱한다. 좌변은 완전미분
$e^{-At}\dot x-e^{-At}Ax=\frac{\d}{\d t}\bigl(e^{-At}x(t)\bigr)$이 되고, $0$에서
$t$까지 적분하면 $e^{-At}x(t)-x_0=\int_0^{t}e^{-A\tau}Bu(\tau)\,\d\tau$이다.
$e^{At}$를 곱하면
\begin{equation}\label{eq:voc}
   \boxed{\;x(t) = e^{At}x_0 + \int_{0}^{t} e^{A(t-\tau)}\,B\,u(\tau)\,\d\tau ,\qquad
   y(t) = Ce^{At}x_0 + C\!\int_{0}^{t} e^{A(t-\tau)}Bu(\tau)\,\d\tau + Du(t).\;}
\end{equation}
이것이 \emph{상수변화 공식}(variation of constants)이다. 첫 항은 초기 상태가 혼자서
하는 일인 \emph{영입력 응답}(zero-input response)이고, 둘째 항은 \emph{영상태
응답}(zero-state response)이다. 행렬 $e^{A(t-t_0)}$은 \emph{상태천이행렬}(state
transition matrix)이다. 입력이 없을 때 이것이 시각 $t_0$의 상태를 시각 $t$의 상태로
옮긴다, $x(t)=e^{A(t-t_0)}x(t_0)$. 적분은 입력에 대해 같은 말을 한다. 시각 $\tau$에
가해진 작은 밀기 $Bu(\tau)\,\d\tau$가 $e^{A(t-\tau)}$에 실려 시각 $t$까지 옮겨지고,
그 밀기들이 모두 더해진다.

대신 $\dot x=Ax+Bu$를 라플라스 변환하면
$X(s)=(sI-A)^{-1}x_0+(sI-A)^{-1}BU(s)$이고, 이는 \eqref{eq:voc}와 항별로 같다.
시간 영역의 합성곱이 $s$ 영역의 곱이기 때문이다. 둘을 비교하면
\begin{equation}\label{eq:laplaceexpm}
   \Lap\{e^{At}\} = (sI-A)^{-1}.
\end{equation}
출력의 영상태 부분 $C(sI-A)^{-1}BU+DU$는 다시 한 번 \ref{sec:ss2tf}절의
전달함수이고, 영입력 부분은 전달함수가 빠뜨리는 자유 응답이다.

\subsection{모드: 대각화로 풀기}\label{sec:modes}
$A=V\Lambda V^{-1}$이 대각화 가능하면 \ref{sec:diag}절의 좌표 $z=V^{-1}x$로
바꾼다. 그러면 $\dot z=\Lambda z$는 따로따로인 $n$개의 스칼라 방정식
$\dot z_i=\lambda_iz_i$이므로 $z_i(t)=e^{\lambda_it}z_i(0)$이고, 되돌아가면
\begin{equation}\label{eq:modal}
   x(t) = Vz(t) = z_1(0)\,e^{\lambda_1t}\,v_1 + \dots + z_n(0)\,e^{\lambda_nt}\,v_n .
\end{equation}
각 항이 하나의 \emph{모드}(mode)다. 고정된 방향 $v_i$에 스칼라 지수함수가 곱해진
것이다. 고유벡터는 방향을, 고윳값은 그 성분이 얼마나 빨리 자라거나 줄어드는지를,
초기조건은 각 모드가 얼마만큼 들어 있는지를 정한다(그림~\ref{fig:modes}). 같은
공식을 좌표 없이 쓰면 \emph{좌고유벡터}(left eigenvector) $w_i\tp$가 쓰인다.
$V^{-1}$의 행들로, $w_i\tp A=\lambda_iw_i\tp$와 $w_i\tp v_j=\delta_{ij}$를
만족한다.
\begin{equation}\label{eq:expmdyad}
   e^{At} = Ve^{\Lambda t}V^{-1} = \sum_{i=1}^{n} e^{\lambda_it}\,v_iw_i\tp ,
   \qquad z_i(0) = w_i\tp x_0 .
\end{equation}
그러니 $e^{At}$는 모드마다 하나씩, 계수가 1인 행렬들의 합이다.

\begin{keybox}
\textbf{$\dot x=Ax$의 안정도.} 모든 초기 상태에서 $x(t)\to0$일 필요충분조건은 $A$의
모든 고윳값의 실수부가 엄격히 음수인 것이다.
\end{keybox}
대각화 가능한 $A$에서는 \eqref{eq:modal}에서 이것이 보인다. 복소쌍
$\sigma\pm j\omega$는 $e^{\sigma t}$에 주파수 $\omega$의 회전을 곱한 것을 내놓고,
모든 항이 사라지는 것은 정확히 모든 $\operatorname{Re}\lambda_i<0$일 때다.
\ref{sec:jordanresp}절에서 이 명제가 대각화 가능성 없이도 살아남음을 보인다. 이것은
정의~\ref{def:stable}의 안정한 플랜트 조건의 상태공간판 쌍둥이이고, (약분이 없다면)
$A$의 고윳값이 $P$의 극점이므로 둘은 같은 말을 한다.

\begin{figure}[htb]
\centering
\begin{tikzpicture}[>=Stealth,scale=1.35]
  \draw[->,gray] (-1.4,0) -- (2.6,0) node[right,lbl] {$x_1$};
  \draw[->,gray] (0,-0.9) -- (0,2.8) node[above,lbl] {$x_2$};
  \draw[dashed,gray] (-1.2,-0.48) -- (2.5,1.0);
  \draw[dashed,gray] (0.30,-0.86) -- (-0.92,2.63);
  \draw[->,very thick] (0,0) -- (1,0.4) node[below right,lbl] {$v_1$};
  \draw[->,very thick] (0,0) -- (-0.35,1) node[left,lbl] {$v_2$};
  \draw[->] (0,0) -- (2,0.8) node[right,lbl] {$z_1(0)v_1$};
  \draw[->] (0,0) -- (-0.525,1.5) node[left,lbl] {$z_2(0)v_2$};
  \draw[dotted] (2,0.8) -- (1.475,2.3) -- (-0.525,1.5);
  \fill (1.475,2.3) circle (1.4pt) node[above right,lbl] {$x(0)$};
  \draw[thick] plot[domain=0:9,samples=90,smooth,variable=\t]
    ({2*exp(-0.5*\t)-0.525*exp(-2*\t)},{0.8*exp(-0.5*\t)+1.5*exp(-2*\t)});
  \pgfmathsetmacro{\ax}{2*exp(-0.5*0.45)-0.525*exp(-2*0.45)}
  \pgfmathsetmacro{\ay}{0.8*exp(-0.5*0.45)+1.5*exp(-2*0.45)}
  \pgfmathsetmacro{\bx}{2*exp(-0.5*0.55)-0.525*exp(-2*0.55)}
  \pgfmathsetmacro{\by}{0.8*exp(-0.5*0.55)+1.5*exp(-2*0.55)}
  \draw[->,thick] (\ax,\ay) -- (\bx,\by);
  \fill (0,0) circle (1.2pt);
\end{tikzpicture}
\caption{고윳값 $\lambda_1=-0.5$, $\lambda_2=-2$로 그린 모드 분해. 초기 상태가
고유벡터를 따라 나뉘고(점선 평행사변형), 각 성분은 제 고유벡터를 따라 제 속도로
줄어든다. 빠른 모드 $v_2$가 먼저 사라지므로 궤적은 휘어서 느린 방향 $v_1$을 따라
원점으로 들어온다.}
\label{fig:modes}
\end{figure}

\begin{example}[행렬지수 하나, 세 가지 방법]\label{ex:expm3}
$A=\begin{bmatrix}0&1\\-2&-3\end{bmatrix}$이라 하자.
$\det(\lambda I-A)=\lambda^{2}+3\lambda+2=(\lambda+1)(\lambda+2)$이다.

\emph{대각화로.} $\lambda_1=-1$에 고유벡터 $v_1=(1,-1)\tp$, $\lambda_2=-2$에
$v_2=(1,-2)\tp$이므로
\[
   V=\begin{bmatrix}1&1\\-1&-2\end{bmatrix},\quad
   V^{-1}=\begin{bmatrix}2&1\\-1&-1\end{bmatrix},\quad
   e^{At}=e^{-t}\begin{bmatrix}2&1\\-2&-1\end{bmatrix}
         +e^{-2t}\begin{bmatrix}-1&-1\\2&2\end{bmatrix}
\]
이고, 이는 \eqref{eq:expmdyad}의 계수 1인 두 조각 $v_1w_1\tp$와 $v_2w_2\tp$이다.

\emph{라플라스로.} $(sI-A)^{-1}=\dfrac{1}{(s+1)(s+2)}\begin{bmatrix}s+3&1\\-2&s\end{bmatrix}$이고,
성분마다 \ref{sec:pfe}절의 부분분수, 예컨대
$\frac{s+3}{(s+1)(s+2)}=\frac{2}{s+1}-\frac{1}{s+2}$로 역변환하면 같은 행렬이 된다.

\emph{케일리--해밀턴으로}(\ref{sec:cayley}절). $e^{At}=\alpha_0(t)I+\alpha_1(t)A$로
쓰고 두 고윳값에서 $e^{\lambda t}=\alpha_0+\alpha_1\lambda$를 요구한다.
$e^{-t}=\alpha_0-\alpha_1$, $e^{-2t}=\alpha_0-2\alpha_1$. 따라서
$\alpha_1=e^{-t}-e^{-2t}$, $\alpha_0=2e^{-t}-e^{-2t}$이고,
$\alpha_0I+\alpha_1A$는 다시 위의 행렬이다.
\end{example}

\subsection{고유벡터가 모자랄 때}\label{sec:jordanresp}
조르당 표준형의 행렬에서는 분리가 불완전하지만 여전히 쉽다. \ref{sec:jordan}절의
$3\times3$ $J$로 $\dot z=Jz$를 쓰면
\[
   \dot z_1=\lambda_1z_1,\qquad \dot z_2=\lambda_2z_2+z_3,\qquad \dot z_3=\lambda_2z_3 .
\]
첫째와 셋째는 앞에서처럼 푼다. 가운데 것은 $z_3$이 몰고 가는데, $n=1$인
\eqref{eq:voc}에 따라
\[
   z_2(t) = e^{\lambda_2t}z_2(0)+\int_0^{t}e^{\lambda_2(t-\tau)}e^{\lambda_2\tau}z_3(0)\,\d\tau
          = e^{\lambda_2t}z_2(0) + t\,e^{\lambda_2t}z_3(0).
\]
따라서
\[
   x(t) = z_1(0)e^{\lambda_1t}v_1 + \bigl(z_2(0)e^{\lambda_2t}+z_3(0)\,t\,e^{\lambda_2t}\bigr)v_2
          + z_3(0)e^{\lambda_2t}v_3
\]
이고, 블록 하나 $J_2(\lambda)$에서는 $e^{J_2t}=e^{\lambda t}\begin{bmatrix}1&t\\0&1\end{bmatrix}$이다.
조르당 블록은 지수함수 앞에 인수 $t$를 데려온다. \ref{sec:pfe}절에서 중근 극점이
데려온 바로 그 $t$다. $\operatorname{Re}\lambda<0$이면 언제나
$t^{k}e^{\lambda t}\to0$이므로 위의 안정도 판정은 모든 정방 행렬에서 성립한다.

\subsection{케일리--해밀턴 정리}\label{sec:cayley}
\begin{theorem}[케일리--해밀턴]
모든 정방 행렬은 자기 자신의 특성방정식을 만족한다.
$p(\lambda)=\det(\lambda I-A)=\lambda^{n}+a_{n-1}\lambda^{n-1}+\dots+a_0$이면
\[
   p(A) = A^{n}+a_{n-1}A^{n-1}+\dots+a_1A+a_0I = 0 .
\]
\end{theorem}
\begin{proof}[개요]
$A=V\Lambda V^{-1}$이면 $A^{k}=V\Lambda^{k}V^{-1}$이므로
\[
   p(A)=V\,p(\Lambda)\,V^{-1},\qquad
   p(\Lambda)=\diag\bigl(p(\lambda_1),\dots,p(\lambda_n)\bigr)=0
\]
이다. 각 $\lambda_i$가 근이기 때문이다. 일반적인 $A$에서는 같은 논증을 조르당
표준형에서 블록별로 하면 된다. 거기서 $p$는 각 $k\times k$ 블록을 소멸시킬 만큼 큰
인수 $(\lambda-\lambda_i)^{k}$를 품고 있다.
\end{proof}
그 결과로 $A^{n}$은 $I,A,\dots,A^{n-1}$의 조합이고, 따라서 $A^{n+1}=A\cdot A^{n}$도
그렇고, 계속 그렇다. \emph{$A$의 모든 다항식은 차수가 $n-1$ 이하인 다항식으로
줄어들고}, 수렴하는 모든 멱급수도 마찬가지다. 특히 예제~\ref{ex:expm3}에서 쓴 대로
$e^{At}=\alpha_0(t)I+\alpha_1(t)A+\dots+\alpha_{n-1}(t)A^{n-1}$이다. 고윳값이 모두
다르면 계수는 $n$개의 스칼라 방정식 $e^{\lambda_it}=\sum_k\alpha_k(t)\lambda_i^{k}$에서
나온다.

\begin{example}[행렬 다항식 줄이기]
$A=\begin{bmatrix}2&1\\1&2\end{bmatrix}$이라 하면 $p(\lambda)=\lambda^{2}-4\lambda+3$이고
$A^{2}=4A-3I$이다. 그러면 $A^{3}=4A^{2}-3A=13A-12I$, $A^{4}=13A^{2}-12A=40A-39I$이므로
\[
   A^{4}+3A^{3}+A^{2}+I = 83A-77I .
\]
같은 이치로 어떤 두 수 $a,b$에 대해 $A^{10000}=aA+bI$이다. 스칼라 항등식
$\lambda^{10000}=a\lambda+b$를 고윳값 $1$과 $3$에서 계산하면
$a=\frac{3^{10000}-1}{2}$, $b=\frac{3-3^{10000}}{2}$이다.
\end{example}

\begin{remark}[$e^{At}$를 계산하는 네 가지 방법]
급수 정의; $(sI-A)^{-1}$의 역라플라스 변환 \eqref{eq:laplaceexpm}; $A$가 대각화
가능할 때의 모드 합 $\sum_ie^{\lambda_it}v_iw_i\tp$(가능하지 않으면 조르당 표준형으로);
그리고 차수 $n-1$의 케일리--해밀턴 다항식. 이들은 한 대상을 보는 네 관점이고,
예제~\ref{ex:expm3}는 그중 셋을 서로 맞대어 확인한다.
\end{remark}

%%%%%%%%%%%%%%%%%%%%%%%%%%%%%%%%%%%%%%%%%%%%%%%%%%%%%%%%%%%%%%%%%%%%%%%%%%%%%%
\section{양의 정부호와 특이값 분해}\label{sec:svd}

\subsection{이차형식과 양의 정부호 행렬}\label{sec:pd}
\emph{이차형식}(quadratic form)은 정방 행렬 $P$에 대한 $x\tp Px$다. $x\tp Px$는
스칼라이므로 자기 전치 $x\tp P\tp x$와 같고, 그래서 $P$의 대칭 부분만 기여한다.
\[
   x\tp Px = x\tp\Bigl(\tfrac{P+P\tp}{2}\Bigr)x .
\]
그러므로 $P=P\tp$로 두어도 잃는 것이 없다.
\begin{definition}[양의 정부호]
대칭행렬 $P$가 모든 $x\neq0$에 대해 $x\tp Px>0$이면 \emph{양의 정부호}(positive
definite)라 하고 $P\succ0$으로 쓰며, 모든 $x$에 대해 $x\tp Px\ge0$이면 \emph{양의
준정부호}(positive semidefinite)라 하고 $P\succeq0$으로 쓴다.
\end{definition}
양의 정부호 행렬은 크기를 재는 행렬이다. 에너지, 비용, 거리의 제곱은 모두 이런
이차형식이다.

\begin{proposition}[두 가지 판정법]\label{prop:pdtests}
$P=P\tp\in\R^{n\times n}$에 대해 다음은 동치다.
\begin{enumerate}[label=(\alph*),itemsep=2pt]
\item $P\succ0$;
\item $P$의 모든 고윳값이 양수다;
\item 열 완전 계수인 어떤 $L$에 대해 $P=L\tp L$이다.
\end{enumerate}
\end{proposition}
\begin{proof}
정리~\ref{thm:spectral}에 따라 $P=V\Lambda V\tp$이므로 $y=V\tp x$로 두면
$x\tp Px=y\tp\Lambda y=\sum_i\lambda_iy_i^{2}$이다. $x$가 영 아닌 벡터 전체를 움직이면
$y$도 그렇고, 그 모두에서 합이 양수인 것은 정확히 모든 $\lambda_i>0$일 때다:
(a)$\Leftrightarrow$(b). $P=L\tp L$이면 $x\tp Px=\|Lx\|^{2}\ge0$이고, 등호는
$Lx=0$일 때만, 곧 $\mathcal N(L)=\{0\}$이면 $x=0$일 때만 성립한다:
(c)$\Rightarrow$(a). 거꾸로, (b)이면 $L=\Lambda^{1/2}V\tp$로 잡을 수 있다.
\end{proof}
준정부호는 ``양수''를 ``음이 아닌 수''로 바꾸고 $L$의 계수 조건을 빼면 같다.
기하학적으로, $P\succ0$이면 등위집합 $x\tp Px=1$은 축이 $P$의 고유벡터 방향이고
$v_i$ 방향의 반축이 $1/\sqrt{\lambda_i}$인 타원체다.
$P=\begin{bmatrix}2&1\\1&2\end{bmatrix}$의 고윳값은 $1$과 $3$이므로 $P\succ0$이고,
$2x_1^{2}+2x_1x_2+2x_2^{2}=1$은 긴 축이 $(1,-1)$ 방향인 타원이다.

\subsection{네 기본 부분공간}\label{sec:foursub}
계수가 $r$인 모든 $A\in\R^{m\times n}$은 \emph{네 기본 부분공간}(four fundamental
subspaces)을 가진다. $A$가 잇는 두 공간에 둘씩이다.
\begin{center}\small
\renewcommand{\arraystretch}{1.3}
\begin{tabular}{@{}llll@{}}
\toprule
부분공간 & 속한 공간 & 차원 & 이름\\
\midrule
$\mathcal R(A\tp)$ & $\R^{n}$ & $r$   & 행공간\\
$\mathcal N(A)$    & $\R^{n}$ & $n-r$ & 영공간\\
$\mathcal R(A)$    & $\R^{m}$ & $r$   & 열공간\\
$\mathcal N(A\tp)$ & $\R^{m}$ & $m-r$ & 왼쪽 영공간\\
\bottomrule
\end{tabular}
\end{center}
여기서 $\mathcal R(A\tp)$를 \emph{행공간}(row space), $\mathcal N(A\tp)$를
\emph{왼쪽 영공간}(left null space)이라 한다. 짝을 이루는 두 공간은 차원이 서로
보완적일 뿐 아니라 서로 수직이다. $Ax=0$이고 $y=A\tp z$이면 $y\tp x=z\tp Ax=0$이므로
\[
   \mathcal N(A)\perp\mathcal R(A\tp),\qquad\text{마찬가지로}\qquad
   \mathcal N(A\tp)\perp\mathcal R(A).
\]
그러므로 모든 $x\in\R^{n}$은 행공간의 $x_r$과 영공간의 $x_n$으로 $x=x_r+x_n$처럼
유일하게 갈라지고, $Ax=Ax_r$이다. 영공간 부분은 사라지고, $A$는 행공간을 열공간 위로
일대일로 옮긴다(그림~\ref{fig:foursub}).

\begin{figure}[htb]
\centering
\begin{tikzpicture}[>=Stealth,scale=1.15]
  % ---- R^n ----
  \node[lbl] at (-1.7,1.9) {$\R^{n}$};
  \draw[thick] ($(0,0)+(215:1.3)$) -- ($(0,0)+(35:2.0)$);
  \draw[thick] ($(0,0)+(125:1.5)$) -- ($(0,0)+(305:1.4)$);
  \draw ($(35:0.18)$) -- ($(35:0.18)+(125:0.18)$) -- ($(125:0.18)$);
  \node[lbl,anchor=south west] at (35:2.0) {$\mathcal R(A\tp)$, 차원 $r$};
  \node[lbl,anchor=south east] at (125:1.5) {$\mathcal N(A)$, 차원 $n-r$};
  \fill (0,0) circle (1.2pt);
  \fill (35:1.35) circle (1.4pt) node[above left,lbl] {$x_r$};
  \fill (305:1.0) circle (1.4pt) node[left,lbl] {$x_n$};
  % ---- R^m ----
  \begin{scope}[shift={(6.2,0)}]
    \node[lbl] at (2.1,1.9) {$\R^{m}$};
    \draw[thick] ($(215:1.4)$) -- ($(35:2.0)$);
    \draw[thick] ($(125:1.3)$) -- ($(305:1.4)$);
    \draw ($(35:0.18)$) -- ($(35:0.18)+(125:0.18)$) -- ($(125:0.18)$);
    \node[lbl,anchor=south west] at (35:2.0) {$\mathcal R(A)$, 차원 $r$};
    \node[lbl,anchor=north west] at (305:1.4) {$\mathcal N(A\tp)$, 차원 $m-r$};
    \fill (0,0) circle (1.2pt) node[below left,lbl] {$0$};
    \fill (35:1.55) circle (1.4pt) node[above left,lbl] {$Ax_r$};
  \end{scope}
  \draw[->] ($(35:1.35)+(0.08,0.02)$) to[bend left=12]
        node[above,lbl] {$A$} ($(6.2,0)+(35:1.55)+(-0.08,0)$);
  \draw[->] ($(305:1.0)+(0.08,-0.02)$) to[bend right=14]
        node[below,lbl] {$A$} ($(6.2,0)+(-0.06,-0.04)$);
\end{tikzpicture}
\caption{$A\in\R^{m\times n}$의 네 기본 부분공간. 각각을 직선 하나로 그렸다.
$\R^{n}$에서 행공간과 영공간이 수직이고, $\R^{m}$에서 열공간과 왼쪽 영공간이
수직이다. $A$는 벡터의 행공간 부분 $x_r$을 열공간 위로 일대일로 옮기고, 영공간
부분 $x_n$은 $0$으로 보낸다.}
\label{fig:foursub}
\end{figure}

$A$와 대칭행렬 $A\tp A$를 잇는 사실이 하나 더 있다.
\begin{equation}\label{eq:nullata}
   \mathcal N(A\tp A) = \mathcal N(A),\qquad\text{따라서}\qquad \rank(A\tp A)=\rank A .
\end{equation}
$Ax=0$이면 $A\tp Ax=0$인 것은 분명하다. 거꾸로 $A\tp Ax=0$이면
$0=x\tp A\tp Ax=\|Ax\|^{2}$이므로 $Ax=0$이다. 특히 $A\tp A$가 가역인 것은 정확히 $A$가
열 완전 계수일 때이고, 그때 $L=A$로 둔 명제~\ref{prop:pdtests}(c)에 따라
$A\tp A\succ0$이다. 어느 경우든 $A\tp A\succeq0$이다.

\subsection{특이값 분해}\label{sec:svdthm}
고유벡터는 정방 행렬에서만 뜻이 있고, 그때조차 직교한다는 보장도, 개수가 넉넉하다는
보장도 없다. \emph{특이값 분해}(singular value decomposition, SVD)는 이 모든 것을
모든 행렬에 대해 한꺼번에 바로잡는다.
\begin{theorem}[SVD]\label{thm:svd}
계수가 $r$인 모든 $A\in\R^{m\times n}$은
\[
   A = U\Sigma V\tp
\]
로 쓸 수 있다. 여기서 $U\in\R^{m\times m}$과 $V\in\R^{n\times n}$은 직교행렬($U\tp U=I$,
$V\tp V=I$)이고, $\Sigma\in\R^{m\times n}$은 대각 성분
$\sigma_1\ge\sigma_2\ge\dots\ge\sigma_r>0$을 빼고는 $0$이다. $\sigma_i$를
\emph{특이값}(singular value), $V$의 열 $v_i$를 오른쪽 \emph{특이벡터}(singular
vector), $U$의 열 $u_i$를 왼쪽 특이벡터라 한다. 같은 말로
\begin{equation}\label{eq:svddyad}
   A = \sum_{i=1}^{r}\sigma_i\,u_iv_i\tp ,\qquad Av_i=\sigma_iu_i,\qquad A\tp u_i=\sigma_iv_i .
\end{equation}
\end{theorem}
\begin{proof}[구성]
$A\tp A$는 대칭이고 양의 준정부호이므로 정리~\ref{thm:spectral}에 따라 고윳값
$\lambda_1\ge\dots\ge\lambda_n\ge0$에 속하는 정규직교 고유벡터 $v_1,\dots,v_n$을
가지며, \eqref{eq:nullata}에 따라 그중 정확히 $r$개가 양수다. $i\le r$에 대해
$\sigma_i=\sqrt{\lambda_i}$, $u_i=Av_i/\sigma_i$로 두자. 이들은 정규직교다.
\[
   u_i\tp u_j=\frac{v_i\tp A\tp Av_j}{\sigma_i\sigma_j}
            =\frac{\lambda_j\,v_i\tp v_j}{\sigma_i\sigma_j}=\delta_{ij}.
\]
$u_1,\dots,u_r$을 $\R^{m}$의 정규직교 기저로 채운다. 그러면 열마다 $AV=U\Sigma$이고
($i>r$이면 그 $v_i$가 $\mathcal N(A\tp A)=\mathcal N(A)$에 있으므로 $Av_i=0$),
오른쪽에 $V\tp$를 곱하면 $A=U\Sigma V\tp$이다.
\end{proof}
그러므로 특이값은
\[
   \sigma_i(A) = \sqrt{\lambda_i(A\tp A)}
\]
이고, 특이벡터는 \ref{sec:foursub}절의 네 부분공간 모두의 정규직교 기저를 한꺼번에
준다. $v_1,\dots,v_r$은 행공간을, $v_{r+1},\dots,v_n$은 영공간을 생성하고,
$u_1,\dots,u_r$은 열공간을, $u_{r+1},\dots,u_m$은 왼쪽 영공간을 생성한다. 계수는 영
아닌 특이값의 개수다. 합 \eqref{eq:svddyad}는 \ref{sec:vspace}절의 읽기 (iii)
$AB=\sum a_ib_i\tp$에서 인수들을 정규직교로 고르고 크기순으로 늘어놓은 것이다.

\begin{example}[$2\times2$ SVD]\label{ex:svd}
$A=\begin{bmatrix}3&0\\4&5\end{bmatrix}$이라 하자. 그러면
$A\tp A=\begin{bmatrix}25&20\\20&25\end{bmatrix}$이고, 고윳값은 $45$와 $5$,
고유벡터는 $v_1=\tfrac{1}{\sqrt2}(1,1)\tp$, $v_2=\tfrac1{\sqrt2}(1,-1)\tp$이다.
그러므로 $\sigma_1=3\sqrt5\approx6.71$, $\sigma_2=\sqrt5\approx2.24$이고
\[
   u_1=\frac{Av_1}{\sigma_1}=\frac{1}{\sqrt{10}}\begin{bmatrix}1\\3\end{bmatrix},\qquad
   u_2=\frac{Av_2}{\sigma_2}=\frac{1}{\sqrt{10}}\begin{bmatrix}3\\-1\end{bmatrix}
\]
이며, 실제로 직교한다. 확인해 보면
$\sigma_1u_1v_1\tp+\sigma_2u_2v_2\tp=\tfrac32\begin{bmatrix}1&1\\3&3\end{bmatrix}
+\tfrac12\begin{bmatrix}3&-3\\-1&1\end{bmatrix}=A$이고
$\sigma_1\sigma_2=15=|\det A|$이다. $A$의 고윳값은 $3$과 $5$로, 특이값이 아니다.
대칭이 아닌 행렬에서 둘은 서로 다른 대상이고, 대칭행렬에서는 특이값이 고윳값의
절댓값이다.
\end{example}

\subsection{특이값이 재는 것}\label{sec:svduse}
$A=U\Sigma V\tp$를 오른쪽에서 왼쪽으로 읽자. $V\tp$는 입력을 회전시키고(직교행렬이므로
길이를 보존한다), $\Sigma$는 $i$번째 좌표를 $\sigma_i$배로 늘이며, $U$는 그 결과를
회전시킨다. 그러므로 $A$는 단위 구를 반축이 $\sigma_iu_i$인 타원체 위로 옮기고, 그
축에 떨어지는 입력 방향이 $v_i$다(그림~\ref{fig:svd}).

\begin{figure}[htb]
\centering
\begin{tikzpicture}[>=Stealth]
  % ---- input: unit circle ----
  \begin{scope}[scale=1.15]
    \draw[->,gray] (-1.5,0) -- (1.6,0); \draw[->,gray] (0,-1.5) -- (0,1.6);
    \draw[thick] (0,0) circle (1);
    \draw[->,very thick] (0,0) -- (45:1) node[above right,lbl] {$v_1$};
    \draw[->,very thick] (0,0) -- (-45:1) node[below right,lbl] {$v_2$};
    \node[lbl] at (-1.15,1.35) {$\|x\|=1$};
  \end{scope}
  \draw[->,thick] (2.2,0.35) -- node[above,sig] {$A$} (3.4,0.35);
  % ---- output: ellipse A(circle) ----
  \begin{scope}[xshift=5.6cm,scale=0.33]
    \draw[->,gray] (-4.6,0) -- (4.8,0); \draw[->,gray] (0,-7.2) -- (0,7.6);
    \draw[thick] plot[domain=0:360,samples=120,smooth,variable=\a]
      ({3*cos(\a)},{4*cos(\a)+5*sin(\a)});
    \draw[->,very thick] (0,0) -- (2.121,6.364) node[right,lbl] {$\sigma_1u_1$};
    \draw[->,very thick] (0,0) -- (2.121,-0.707) node[right,lbl] {$\sigma_2u_2$};
  \end{scope}
\end{tikzpicture}
\caption{예제~\ref{ex:svd}의 SVD를 그림으로. 단위원(왼쪽)이 $A$에 실려 반축
$\sigma_1u_1$, $\sigma_2u_2$(길이 $3\sqrt5$, $\sqrt5$)를 가진 타원(오른쪽, 축척
3분의 1로 그림) 위로 옮겨진다. 입력 방향 $v_1$과 $v_2$가 그 축에 떨어지는
방향이다.}
\label{fig:svd}
\end{figure}

여기서 쓸모 있는 수 몇 가지가 바로 나온다.
\begin{itemize}[leftmargin=*,itemsep=2pt]
\item \textbf{이득.} 모든 $x\neq0$에 대해
$\sigma_{\min}\le\|Ax\|/\|x\|\le\sigma_1$이다($A$가 열 완전 계수이면
$\sigma_{\min}=\sigma_n$, 아니면 $0$). 가장 큰 증폭은 $\sigma_1$이고 $x$가 $v_1$
방향일 때 얻어진다. 이것이 \emph{유도 $2$-노름}(induced $2$-norm)
$\|A\|_2=\sigma_1$이다. 그러니 특이값은 행렬이 얼마나 증폭할 수 있는지, 어느 입력
방향에서 그런지를 알려 준다. 고윳값이 각 모드가 얼마나 빨리 변하는지를 알려 주는
것과 대비된다.
\item \textbf{조건.} 비 $\kappa(A)=\sigma_1/\sigma_r$이 \emph{조건수}(condition
number)다. 이것이 크면 어떤 입력 방향은 다른 방향보다 훨씬 크게 증폭되고, $Ax=b$를
풀 때 $b$의 오차가 최대 $\kappa$배까지 커진다. 서류상으로는 가역이지만 $\sigma_r$이
아주 작은 행렬은 수치적으로는 거의 계수가 모자란 행렬이다.
\item \textbf{저계수 근사.} \eqref{eq:svddyad}를 $k$개 항에서 자른
$A_k=\sum_{i\le k}\sigma_iu_iv_i\tp$는 $A$에 가장 가까운 계수 $k$의 행렬이고, 오차는
$\|A-A_k\|_2=\sigma_{k+1}$이다. 이것이 \emph{저계수 근사}(low-rank
approximation)다. 가장 강한 방향만 남길 때 얼마나 잃는지를 특이값이 정확히 알려 준다.
\end{itemize}

\begin{remark}[최소제곱]
$Ax=b$의 해가 없을 때($b\notin\mathcal R(A)$)에도 SVD는 할 수 있는 최선을 준다.
\emph{유사역행렬}(pseudo-inverse) $A^{+}=V\Sigma^{+}U\tp$, 여기서 $\Sigma^{+}$는 영
아닌 특이값을 역수로 바꾸고 모양을 전치한 것인데, 이것이 주는 $x=A^{+}b$는
$\|Ax-b\|$를 최소로 만들고, 그런 최소화 해 가운데 노름이 가장 작다. 이것이
\emph{최소제곱}(least squares) 해다. $A$가 열 완전 계수이면
$A^{+}=(A\tp A)^{-1}A\tp$이고, \eqref{eq:nullata}에 따라 $A\tp A$는 가역이다.
\end{remark}

\end{document}
